% Équations, inéquations et systèmes — Mathématiques Tle A — Cours signé OnBuch+
% Programme officiel MINESEC. Compiler avec : tectonic MA01-equations-inequations-systemes-a.tex
\newcommand{\DOCMATIERE}{Mathématiques}
\newcommand{\DOCNIVEAU}{Tle A}
\newcommand{\DOCMODULE}{Relations et opérations fondamentales dans l'ensemble des nombres réels}
\newcommand{\DOCLECON}{A01}
\newcommand{\DOCTITRE}{Équations, inéquations et systèmes}
% OnBuch+ — préambule « cours signés » ENRICHI (compilé avec tectonic / XeLaTeX)
% Système visuel « Pop » d'OnBuch+ : fond crème (jamais blanc), bordures encre
% épaisses, ombres « dures » décalées, titres Archivo Black en capitales.
% Version MATHÉMATIQUES : graphes (pgfplots/tikz), boîtes pédagogiques
% (définition, propriété, méthode, attention, le savais-tu, exercice résolu,
% exercice, TP, bilan), page de garde et signature OnBuch+ ancrée en bas.
%
% Macros attendues dans main.tex AVANT \input{preamble} :
%   \DOCMATIERE  \DOCNIVEAU  \DOCMODULE  \DOCLECON (numéro)  \DOCTITRE
\documentclass[11pt]{article}

\usepackage{fontspec}
\usepackage[french]{babel}
\usepackage[a4paper,margin=2cm,top=3.4cm,bottom=2.8cm,headheight=1.4cm,footskip=1.3cm]{geometry}
\usepackage{amsmath,amssymb,amsfonts}
\usepackage{mathtools} % \xmapsto, \coloneqq... souvent utilisés par les modèles
\usepackage{cancel} % simplifications barrées
% \abs{z} (module, valeur absolue) et \norm{u} (norme), utilisés par les modèles
\providecommand{\abs}{}\let\abs\relax\DeclarePairedDelimiter{\abs}{\lvert}{\rvert}
\providecommand{\norm}{}\let\norm\relax\DeclarePairedDelimiter{\norm}{\lVert}{\rVert}
\usepackage[table]{xcolor} % [table] : fournit aussi \rowcolors (alternance de lignes)
\usepackage{pagecolor}
\usepackage{tikz}
\usetikzlibrary{arrows.meta,positioning,calc,shapes.geometric,shapes.misc,shapes.arrows,decorations.pathmorphing,decorations.markings,patterns,shadows,mindmap,trees,fit,backgrounds,matrix,intersections}
\usepackage{pgfplots}
\usepackage{pgf-pie} % diagrammes circulaires \pie (statistiques)
\pgfplotsset{compat=1.18}
\usepgfplotslibrary{fillbetween,groupplots} % aires entre courbes (calcul intégral)
\usepackage{enumitem}
\usepackage{fancyhdr}
% [most] charge toutes les bibliothèques tcolorbox (listings, minted, raster,
% documentation...) et gonfle inutilement la mémoire TeX sur un document long
% (~300 boîtes) : on ne charge que ce qui est réellement utilisé.
\usepackage[skins,breakable]{tcolorbox}
\usepackage{graphicx}
\usepackage{colortbl}
\usepackage{booktabs}
\usepackage{tabularx}
\usepackage{diagbox} % case d'en-tête diagonale des tableaux à double entrée
% Colonnes extensibles centrée / gauche / droite (C, L, R), souvent utilisées par les modèles
\newcolumntype{C}{>{\centering\arraybackslash}X}
\newcolumntype{L}{>{\raggedright\arraybackslash}X}
\newcolumntype{R}{>{\raggedleft\arraybackslash}X}
\usepackage{array}
\usepackage{multicol}
\usepackage{siunitx}
% parse-numbers=false : accepte aussi \SI{2,5 \times 10^{-3}}{...}
\sisetup{locale=FR,output-decimal-marker={,},group-minimum-digits=4,parse-numbers=false}
\DeclareSIUnit\ohm{\ensuremath{\symup{\Omega}}} % U+2126 absent de la police
\usepackage{qrcode}
% \degree : commande fréquemment utilisée par les modèles pour un angle ou une
% température en dehors de \SI{}{} (non fournie par les paquets chargés).
\providecommand{\degree}{^{\circ}}
% \qed (fin de démonstration), utilisé par les modèles sans amsthm
\providecommand{\qed}{\hfill$\square$}
% \barre : double barre des valeurs interdites dans un tableau de variations (tkz-tab)
\providecommand{\barre}{\ensuremath{\|}}
% environnement proof (amsthm non chargé) utilisé par les modèles
\ifdefined\proof\else\newenvironment{proof}[1][Démonstration]{\par\noindent\textit{#1.}\ }{\qed\par}\fi
\usepackage{lastpage}
\usepackage{hyperref}
\hypersetup{hidelinks}

% ============================================================
% Palette « Pop » OnBuch+ — mêmes hex que la classe P (plus_ui.dart)
% ============================================================
\definecolor{poporange}{HTML}{F59321}
\definecolor{popdark}{HTML}{E07A0C}
\definecolor{popcream}{HTML}{FFF6E8}
\definecolor{popink}{HTML}{1C1714}
\definecolor{poppurple}{HTML}{7A5AE0}
\definecolor{popgreen}{HTML}{2E9E6A}
\definecolor{poppink}{HTML}{E0457B}
\definecolor{popblue}{HTML}{2F7DE1}
\definecolor{popgold}{HTML}{C98A12}
\definecolor{popmuted}{HTML}{8A7F76}
\definecolor{popline}{HTML}{EADFD2}
\definecolor{popgray}{HTML}{8A7F76}
\colorlet{popgrey}{popgray}
% Alias tolérants (le modèle nomme parfois les couleurs différemment)
\colorlet{popwhite}{white}
\colorlet{popblack}{popink}
\colorlet{popred}{poppink}
\colorlet{popyellow}{popgold}
% Teintes claires pour fonds de tableaux / zones
\colorlet{popblueL}{popblue!10!white}
\colorlet{poporangeL}{poporange!14!white}
\colorlet{popgreenL}{popgreen!12!white}
\colorlet{poppurpleL}{poppurple!12!white}
\colorlet{poppinkL}{poppink!12!white}
\colorlet{popgoldL}{popgold!14!white}

\pagecolor{popcream}

% ============================================================
% Typographie
% ============================================================
\defaultfontfeatures{Ligatures=TeX}
\setmainfont{PlusJakartaSans-Medium.ttf}[Path=../../../fonts/,BoldFont=PlusJakartaSans-Bold.ttf]
% Maths en sans-serif (Fira Math) avec chiffres et lettres droites de Plus
% Jakarta, pour que les formules se fondent dans le texte.
\usepackage[math-style=ISO,bold-style=ISO]{unicode-math}
\setmathfont{FiraMath-Regular.otf}[Path=../../../fonts/]
\setmathfont{PlusJakartaSans-Medium.ttf}[Path=../../../fonts/,range={up/num,up/{Latin,latin},\mathup}]
\usepackage{icomma} % virgule décimale sans espace en mode math
% Caractères mathématiques Unicode absents des polices de texte (Plus Jakarta,
% Archivo Black) : rendus par leur équivalent mathématique, en texte comme en
% formule — sinon ils s'affichent en carré vide (ex. « Arithmétique dans ℤ »).
\usepackage{newunicodechar}
\newunicodechar{ℝ}{\ensuremath{\mathbb{R}}}
\newunicodechar{ℤ}{\ensuremath{\mathbb{Z}}}
\newunicodechar{ℂ}{\ensuremath{\mathbb{C}}}
\newunicodechar{ℕ}{\ensuremath{\mathbb{N}}}
\newunicodechar{ℚ}{\ensuremath{\mathbb{Q}}}
\newunicodechar{𝔻}{\ensuremath{\mathbb{D}}}
\newunicodechar{α}{\ensuremath{\alpha}}
\newunicodechar{β}{\ensuremath{\beta}}
\newunicodechar{γ}{\ensuremath{\gamma}}
\newunicodechar{δ}{\ensuremath{\delta}}
\newunicodechar{ε}{\ensuremath{\varepsilon}}
\newunicodechar{θ}{\ensuremath{\theta}}
\newunicodechar{λ}{\ensuremath{\lambda}}
\newunicodechar{μ}{\ensuremath{\mu}}
\newunicodechar{σ}{\ensuremath{\sigma}}
\newunicodechar{τ}{\ensuremath{\tau}}
\newunicodechar{φ}{\ensuremath{\varphi}}
\newunicodechar{ψ}{\ensuremath{\psi}}
\newunicodechar{ω}{\ensuremath{\omega}}
\newunicodechar{Σ}{\ensuremath{\Sigma}}
% majuscules produites par \MakeUppercase dans les titres (α-aminés -> Α)
\newunicodechar{Α}{\ensuremath{\alpha}}
\newunicodechar{Β}{\ensuremath{\beta}}
\newunicodechar{Γ}{\ensuremath{\Gamma}}
\newunicodechar{Δ}{\ensuremath{\Delta}}
\newunicodechar{Ε}{\ensuremath{\varepsilon}}
\newunicodechar{Θ}{\ensuremath{\Theta}}
\newunicodechar{Λ}{\ensuremath{\Lambda}}
\newunicodechar{Μ}{\ensuremath{\mu}}
\newunicodechar{Τ}{\ensuremath{\tau}}
\newunicodechar{Φ}{\ensuremath{\Phi}}
\newunicodechar{Ψ}{\ensuremath{\Psi}}
\newunicodechar{Ω}{\ensuremath{\Omega}}
\newunicodechar{↦}{\ensuremath{\mapsto}}
\newunicodechar{∈}{\ensuremath{\in}}
\newunicodechar{∉}{\ensuremath{\notin}}
\newunicodechar{∧}{\ensuremath{\wedge}}
\newunicodechar{⇔}{\ensuremath{\Leftrightarrow}}
\newunicodechar{⟺}{\ensuremath{\Longleftrightarrow}}
\newunicodechar{⇒}{\ensuremath{\Rightarrow}}
\newunicodechar{⟹}{\ensuremath{\Longrightarrow}}
\newunicodechar{∘}{\ensuremath{\circ}}
\newunicodechar{∪}{\ensuremath{\cup}}
\newunicodechar{∩}{\ensuremath{\cap}}
\newunicodechar{≡}{\ensuremath{\equiv}}
\newunicodechar{⊂}{\ensuremath{\subset}}
\newunicodechar{⊥}{\ensuremath{\perp}}
\newunicodechar{∃}{\ensuremath{\exists}}
\newunicodechar{∀}{\ensuremath{\forall}}
\newunicodechar{∛}{\ensuremath{\sqrt[3]{\,}}}
\newunicodechar{∜}{\ensuremath{\sqrt[4]{\,}}}
\newunicodechar{⃗}{\ensuremath{{}^{\to}}}
\newunicodechar{ⁿ}{\ifmmode^{n}\else\textsuperscript{\ensuremath{n}}\fi}
\newunicodechar{ˣ}{\ifmmode^{x}\else\textsuperscript{\ensuremath{x}}\fi}
\newunicodechar{ᵇ}{\ifmmode^{b}\else\textsuperscript{\ensuremath{b}}\fi}
\newunicodechar{ʳ}{\ifmmode^{r}\else\textsuperscript{\ensuremath{r}}\fi}
\newunicodechar{ᵘ}{\ifmmode^{u}\else\textsuperscript{\ensuremath{u}}\fi}
\newunicodechar{ʸ}{\ifmmode^{y}\else\textsuperscript{\ensuremath{y}}\fi}
\newunicodechar{ᵃ}{\ifmmode^{a}\else\textsuperscript{\ensuremath{a}}\fi}
\newunicodechar{ᵉ}{\ifmmode^{e}\else\textsuperscript{\ensuremath{e}}\fi}
\newunicodechar{ᵏ}{\ifmmode^{k}\else\textsuperscript{\ensuremath{k}}\fi}
\newunicodechar{ᵖ}{\ifmmode^{p}\else\textsuperscript{\ensuremath{p}}\fi}
\newunicodechar{ᴺ}{\ifmmode^{N}\else\textsuperscript{\ensuremath{N}}\fi}
\newunicodechar{ᶜ}{\ifmmode^{c}\else\textsuperscript{\ensuremath{c}}\fi}
\newunicodechar{ʰ}{\ifmmode^{h}\else\textsuperscript{\ensuremath{h}}\fi}
\newunicodechar{ᵠ}{\ifmmode^{q}\else\textsuperscript{\ensuremath{q}}\fi}
\newunicodechar{ₙ}{\ifmmode_{n}\else\textsubscript{\ensuremath{n}}\fi}
\newunicodechar{ₐ}{\ifmmode_{a}\else\textsubscript{\ensuremath{a}}\fi}
\newunicodechar{ᵢ}{\ifmmode_{i}\else\textsubscript{\ensuremath{i}}\fi}
\newunicodechar{ᵣ}{\ifmmode_{r}\else\textsubscript{\ensuremath{r}}\fi}
\newunicodechar{ₖ}{\ifmmode_{k}\else\textsubscript{\ensuremath{k}}\fi}
\newunicodechar{ᵦ}{\ifmmode_{\beta}\else\textsubscript{\ensuremath{\beta}}\fi}
\newunicodechar{ₓ}{\ifmmode_{x}\else\textsubscript{\ensuremath{x}}\fi}
\newfontfamily\popdisplay{ArchivoBlack-Regular.ttf}[Path=../../../fonts/]
\newfontfamily\poplogo{SpaceGrotesk-Bold.ttf}[Path=../../../fonts/]
\newfontfamily\popsemi{PlusJakartaSans-SemiBold.ttf}[Path=../../../fonts/]
\newfontfamily\popxbold{PlusJakartaSans-ExtraBold.ttf}[Path=../../../fonts/]
\newfontfamily\popxboldls{PlusJakartaSans-ExtraBold.ttf}[Path=../../../fonts/,LetterSpace=3.0]

\setlist[enumerate]{leftmargin=1.7em,itemsep=5pt,topsep=4pt}
\setlist[itemize]{leftmargin=1.7em,itemsep=4pt,topsep=4pt,label={\color{poporange}\textbullet}}
\setlength{\parindent}{0pt}
\setlength{\parskip}{5pt}
\renewcommand{\arraystretch}{1.25}

% pgfplots : style Pop par défaut
\pgfplotsset{
  every axis/.append style={axis line style={popink,line width=1pt},tick style={popink},
    grid style={popline,line width=0.5pt},label style={font=\small},tick label style={font=\footnotesize}},
  popcurve/.style={poporange,line width=1.8pt,no markers,smooth},
}
% Même style disponible dans un \draw TikZ simple (hors pgfplots)
\tikzset{popcurve/.style={poporange,line width=1.8pt}}

% ============================================================
% Pill / squiggle
% ============================================================
\newcommand{\poppill}[3]{%
  \tikz[baseline=-0.55ex]\node[draw=popink,line width=1.2pt,rounded corners=2.6pt,%
    fill=#2,inner xsep=6.5pt,inner ysep=3pt]{%
    \popxboldls\fontsize{7}{7}\selectfont\color{#3}\MakeUppercase{#1}};%
}
\newcommand{\popsquiggle}[1]{%
  \tikz[baseline=-0.4ex]{
    \draw[#1,line width=2.6pt,line cap=round]
      plot [smooth] coordinates {(0,0.09) (0.34,0.01) (0.68,0.21) (1.02,0.01) (1.36,0.10)};
  }%
}
% Mot-clé surligné (marqueur orange sous le texte)
\newcommand{\cle}[1]{{\popxbold\color{popdark}#1}}

% ============================================================
% En-tête / pied
% ============================================================
\pagestyle{fancy}
\fancyhf{}
\renewcommand{\headrulewidth}{0pt}
\renewcommand{\footrulewidth}{0pt}
\lhead{\raisebox{-0.15ex}{\poplogo\fontsize{13}{13}\selectfont\color{popink}OnBuch\color{poporange}+}}
\rhead{\poppill{\DOCMATIERE\ \textbullet\ \DOCNIVEAU\ \textbullet\ Leçon \DOCLECON}{white}{popink}}
\lfoot{\popxbold\fontsize{7.6}{7.6}\selectfont\color{popmuted}\MakeUppercase{Cours signé OnBuch+}\ \textbullet\ {\popsemi\DOCTITRE}}
\rfoot{\tikz[baseline=-0.6ex]\node[rounded corners=8.5pt,fill=popink,minimum height=17pt,minimum width=17pt,inner xsep=5pt,inner sep=0]{\popxbold\fontsize{8}{8}\selectfont\color{popcream}\thepage\,/\,\pageref*{LastPage}};}

% ============================================================
% Titres
% ============================================================
\newcommand{\chaptitre}[2]{%
  \begin{tcolorbox}[enhanced,colback=poporange,colframe=popink,boxrule=2.6pt,arc=11pt,
      drop shadow=popink,left=15pt,right=15pt,top=12pt,bottom=11pt,before skip=3pt,after skip=18pt]
    \poppill{#2}{popink}{popcream}\par\vspace{9pt}
    {\raggedright\hyphenpenalty=10000\exhyphenpenalty=10000\popdisplay\fontsize{22}{24}\selectfont\color{popcream}\MakeUppercase{#1}\par}
  \end{tcolorbox}
}
% Section numérotée : pastille orange avec le numéro + titre Archivo
\newcounter{popsec}
\newcommand{\coursec}[1]{%
  \stepcounter{popsec}\setcounter{popsub}{0}%
  \par\vspace{16pt}\noindent
  \begin{minipage}{\linewidth}
  \tikz[baseline=-0.7ex]\node[draw=popink,line width=1.6pt,fill=poporange,rounded corners=4pt,minimum size=22pt,inner sep=0]{\popdisplay\fontsize{12}{12}\selectfont\color{popcream}\thepopsec};%
  \hspace{8pt}{\popdisplay\fontsize{15}{17}\selectfont\color{popink}\MakeUppercase{#1}}\par
  \vspace{4pt}\noindent{\color{poporange}\rule{36pt}{3.6pt}}
  \end{minipage}\par\nopagebreak\vspace{8pt}
}
\newcounter{popsub}
\newcommand{\courssub}[1]{%
  \stepcounter{popsub}%
  \par\vspace{9pt}\noindent{\popxbold\fontsize{11.5}{13}\selectfont\color{popink}{\color{poporange}\thepopsec.\thepopsub}\hspace{6pt}#1}\par\nopagebreak\vspace{4pt}
}

% ============================================================
% Boîtes pédagogiques — même recette : fond blanc, bordure épaisse colorée,
% ombre dure encre, badge plein. Titre optionnel [#1].
% ============================================================
\tcbset{popbase/.style={breakable,enhanced,colback=white,boxrule=2.3pt,arc=9pt,
  shadow={3pt}{-3pt}{0pt}{popink},left=14pt,right=14pt,top=11pt,bottom=11pt,before skip=11pt,after skip=15pt,
  fontupper=\popsemi\fontsize{10.6}{14.5}\selectfont\color{popink},
  fontlower=\popsemi\fontsize{10.6}{14.5}\selectfont\color{popink}}}
% Coche dessinée (le glyphe ✓ n'existe pas dans les polices du document)
\newcommand{\popcheck}[1]{\tikz[baseline=-0.5ex]\draw[#1,line width=1.6pt,line cap=round,line join=round](-0.13,0.0)--(-0.04,-0.09)--(0.13,0.1);}
\newcommand{\popboxhead}[3]{\poppill{#1}{#2}{white}\ifx\relax#3\relax\else\hspace{6pt}{\popxbold\fontsize{10.6}{12}\selectfont\color{popink}#3}\fi\par\vspace{7pt}}

\newtcolorbox{aretenir}[1][]{popbase,colframe=popgreen,before upper={\popboxhead{À retenir}{popgreen}{#1}}}
\newtcolorbox{exemplebox}[1][]{popbase,colframe=poporange,before upper={\popboxhead{Exemple}{poporange}{#1}}}
\newtcolorbox{definition}[1][]{popbase,colframe=popblue,before upper={\popboxhead{Définition}{popblue}{#1}}}
\newtcolorbox{propriete}[1][]{popbase,colframe=poppurple,before upper={\popboxhead{Propriété}{poppurple}{#1}}}
\newtcolorbox{methode}[1][]{popbase,colframe=popgold,before upper={\popboxhead{Méthode}{popgold}{#1}}}
\newtcolorbox{attention}[1][]{popbase,colframe=poppink,before upper={\popboxhead{Attention}{poppink}{#1}}}
\newtcolorbox{experience}[1][]{popbase,colframe=popblue,colback=popblueL,before upper={\popboxhead{Expérience}{popblue}{#1}}}
% « Le savais-tu ? » : carte violette pleine (contexte, culture, Cameroun)
\newtcolorbox{savaistu}[1][]{popbase,colback=poppurple,colframe=popink,
  fontupper=\popsemi\fontsize{10.4}{14.2}\selectfont\color{white},
  before upper={\poppill{Le savais-tu\,?}{white}{poppurple}\ifx\relax#1\relax\else\hspace{6pt}{\popxbold\color{white}#1}\fi\par\vspace{7pt}}}
% Exercice résolu : énoncé en haut, \tcblower puis la solution sur fond crème
\newcounter{popexo}\newcounter{popexr}
\newtcolorbox{exoresolu}[1][]{popbase,colframe=poporange,colbacklower=popcream,
  segmentation style={popink,line width=1.2pt,dashed},
  before upper={\stepcounter{popexr}\popboxhead{Exercice résolu \thepopexr}{poporange}{#1}},
  before lower={\poppill{Solution}{popink}{popcream}\par\vspace{6pt}}}
% Exercice (non résolu en place) — la correction est dans la partie Corrigés
\newtcolorbox{exercice}[1][]{popbase,colframe=popink,
  before upper={\stepcounter{popexo}\popboxhead{Exercice \thepopexo}{popink}{#1}}}
% Corrigé
\newtcolorbox{corrige}[1][]{popbase,colframe=popgreen,colback=popgreenL,
  before upper={\popboxhead{Corrigé}{popgreen}{#1}}}
% Objectifs / prérequis / bilan
\newtcolorbox{objectifs}{popbase,colframe=popink,colback=poporangeL,
  before upper={\popboxhead{Objectifs de la leçon}{popink}{}}}
\newtcolorbox{prerequis}{popbase,colframe=popink,colback=popblueL,
  before upper={\popboxhead{Prérequis}{popblue}{}}}
\newtcolorbox{bilan}{popbase,colframe=popink,colback=white,boxrule=2.8pt,
  before upper={\popboxhead{Fiche bilan}{poporange}{L'essentiel à connaître pour l'examen}}}
% Figure encadrée avec légende
\newtcolorbox{popfigure}[1][]{enhanced,breakable=false,colback=white,colframe=popline,boxrule=1.4pt,arc=8pt,
  left=10pt,right=10pt,top=10pt,bottom=8pt,before skip=10pt,after skip=12pt,halign=center,
  lower separated=false,halign lower=center,
  fontlower=\popsemi\fontsize{9}{11.5}\selectfont\color{popmuted},
  #1}
\newcommand{\legende}[1]{\par\vspace{6pt}{\popsemi\fontsize{9}{11.5}\selectfont\color{popmuted}#1}}

% ============================================================
% Signature OnBuch+
% ============================================================
\newcommand{\poplogoinline}[1]{{\poplogo\fontsize{#1}{#1}\selectfont\color{popink}OnBuch\color{poporange}+}}
\newcommand{\signatureonbuch}{%
  \begin{tcolorbox}[enhanced,colback=popink,colframe=popink,boxrule=2.2pt,arc=10pt,
      drop shadow=poporange,left=14pt,right=14pt,top=10pt,bottom=10pt,before skip=0pt,after skip=0pt]
    \tikz[baseline=-0.7ex]\node[circle,fill=popgreen,draw=popcream,line width=1.3pt,minimum size=22pt,inner sep=0]{\popcheck{white}};%
    \hspace{9pt}\begin{minipage}[c]{0.62\linewidth}
      {\popxbold\fontsize{12}{13}\selectfont\color{popcream}Signé\ \textbullet\ {\poplogo OnBuch\color{poporange}+}}\par\vspace{2pt}
      {\raggedright\popsemi\fontsize{8}{10.5}\selectfont\color{popline}\MakeUppercase{Équipe pédagogique OnBuch+}\\\MakeUppercase{Conforme au programme officiel MINESEC}\par}
    \end{minipage}\hfill
    {\poplogo\fontsize{22}{22}\selectfont\color{popcream}OnBuch\color{poporange}+}
  \end{tcolorbox}%
}

% ============================================================
% Page de garde
% #1 titre · #2 matière · #3 niveau · #4 module · #5 n° leçon · #6 durée ·
% #7 description / objectifs (texte ou itemize)
% ============================================================
\newcommand{\pagegarde}[7]{%
  \thispagestyle{empty}
  \newgeometry{a4paper,margin=1.7cm,top=1.6cm,bottom=1.4cm}%
  \begin{tikzpicture}[remember picture,overlay]
    \fill[poporange!18] ([xshift=-3.2cm,yshift=-3.6cm]current page.north east) circle (4.6cm);
    \fill[poppurple!14] ([xshift=2.4cm,yshift=7.6cm]current page.south west) circle (3.8cm);
  \end{tikzpicture}%
  {\poplogo\fontsize{30}{30}\selectfont\color{popink}OnBuch\color{poporange}+}\hfill
  \raisebox{0.6ex}{\poppill{Cours signé \textbullet\ Premium}{popink}{popcream}}\par
  \vspace{1pt}\popsquiggle{poppurple}\par
  \vspace{8pt}
  \begin{tcolorbox}[enhanced,colback=poporange,colframe=popink,boxrule=2.8pt,arc=13pt,
      drop shadow=popink,left=18pt,right=18pt,top=12pt,bottom=13pt,after skip=10pt]
    \poppill{#2 \textbullet\ #3}{popink}{popcream}\hspace{5pt}\poppill{#4}{white}{popink}\par\vspace{8pt}
    {\popdisplay\fontsize{14}{14}\selectfont\color{popink}LEÇON #5}\par\vspace{4pt}
    {\raggedright\hyphenpenalty=10000\exhyphenpenalty=10000\popdisplay\fontsize{25}{27}\selectfont\color{popcream}\MakeUppercase{#1}\par}
    \vspace{8pt}
    {\popxbold\fontsize{9.5}{9.5}\selectfont\color{popink}\MakeUppercase{Durée conseillée : #6}}
  \end{tcolorbox}
  \begin{tcolorbox}[enhanced,colback=white,colframe=popink,boxrule=2.2pt,arc=10pt,
      drop shadow=popink,left=16pt,right=16pt,top=10pt,bottom=10pt,after skip=8pt]
    \poppill{Dans cette leçon}{poppurple}{white}\par\vspace{5pt}
    \popsemi\fontsize{9.3}{12.6}\selectfont\color{popink}#7
  \end{tcolorbox}
  \noindent\begin{minipage}[t]{0.487\linewidth}
    \begin{tcolorbox}[enhanced,colback=popgreen,colframe=popink,boxrule=2.2pt,arc=10pt,
        drop shadow=popink,left=11pt,right=11pt,top=10pt,bottom=10pt,height=3.1cm,valign=center]
      \tcbox[colback=white,colframe=popink,boxrule=1.3pt,arc=5pt,boxsep=0pt,nobeforeafter,
        left=3pt,right=3pt,top=3pt,bottom=3pt,valign=center]{\qrcode[height=1.9cm]{https://play.google.com/store/apps/details?id=com.luvvix.onbuch&hl=fr}}%
      \hspace{8pt}\begin{minipage}[c]{0.5\linewidth}
        {\popdisplay\fontsize{11}{12.5}\selectfont\color{white}\MakeUppercase{Télécharge OnBuch}}\par\vspace{4pt}
        {\raggedright\popsemi\fontsize{8.4}{11}\selectfont\color{white}Gratuit \textbullet\ Google Play\\Tuteur IA, annales, concours, résultats\par}
      \end{minipage}
    \end{tcolorbox}
  \end{minipage}\hfill
  \begin{minipage}[t]{0.487\linewidth}
    \begin{tcolorbox}[enhanced,colback=poppurple,colframe=popink,boxrule=2.2pt,arc=10pt,
        drop shadow=popink,left=11pt,right=11pt,top=10pt,bottom=10pt,height=3.1cm,valign=center]
      \tikz[baseline=-0.7ex]\node[circle,fill=white,draw=popink,line width=1.3pt,minimum size=1.6cm,inner sep=0]{\popdisplay\fontsize{24}{24}\selectfont\color{poppurple}+};%
      \hspace{8pt}\begin{minipage}[c]{0.6\linewidth}
        {\popdisplay\fontsize{11}{12.5}\selectfont\color{white}\MakeUppercase{Passe à OnBuch+}}\par\vspace{4pt}
        {\raggedright\popsemi\fontsize{8.4}{11}\selectfont\color{white}Tous les cours signés, Tuteur IA illimité, fiches PDF — onglet OnBuch+ de l'app\par}
      \end{minipage}
    \end{tcolorbox}
  \end{minipage}
  \vspace{8pt}
  \popcontenu
  \vfill
  \signatureonbuch
  \restoregeometry
  \newpage
}

% Rangée « Ce cours contient » (page de garde)
\newcommand{\poptile}[3]{% #1 couleur #2 gros texte #3 légende
  \begin{tcolorbox}[enhanced,colback=#1,colframe=popink,boxrule=2pt,arc=9pt,shadow={2.5pt}{-2.5pt}{0pt}{popink},
      left=6pt,right=6pt,top=9pt,bottom=9pt,halign=center,valign=center,height=1.75cm,width=\linewidth,nobeforeafter]
    {\popdisplay\fontsize{12.5}{13}\selectfont\color{white}\MakeUppercase{#2}}\par\vspace{3pt}
    {\centering\popsemi\fontsize{7.8}{9.5}\selectfont\color{white}#3\par}
  \end{tcolorbox}}
\newcommand{\popcontenu}{%
  \par\noindent{\popxboldls\fontsize{7.5}{7.5}\selectfont\color{popmuted}CE COURS SIGNÉ CONTIENT}\par\vspace{7pt}
  \noindent\begin{minipage}[t]{0.235\linewidth}\poptile{popblue}{Cours}{complet, illustré, pas à pas}\end{minipage}\hfill
  \begin{minipage}[t]{0.235\linewidth}\poptile{poporange}{Méthodes}{et exercices résolus}\end{minipage}\hfill
  \begin{minipage}[t]{0.235\linewidth}\poptile{poppink}{Exercices}{type examen + corrigés}\end{minipage}\hfill
  \begin{minipage}[t]{0.235\linewidth}\poptile{popgreen}{Fiche bilan}{l'essentiel à retenir}\end{minipage}\par}

% Sommaire
\newtcolorbox{sommaire}{enhanced,colback=white,colframe=popink,boxrule=2.2pt,arc=10pt,
  drop shadow=popink,left=18pt,right=18pt,top=13pt,bottom=13pt,before skip=6pt,after skip=20pt,
  before upper={\poppill{Sommaire}{poppurple}{white}\par\vspace{9pt}\popsemi\fontsize{10.6}{14.5}\selectfont\color{popink}}}

% Signature de fin de cours — poussée tout en bas de la dernière page
\newcommand{\findecours}{%
  \par\vfill\nopagebreak
  \begin{center}\popsquiggle{poporange}\end{center}\vspace{4pt}
  \signatureonbuch
}
\AtBeginDocument{\def\llbracket{\mathopen{[\mkern-3.5mu[}}\def\rrbracket{\mathclose{]\mkern-3.5mu]}}} % intervalles d'entiers (glyphe ⟦ absent de la police)
% Glyphes absents de Fira Math : redéfinis à partir de glyphes disponibles
\AtBeginDocument{%
  \def\setminus{\mathbin{\backslash}}\let\smallsetminus\setminus
  \def\vdots{\mathord{\vbox{\baselineskip3.2pt\lineskiplimit0pt\kern4pt\hbox{.}\hbox{.}\hbox{.}}}}%
  \def\ddots{\mathinner{\mkern1mu\raise6.5pt\hbox{.}\mkern2mu\raise3.6pt\hbox{.}\mkern2mu\raise0.7pt\hbox{.}\mkern1mu}}%
  \def\triangle{\mathord{\scalebox{0.9}{$\symup{\Delta}$}}}%
  \def\nabla{\mathord{\raisebox{\depth}{\scalebox{1}[-1]{$\symup{\Delta}$}}}}%
  \def\checkmark{\popcheck{popdark}}%
  \def\star{\mathbin{\ast}}\def\bigstar{\mathord{\ast}}%
  \def\diamond{\mathbin{\scalebox{0.6}{$\Diamond$}}}%
}

%% FIGFIX-BEGIN v5 — correctifs globaux des figures (docs/onbuch-generation/tools/figfix)
\usepackage{adjustbox}\usepackage{etoolbox}\tcbuselibrary{raster}
% toute figure TikZ/pgfplots plus large que la ligne est réduite à la largeur disponible
\BeforeBeginEnvironment{tikzpicture}{\begin{adjustbox}{max width=\linewidth}}
\AfterEndEnvironment{tikzpicture}{\end{adjustbox}}
% légendes pgfplots lisibles par-dessus les courbes
\pgfplotsset{every axis/.append style={legend style={fill=white,fill opacity=0.92,text opacity=1,draw=black!15,inner sep=3pt}}}
% étiquettes d'arêtes (syntaxe edge["..."]) avec fond blanc
\tikzset{every edge quotes/.append style={fill=white,inner sep=1.2pt}}
%% FIGFIX-END
\begin{document}

\pagegarde{Équations, inéquations et systèmes}{Mathématiques}{Tle A}{Relations et opérations fondamentales dans l'ensemble des nombres réels}{A01}{7 h}{Cette leçon outille l'élève pour résoudre des équations et inéquations de degré 3 dans ℝ, utiliser une inconnue auxiliaire, et résoudre des systèmes linéaires à deux ou trois inconnues. Elle développe des techniques algébriques fondamentales (factorisation, pivot de Gauss) indispensables au Baccalauréat et à la modélisation de situations concrètes.
\vspace{4pt}
\begin{multicols}{2}\small
\begin{itemize}
\item À la fin de cette leçon, je sais trouver une racine évidente d'une équation de degré 3 et factoriser le polynôme sous la forme (x − α)(ax² + bx + c) par identification ou par division euclidienne
\item À la fin de cette leçon, je sais résoudre complètement une équation de degré 3 dans ℝ
\item À la fin de cette leçon, je sais résoudre une inéquation de degré 3 à l'aide d'un tableau de signes
\item À la fin de cette leçon, je sais résoudre une équation ou une inéquation en posant une inconnue auxiliaire (bicarrée, en √x, en 1/x...)
\item À la fin de cette leçon, je sais résoudre un système linéaire 2×2 par substitution, par combinaison et par le déterminant
\item À la fin de cette leçon, je sais résoudre un système linéaire 3×3 par substitution, par combinaisons et par le pivot de Gauss
\end{itemize}\end{multicols}}

\begin{sommaire}
\begin{multicols}{2}
\begin{enumerate}
\item Équations de degré 3 : racine évidente et factorisation
\item Inéquations de degré 3 et tableaux de signes
\item Résolution par inconnue auxiliaire
\item Systèmes linéaires 2×2
\item Systèmes linéaires 3×3 et pivot de Gauss
\item Problèmes concrets et modélisation
\item Activité d'intégration
\item Méthodes et exercices résolus
\item Exercices
\item Corrigés des exercices
\item Fiche bilan
\end{enumerate}
\end{multicols}
\end{sommaire}

\begin{objectifs}
\begin{itemize}
\item À la fin de cette leçon, je sais trouver une racine évidente d'une équation de degré 3 et factoriser le polynôme sous la forme (x − α)(ax² + bx + c) par identification ou par division euclidienne
\item À la fin de cette leçon, je sais résoudre complètement une équation de degré 3 dans ℝ
\item À la fin de cette leçon, je sais résoudre une inéquation de degré 3 à l'aide d'un tableau de signes
\item À la fin de cette leçon, je sais résoudre une équation ou une inéquation en posant une inconnue auxiliaire (bicarrée, en √x, en 1/x...)
\item À la fin de cette leçon, je sais résoudre un système linéaire 2×2 par substitution, par combinaison et par le déterminant
\item À la fin de cette leçon, je sais résoudre un système linéaire 3×3 par substitution, par combinaisons et par le pivot de Gauss
\item À la fin de cette leçon, je sais discuter l'existence et l'unicité des solutions d'un système 2×2 à l'aide du déterminant
\item À la fin de cette leçon, je sais modéliser un problème concret par une équation, une inéquation ou un système et interpréter les solutions obtenues
\end{itemize}
\end{objectifs}

\begin{prerequis}
\begin{itemize}
\item Résolution des équations et inéquations du premier et du second degré dans ℝ (classe de Première)
\item Factorisation et tableau de signes d'un trinôme du second degré
\item Développement, identités remarquables et calcul algébrique
\item Notion de polynôme et évaluation d'un polynôme en un réel
\item Résolution de systèmes de deux équations à deux inconnues (classe de Seconde)
\end{itemize}
\end{prerequis}

\newpage

\coursec{Équations de degré 3 : racine évidente et factorisation}

\textbf{Situation d'accroche.} Au marché Mokolo à Yaoundé, Mme Etoundi vend trois produits : du manioc, des arachides et de l'huile de palme. Ce matin, elle a encaissé trois ventes différentes dont elle connaît les montants totaux, mais — catastrophe — elle a oublié le prix unitaire de chaque produit ! Comment les retrouver ? Quelques kilomètres plus loin, un transporteur de la ligne Yaoundé--Douala sait que le volume de sa citerne, exprimé en fonction de ses dimensions, vérifie une équation du type $x^3 - 6x^2 + 11x - 6 = 0$. Ces problèmes, très concrets, se ramènent à des \cle{équations}, des \cle{inéquations} et des \cle{systèmes} que nous allons apprendre à résoudre méthodiquement tout au long de cette leçon.

\textbf{Problématique.} Tu sais résoudre une équation du premier degré $ax + b = 0$ et une équation du second degré $ax^2 + bx + c = 0$ grâce au discriminant. Mais que faire face à une équation de degré 3 ? Existe-t-il une méthode générale accessible ? La réponse est oui, à condition de connaître une astuce fondamentale : la \cle{racine évidente}, qui permet de \emph{ramener} le degré 3 au degré 2. C'est toute l'ambition de cette première section.

\courssub{Équations polynomiales de degré 3 et notion de racine}

\textbf{Intuition.} Une équation de degré 3, c'est une équation où l'inconnue apparaît élevée à la puissance 3 (et pas plus). Géométriquement, résoudre $P(x) = 0$, c'est chercher les abscisses des points où la courbe de $P$ \emph{coupe} l'axe des abscisses. Contrairement au second degré, il n'existe pas de formule simple au programme : nous allons donc construire une stratégie.

\begin{definition}[Équation polynomiale de degré 3]
On appelle \cle{équation polynomiale de degré 3} (ou équation cubique) toute équation qui, après développement et réduction, s'écrit sous la forme
\[ P(x) = 0 \quad \text{avec} \quad P(x) = ax^3 + bx^2 + cx + d, \]
où $a$, $b$, $c$, $d$ sont des nombres réels et $a \neq 0$. Le polynôme $P$ est appelé \cle{polynôme de degré 3} ; $a$ est le \emph{coefficient dominant} et $d$ le \emph{terme constant}.
\end{definition}

\begin{exemplebox}[Reconnaître une équation de degré 3]
\begin{itemize}
\item $2x^3 - 5x^2 + x - 7 = 0$ est une équation de degré 3 ($a = 2$, $b = -5$, $c = 1$, $d = -7$).
\item $x^3 - 6x^2 + 11x - 6 = 0$ est de degré 3 ($a = 1$).
\item $(x-1)(x^2 + 3x + 2) = 0$ est aussi une équation de degré 3, car en développant on obtient $x^3 + 2x^2 - x - 2 = 0$.
\item $x^3 + 2x = x^3 - 5$ n'est \textbf{pas} de degré 3 : les termes en $x^3$ s'éliminent, il reste $2x = -5$, une équation de degré 1. La condition $a \neq 0$ s'apprécie \emph{après réduction} !
\end{itemize}
\end{exemplebox}

\begin{definition}[Racine d'un polynôme]
Soit $P$ un polynôme et $\alpha$ un nombre réel. On dit que $\alpha$ est une \cle{racine} de $P$ (ou une \cle{solution} de l'équation $P(x) = 0$) lorsque
\[ P(\alpha) = 0. \]
\end{definition}

\begin{exemplebox}[Vérifier qu'un nombre est racine]
Soit $P(x) = x^3 - 6x^2 + 11x - 6$. Testons $x = 1$ :
\[ P(1) = 1^3 - 6 \times 1^2 + 11 \times 1 - 6 = 1 - 6 + 11 - 6 = 0. \]
Donc $1$ est une racine de $P$. En revanche $P(4) = 64 - 96 + 44 - 6 = 6 \neq 0$, donc $4$ n'est pas racine.
\end{exemplebox}

\courssub{La recherche d'une racine évidente}

\textbf{Intuition.} Puisqu'il n'existe pas de formule simple pour le degré 3, la première chose à faire est de \emph{deviner} une solution en testant des valeurs simples. Au Baccalauréat, les sujets sont construits de sorte qu'il existe toujours une racine « évidente » : un petit entier facile à tester.

\begin{methode}[Trouver une racine évidente]
Pour chercher une racine de $P(x) = ax^3 + bx^2 + cx + d$ :
\begin{enumerate}
\item Tester successivement les valeurs simples : $0$, $1$, $-1$, $2$, $-2$, $3$, $-3$.
\item Si aucune ne convient, tester les \cle{diviseurs du terme constant} $d$ (positifs et négatifs) : en effet, si une racine entière $\alpha$ existe, alors $\alpha$ divise $d$ (lorsque $a = 1$).
\item Dès qu'une valeur $\alpha$ vérifie $P(\alpha) = 0$, \textbf{s'arrêter} : une seule racine évidente suffit pour factoriser.
\item Toujours \textbf{écrire la vérification} $P(\alpha) = \ldots = 0$ sur la copie : une racine annoncée sans vérification ne vaut aucun point.
\end{enumerate}
\end{methode}

\begin{exemplebox}[Recherche en pratique]
Pour $P(x) = x^3 - 6x^2 + 11x - 6$ : on teste $P(0) = -6 \neq 0$, puis $P(1) = 0$. La valeur $\alpha = 1$ est notre racine évidente. Inutile de chercher plus loin.

Pour $Q(x) = x^3 + 2x^2 - 5x - 6$ : $Q(0) = -6$, $Q(1) = -8$, $Q(-1) = 0$. La racine évidente est $\alpha = -1$ (qui divise bien le terme constant $-6$).
\end{exemplebox}

\courssub{Le théorème de factorisation}

Voici le résultat central de cette section : c'est lui qui transforme la racine évidente en outil de résolution.

\begin{propriete}[Théorème de factorisation par $(x - \alpha)$]
Soit $P$ un polynôme et $\alpha$ un nombre réel.
\[ P(\alpha) = 0 \iff \text{il existe un polynôme } Q \text{ tel que } P(x) = (x - \alpha)\,Q(x). \]
De plus, si $P$ est de degré 3, alors $Q$ est de degré 2 : on peut écrire
\[ P(x) = (x - \alpha)(ax^2 + b'x + c'), \]
où $a$ est le coefficient dominant de $P$.
\end{propriete}

\begin{experience}[Démonstration guidée du sens direct (cas du degré 3)]
Supposons que $P(x) = ax^3 + bx^2 + cx + d$ admette $\alpha$ pour racine, c'est-à-dire $P(\alpha) = 0$. Montrons qu'on peut factoriser par $(x - \alpha)$.

\textbf{Étape 1 : l'astuce de la soustraction.} Comme $P(\alpha) = 0$, on peut écrire
\[ P(x) = P(x) - P(\alpha) = a(x^3 - \alpha^3) + b(x^2 - \alpha^2) + c(x - \alpha). \]

\textbf{Étape 2 : factoriser chaque morceau.} Rappelons les identités
\[ x^2 - \alpha^2 = (x - \alpha)(x + \alpha) \quad \text{et} \quad x^3 - \alpha^3 = (x - \alpha)(x^2 + \alpha x + \alpha^2). \]
(La seconde se vérifie en développant le membre de droite.)

\textbf{Étape 3 : mettre $(x - \alpha)$ en facteur.}
\begin{align*}
P(x) &= a(x - \alpha)(x^2 + \alpha x + \alpha^2) + b(x - \alpha)(x + \alpha) + c(x - \alpha) \\
&= (x - \alpha)\left[ a(x^2 + \alpha x + \alpha^2) + b(x + \alpha) + c \right].
\end{align*}
Le crochet est bien un polynôme de degré 2, de coefficient dominant $a$. On a donc $P(x) = (x - \alpha)(ax^2 + b'x + c')$. \hfill $\square$

\textbf{Réciproque} : si $P(x) = (x - \alpha)Q(x)$, alors $P(\alpha) = (\alpha - \alpha)Q(\alpha) = 0$, donc $\alpha$ est racine de $P$.
\end{experience}

\begin{aretenir}[La stratégie de résolution du degré 3]
Pour résoudre $P(x) = 0$ où $P$ est de degré 3 :
\begin{enumerate}
\item Trouver une racine évidente $\alpha$ (et la vérifier par écrit).
\item Factoriser : $P(x) = (x - \alpha)(ax^2 + b'x + c')$.
\item Résoudre l'équation produit : $x = \alpha$ \textbf{ou} $ax^2 + b'x + c' = 0$ (discriminant).
\end{enumerate}
Le problème de degré 3 est ainsi \emph{ramené} à un problème de degré 2, que tu maîtrises déjà.
\end{aretenir}

\courssub{Méthode 1 : identification des coefficients}

\begin{methode}[Factoriser par identification des coefficients]
On cherche $b'$ et $c'$ tels que $ax^3 + bx^2 + cx + d = (x - \alpha)(ax^2 + b'x + c')$.
\begin{enumerate}
\item \textbf{Développer} le membre de droite :
\[ (x - \alpha)(ax^2 + b'x + c') = ax^3 + (b' - a\alpha)x^2 + (c' - b'\alpha)x - c'\alpha. \]
\item \textbf{Comparer} les coefficients avec ceux de $P$ : deux polynômes sont égaux si et seulement si leurs coefficients sont égaux. On obtient le système
\[ \begin{cases} b' - a\alpha = b \\ c' - b'\alpha = c \\ -c'\alpha = d \end{cases} \]
\item \textbf{Résoudre} : la première équation donne $b' = b + a\alpha$, la deuxième donne $c'$, et la troisième sert de \emph{vérification}.
\item \textbf{Conclure} en écrivant la factorisation.
\end{enumerate}
\end{methode}

\begin{exemplebox}[Identification sur notre exemple fil rouge]
Factorisons $P(x) = x^3 - 6x^2 + 11x - 6$ par $(x - 1)$. On cherche $b'$ et $c'$ tels que
\[ x^3 - 6x^2 + 11x - 6 = (x - 1)(x^2 + b'x + c'). \]
Développons : $(x - 1)(x^2 + b'x + c') = x^3 + (b' - 1)x^2 + (c' - b')x - c'$.

Par identification :
\[ \begin{cases} b' - 1 = -6 \\ c' - b' = 11 \\ -c' = -6 \end{cases} \iff \begin{cases} b' = -5 \\ c' = 11 + b' = 6 \\ c' = 6 \quad \text{(cohérent !)} \end{cases} \]
Donc $P(x) = (x - 1)(x^2 - 5x + 6)$.
\end{exemplebox}

\courssub{Méthode 2 : division euclidienne}

On peut aussi obtenir le quotient $Q(x)$ en \emph{posant la division} de $P(x)$ par $(x - \alpha)$, exactement comme une division de nombres entiers : on élimine les termes de plus haut degré un par un.

\begin{exemplebox}[Division posée de $x^3 - 6x^2 + 11x - 6$ par $(x - 1)$]
On procède terme à terme :
\begin{itemize}
\item $x^3 \div x = x^2$ ; on multiplie $x^2(x-1) = x^3 - x^2$ et on soustrait : reste $-5x^2 + 11x$.
\item $-5x^2 \div x = -5x$ ; on multiplie $-5x(x-1) = -5x^2 + 5x$ et on soustrait : reste $6x - 6$.
\item $6x \div x = 6$ ; on multiplie $6(x-1) = 6x - 6$ et on soustrait : reste $0$.
\end{itemize}
Le reste est nul (ce qui confirme que $1$ est racine) et le quotient est $x^2 - 5x + 6$.
\end{exemplebox}

\begin{popfigure}
\begin{center}
\begin{tikzpicture}[scale=0.92, every node/.style={font=\small}]
% Dividende et diviseur
\node[anchor=east] (d0) at (0,0) {$x^3 - 6x^2 + 11x - 6$};
\node[anchor=west] (v0) at (4.6,0) {$x - 1$};
\draw[popdark, thick] (4.4,0.35) -- (4.4,-2.6);
\draw[popdark, thick] (4.4,0.35) -- (7.6,0.35);
% Quotient
\node[anchor=west, popblue] at (4.6,-0.6) {$x^2 - 5x + 6$};
% Étapes
\node[anchor=east, popmuted] at (0,-0.9) {$x^3 - \phantom{1}x^2$};
\node[anchor=east, popmuted] at (0.4,-1.3) {$-$};
\node[anchor=east] at (0,-1.8) {$-5x^2 + 11x$};
\node[anchor=east, popmuted] at (0,-2.5) {$-5x^2 + \phantom{1}5x$};
\node[anchor=east, popmuted] at (0.4,-2.9) {$-$};
\node[anchor=east] at (0,-3.4) {$6x - 6$};
\node[anchor=east, popmuted] at (0,-4.1) {$6x - 6$};
\node[anchor=east, popmuted] at (0.4,-4.5) {$-$};
\node[anchor=east, popink] at (0,-5.0) {$0$};
% Annotations
\node[popblue, anchor=west] at (7.8,-0.6) {$\leftarrow$ quotient};
\node[popink, anchor=west] at (0.8,-5.0) {$\leftarrow$ reste nul : $1$ est bien racine};
\end{tikzpicture}
\end{center}
\legende{Division euclidienne posée de $x^3 - 6x^2 + 11x - 6$ par $(x-1)$ : le quotient est $x^2 - 5x + 6$ et le reste est nul.}
\end{popfigure}

\begin{attention}[Le reste doit être nul]
Si, à la fin de la division de $P(x)$ par $(x - \alpha)$, le reste n'est pas nul, c'est que $\alpha$ \textbf{n'est pas} racine de $P$ (erreur de calcul dans la racine évidente ou dans la division). Le reste nul est un excellent auto-contrôle. De même, dans l'identification, la dernière équation $-c'\alpha = d$ doit être vérifiée : si elle ne l'est pas, cherche l'erreur !
\end{attention}

\courssub{Résolution complète et discussion du nombre de solutions}

Une fois la factorisation obtenue, on applique la \cle{règle du produit nul} :
\[ P(x) = 0 \iff (x - \alpha)(ax^2 + b'x + c') = 0 \iff x = \alpha \ \text{ou}\ ax^2 + b'x + c' = 0. \]
Le nombre de solutions réelles dépend alors du \cle{discriminant} $\Delta = b'^2 - 4ac'$ du facteur du second degré.

\begin{propriete}[Nombre de solutions d'une équation de degré 3 factorisée]
Soit $P(x) = (x - \alpha)(ax^2 + b'x + c')$ et $\Delta = b'^2 - 4ac'$.
\begin{itemize}
\item Si $\Delta > 0$ : l'équation $P(x) = 0$ admet \textbf{trois solutions réelles} (distinctes si $\alpha$ n'est pas racine du trinôme) : $\alpha$, $x_1 = \dfrac{-b' - \sqrt{\Delta}}{2a}$, $x_2 = \dfrac{-b' + \sqrt{\Delta}}{2a}$.
\item Si $\Delta = 0$ : l'équation admet \textbf{deux solutions réelles} : $\alpha$ et la racine double $x_0 = \dfrac{-b'}{2a}$ (si $\alpha \neq x_0$).
\item Si $\Delta < 0$ : l'équation admet \textbf{une unique solution réelle} : $x = \alpha$.
\end{itemize}
Une équation de degré 3 à coefficients réels possède donc \emph{toujours au moins une solution réelle}.
\end{propriete}

\begin{exoresolu}[Résoudre $x^3 - 6x^2 + 11x - 6 = 0$]
Résoudre dans $\mathbb{R}$ l'équation $x^3 - 6x^2 + 11x - 6 = 0$.
\tcblower
\textbf{Étape 1 : racine évidente.} On teste $x = 1$ :
\[ P(1) = 1 - 6 + 11 - 6 = 0. \]
Donc $1$ est racine de $P$.

\textbf{Étape 2 : factorisation.} Par identification, $P(x) = (x - 1)(x^2 + b'x + c')$ donne $b' - 1 = -6$ et $-c' = -6$, soit $b' = -5$ et $c' = 6$ (vérification : $c' - b' = 6 + 5 = 11$ ✓). Ainsi
\[ P(x) = (x - 1)(x^2 - 5x + 6). \]

\textbf{Étape 3 : résolution du trinôme.} $\Delta = (-5)^2 - 4 \times 1 \times 6 = 25 - 24 = 1 > 0$. Deux racines :
\[ x_1 = \frac{5 - 1}{2} = 2 \quad \text{et} \quad x_2 = \frac{5 + 1}{2} = 3. \]

\textbf{Conclusion.} L'ensemble des solutions est $\mathcal{S} = \{1 \,;\, 2 \,;\, 3\}$. Vérification rapide : $P(2) = 8 - 24 + 22 - 6 = 0$ et $P(3) = 27 - 54 + 33 - 6 = 0$. ✓
\end{exoresolu}

\begin{popfigure}
\begin{center}
\begin{tikzpicture}
\begin{axis}[
    width=0.85\linewidth, height=6.5cm,
    axis lines=middle,
    xlabel={$x$}, ylabel={$y$},
    xmin=-0.5, xmax=4, ymin=-2.5, ymax=3,
    xtick={1,2,3}, ytick=\empty,
    grid=both, grid style={popline!40},
    tick label style={font=\small},
    label style={font=\small}
]
\addplot[popcurve, domain=0.35:3.65, samples=200]{x^3 - 6*x^2 + 11*x - 6};
\addplot[only marks, mark=*, mark size=2.2pt, popink] coordinates {(1,0) (2,0) (3,0)};
\node[popink, above, font=\small] at (axis cs:1,0) {$1$};
\node[popink, below, font=\small] at (axis cs:2,0) {$2$};
\node[popink, above, font=\small] at (axis cs:3,0) {$3$};
\node[popblue, font=\small] at (axis cs:3.45,1.6) {$y = P(x)$};
\end{axis}
\end{tikzpicture}
\end{center}
\legende{Courbe de $P(x) = x^3 - 6x^2 + 11x - 6$. Les solutions de $P(x) = 0$ sont les abscisses des points d'intersection de la courbe avec l'axe des abscisses : $1$, $2$ et $3$.}
\end{popfigure}

\begin{attention}[Les deux erreurs classiques au Bac]
\begin{itemize}
\item \textbf{Oublier le facteur du second degré.} Écrire « $P(1) = 0$ donc $\mathcal{S} = \{1\}$ » est faux : la factorisation fait apparaître une équation du second degré qu'il faut \emph{toujours} résoudre. Une équation de degré 3 peut avoir jusqu'à trois solutions !
\item \textbf{Annoncer une racine sans la vérifier.} Écrire « la racine évidente est $1$ » sans calculer $P(1)$ est une faute de rédaction sanctionnée. La vérification $P(1) = 1 - 6 + 11 - 6 = 0$ doit figurer explicitement sur la copie.
\item Piège supplémentaire : dans la factorisation $(x - \alpha)$, le signe dépend de $\alpha$. Si $\alpha = -2$, le facteur est $(x - (-2)) = (x + 2)$, et non $(x - 2)$.
\end{itemize}
\end{attention}

\begin{exoresolu}[Une équation à une seule solution]
Résoudre dans $\mathbb{R}$ l'équation $x^3 + x^2 + x - 3 = 0$.
\tcblower
\textbf{Racine évidente.} $P(1) = 1 + 1 + 1 - 3 = 0$, donc $1$ est racine.

\textbf{Factorisation.} $P(x) = (x - 1)(x^2 + b'x + c')$ avec $b' - 1 = 1$ et $-c' = -3$, soit $b' = 2$, $c' = 3$ (vérification : $c' - b' = 3 - 2 = 1$ ✓). Donc
\[ P(x) = (x - 1)(x^2 + 2x + 3). \]

\textbf{Trinôme.} $\Delta = 4 - 12 = -8 < 0$ : pas de racine réelle.

\textbf{Conclusion.} $\mathcal{S} = \{1\}$. L'équation n'admet qu'une seule solution réelle.
\end{exoresolu}

\begin{savaistu}[Les formules de Cardan : une histoire de duels mathématiques]
Au XVI\textsuperscript{e} siècle, les mathématiciens italiens Tartaglia et Cardan se sont disputé la paternité de formules générales donnant les solutions de \emph{toute} équation de degré 3, à l'aide de racines carrées et cubiques — les célèbres \emph{formules de Cardan}, publiées en 1545. Ces formules, spectaculaires mais lourdes, sont hors programme au lycée. Elles ont d'ailleurs joué un rôle historique capital : pour les manipuler, il a fallu accepter des « racines carrées de nombres négatifs », ce qui a mené à l'invention des nombres complexes ! Au Baccalauréat, pas besoin de Cardan : la \cle{racine évidente} est notre clé, et elle ouvre toutes les portes des sujets officiels.
\end{savaistu}

\begin{exercice}[À toi de jouer]
\begin{enumerate}
\item Résoudre dans $\mathbb{R}$ l'équation $x^3 - 2x^2 - 5x + 6 = 0$ (chercher une racine parmi $1$, $-2$, $3$).
\item Résoudre dans $\mathbb{R}$ l'équation $2x^3 - 3x^2 - 3x + 2 = 0$.
\item Un réservoir a un volume (en m$^3$) donné par $V(x) = x^3 - 7x^2 + 14x - 8$ où $x$ est sa profondeur en mètres. Pour quelles profondeurs le volume est-il nul ?
\end{enumerate}
\end{exercice}

\begin{corrige}[Exercice]
\textbf{1.} $P(1) = 1 - 2 - 5 + 6 = 0$, donc $1$ est racine. Identification : $P(x) = (x-1)(x^2 - x - 6)$. Discriminant : $\Delta = 1 + 24 = 25$, racines $\dfrac{1-5}{2} = -2$ et $\dfrac{1+5}{2} = 3$. Donc $\mathcal{S} = \{-2 \,;\, 1 \,;\, 3\}$.

\textbf{2.} $P(-1) = -2 - 3 + 3 + 2 = 0$, donc $-1$ est racine. On écrit $P(x) = (x+1)(2x^2 + b'x + c')$ : par identification $b' + 2 = -3$ donne $b' = -5$, et $c' = 2$ (vérification : $c' + b' = 2 - 5 = -3$ ✓). Donc $P(x) = (x+1)(2x^2 - 5x + 2)$. Discriminant : $\Delta = 25 - 16 = 9$, racines $\dfrac{5-3}{4} = \dfrac{1}{2}$ et $\dfrac{5+3}{4} = 2$. Donc $\mathcal{S} = \left\{-1 \,;\, \dfrac{1}{2} \,;\, 2\right\}$.

\textbf{3.} On résout $x^3 - 7x^2 + 14x - 8 = 0$. $V(1) = 1 - 7 + 14 - 8 = 0$, donc $1$ est racine. Factorisation : $V(x) = (x-1)(x^2 - 6x + 8)$. Discriminant : $\Delta = 36 - 32 = 4$, racines $\dfrac{6-2}{2} = 2$ et $\dfrac{6+2}{2} = 4$. Les profondeurs sont $1$ m, $2$ m et $4$ m.
\end{corrige}

\coursec{Inéquations de degré 3 et tableaux de signes}

Dans la section précédente, tu as appris à résoudre une équation de degré 3 : on cherche une racine évidente $\alpha$, on factorise $P(x) = (x-\alpha)(ax^2+bx+c)$, puis on résout le trinôme. Nous savons donc désormais trouver \emph{où} un polynôme de degré 3 s'annule. Mais dans la vie courante comme au Baccalauréat, on demande souvent autre chose : \emph{quand} une quantité est-elle positive ? À partir de quel niveau de production une entreprise est-elle bénéficiaire ? Ces questions se traduisent par des \cle{inéquations} : $P(x) \geq 0$, $P(x) < 0$, etc. Toute cette section repose sur une idée simple et puissante : \textbf{une fois le polynôme factorisé, le signe d'un produit se déduit des signes de ses facteurs}, et l'outil qui organise ce raisonnement sans erreur est le \cle{tableau de signes}.

\courssub{De l'équation à l'inéquation : l'intuition}

\begin{definition}[Inéquation polynomiale de degré 3]
Soit $P(x) = ax^3+bx^2+cx+d$ un polynôme de degré 3 ($a \neq 0$). Résoudre dans $\mathbb{R}$ une inéquation du type
\[ P(x) \geq 0, \qquad P(x) > 0, \qquad P(x) \leq 0, \qquad P(x) < 0, \]
c'est déterminer l'\cle{ensemble des valeurs} de $x$ pour lesquelles $P(x)$ est positif, négatif ou nul selon le cas. Cet ensemble s'exprime toujours comme un intervalle ou une réunion d'intervalles.
\end{definition}

L'idée directrice est la suivante. Les racines de $P$ sont les \emph{frontières} : ce sont les seuls points où le signe de $P(x)$ peut changer. Entre deux racines consécutives, le polynôme garde un signe constant. Résoudre l'inéquation revient donc à :
\begin{enumerate}
\item trouver les racines (c'est le travail de la section 1) ;
\item déterminer le signe de $P(x)$ sur chaque intervalle délimité par ces racines.
\end{enumerate}

\begin{propriete}[Changement de signe en une racine simple — admise]
Soit $P$ un polynôme de degré 3 et $\alpha$ une racine \cle{simple} de $P$ (c'est-à-dire que $(x-\alpha)$ apparaît une seule fois dans la factorisation). Alors $P(x)$ \textbf{change de signe} lorsque $x$ traverse $\alpha$.

En revanche, si $\alpha$ est une racine \cle{double} (facteur $(x-\alpha)^2$), alors $P(x)$ \textbf{garde le même signe} de part et d'autre de $\alpha$ : seul le point $x=\alpha$ annule $P$.
\end{propriete}

Cette propriété est admise en Terminale A, mais elle est très intuitive graphiquement : la courbe d'un polynôme est un tracé continu, sans saut. Pour passer d'une valeur positive à une valeur négative, elle est obligée de \emph{couper} l'axe des abscisses, et ce point de coupure est exactement une racine. Nous l'illustrerons graphiquement à la fin de la section.

\courssub{Rappels indispensables : signe des facteurs}

Tout le travail repose sur deux résultats que tu connais depuis la classe de Première. Rappelons-les avec précision, car la moindre hésitation ici ruine tout le tableau de signes.

\begin{propriete}[Signe d'un binôme du premier degré]
Le facteur $(x - \alpha)$ est :
\begin{itemize}
\item \textbf{négatif} pour $x < \alpha$ ;
\item \textbf{nul} pour $x = \alpha$ ;
\item \textbf{positif} pour $x > \alpha$.
\end{itemize}
Autrement dit, un binôme $(x-\alpha)$ est du signe de ce qui suit la racine : il s'annule en $\alpha$ et devient positif après.
\end{propriete}

\begin{propriete}[Signe d'un trinôme du second degré]
Soit $T(x) = ax^2+bx+c$ ($a\neq 0$) de discriminant $\Delta = b^2 - 4ac$.
\begin{itemize}
\item Si $\Delta > 0$ : $T$ a deux racines $x_1 < x_2$ et $T(x)$ est \textbf{du signe de $a$ à l'extérieur des racines}, du signe contraire entre les racines.
\item Si $\Delta = 0$ : $T$ a une racine double $x_0$ et $T(x)$ est \textbf{toujours du signe de $a$}, en s'annulant seulement en $x_0$.
\item Si $\Delta < 0$ : $T$ ne s'annule jamais et est \textbf{toujours du signe de $a$} sur $\mathbb{R}$.
\end{itemize}
\end{propriete}

\begin{attention}[Le cas de la racine double ne doit pas être oublié]
Quand le facteur quadratique a un discriminant nul, par exemple $P(x) = (x-1)(x-2)^2$, le facteur $(x-2)^2$ est \textbf{toujours positif ou nul}. Le signe de $P$ est donc entièrement commandé par $(x-1)$, et $P$ ne change \emph{pas} de signe en $x=2$. Ainsi, pour $P(x) \geq 0$, la solution est $[1\,;\,+\infty[$ : le point $x=2$ est déjà inclus dans cet intervalle et ne crée aucune coupure. Oublier cette subtilité conduit à découper le tableau en trop de colonnes et à conclure faux.
\end{attention}

\courssub{La méthode en quatre étapes}

\begin{methode}[Résoudre une inéquation de degré 3]
Pour résoudre $P(x) \geq 0$ (ou $>0$, $\leq 0$, $<0$) :
\begin{enumerate}
\item \textbf{Factoriser} $P(x)$ complètement : chercher une racine évidente $\alpha$, écrire $P(x)=(x-\alpha)(ax^2+bx+c)$, puis factoriser le trinôme si possible (discriminant).
\item \textbf{Repérer les racines} de chaque facteur et les ranger dans l'ordre croissant : ce sont les bornes des colonnes du tableau.
\item \textbf{Construire le tableau de signes} : une ligne par facteur, une ligne pour le produit $P(x)$ obtenue par la règle des signes. On place un $0$ sous chaque racine.
\item \textbf{Conclure} : lire les colonnes où $P(x)$ a le signe demandé et écrire l'ensemble solution sous forme d'intervalles. Crochets \textbf{fermés} si l'inégalité est large ($\geq$ ou $\leq$), crochets \textbf{ouverts} si elle est stricte ($>$ ou $<$).
\end{enumerate}
\end{methode}

\begin{attention}[Piège capital : ne jamais diviser par $(x-\alpha)$ sans discuter]
Face à une inéquation comme $(x-1)(x^2-5x+6) \geq 0$, beaucoup d'élèves sont tentés de « simplifier » en divisant les deux membres par $(x-1)$. C'est une \textbf{erreur grave} : le sens d'une inégalité change lorsqu'on multiplie ou divise par un nombre \emph{négatif}, et le signe de $(x-1)$ dépend justement de $x$ ! Il faudrait distinguer les cas $x>1$ et $x<1$, ce qui complique inutilement. La bonne attitude en Terminale A est systématique : \textbf{on ne divise jamais, on factorise et on fait un tableau de signes}. C'est plus sûr, plus rapide, et c'est ce qui est attendu au Baccalauréat.
\end{attention}

\courssub{Exemple fondateur : résoudre $x^3 - 6x^2 + 11x - 6 \geq 0$}

\begin{exoresolu}[Résolution complète d'une inéquation de degré 3]
Résoudre dans $\mathbb{R}$ l'inéquation $x^3 - 6x^2 + 11x - 6 \geq 0$.
\tcblower
\textbf{Étape 1 : factorisation.} Posons $P(x) = x^3 - 6x^2 + 11x - 6$. On teste les valeurs simples : $P(1) = 1 - 6 + 11 - 6 = 0$. Donc $1$ est racine évidente et $P(x) = (x-1)(ax^2+bx+c)$. Par identification :
\[ (x-1)(x^2+bx+c) = x^3 + (b-1)x^2 + (c-b)x - c. \]
En comparant avec $x^3 - 6x^2 + 11x - 6$ : $b-1 = -6$ donne $b=-5$ ; $-c = -6$ donne $c=6$ ; vérification : $c-b = 6-(-5) = 11$. D'où
\[ P(x) = (x-1)(x^2 - 5x + 6). \]
Le trinôme $x^2-5x+6$ a pour discriminant $\Delta = 25 - 24 = 1 > 0$, de racines $x_1 = \dfrac{5-1}{2} = 2$ et $x_2 = \dfrac{5+1}{2} = 3$. Finalement :
\[ P(x) = (x-1)(x-2)(x-3). \]

\textbf{Étape 2 : racines.} Les trois racines, rangées dans l'ordre croissant, sont $1$, $2$, $3$.

\textbf{Étape 3 : tableau de signes.} Chaque facteur $(x-\alpha)$ est négatif avant $\alpha$, positif après. On applique la règle des signes ligne par ligne :
\begin{center}
\begin{tabular}{|c|ccccccccc|}
\hline
\rowcolor{popblueL}
$x$ & $-\infty$ & & $1$ & & $2$ & & $3$ & & $+\infty$ \\
\hline
$x-1$ & & $-$ & $0$ & $+$ & $\vert$ & $+$ & $\vert$ & $+$ & \\
\hline
$x-2$ & & $-$ & $\vert$ & $-$ & $0$ & $+$ & $\vert$ & $+$ & \\
\hline
$x-3$ & & $-$ & $\vert$ & $-$ & $\vert$ & $-$ & $0$ & $+$ & \\
\hline
\rowcolor{popgreenL}
$P(x)$ & & $-$ & $0$ & $+$ & $0$ & $-$ & $0$ & $+$ & \\
\hline
\end{tabular}
\end{center}
Vérifions la ligne produit : avant $1$, trois facteurs négatifs donnent $(-)\times(-)\times(-) = -$ ; entre $1$ et $2$, un seul facteur négatif donne $+$ ; entre $2$ et $3$, deux facteurs négatifs donnent $-$ ; après $3$, tout est positif.

\textbf{Étape 4 : conclusion.} On cherche où $P(x) \geq 0$ (inégalité large, donc crochets fermés sur les racines) :
\[ \mathcal{S} = [1\,;\,2] \cup [3\,;\,+\infty[. \]
\end{exoresolu}

\begin{aretenir}[L'alternance des signes]
Quand un polynôme de degré 3 possède \textbf{trois racines simples} $x_1 < x_2 < x_3$ et un coefficient dominant $a > 0$, ses signes alternent en partant du négatif à gauche :
\[ - \;\Big|\; x_1 \;\Big|\; + \;\Big|\; x_2 \;\Big|\; - \;\Big|\; x_3 \;\Big|\; + \]
C'est une conséquence directe de la propriété de changement de signe en chaque racine simple. Si $a < 0$, tous les signes sont inversés. Ce schéma permet de vérifier instantanément ton tableau.
\end{aretenir}

\courssub{Lecture graphique : la courbe par rapport à l'axe des abscisses}

Le tableau de signes admet une traduction géométrique immédiate : dire $P(x) \geq 0$, c'est dire que le point de la courbe d'abscisse $x$ se trouve \textbf{au-dessus} de l'axe des abscisses (ou dessus). L'ensemble solution est donc l'ensemble des abscisses des points de la courbe situés au-dessus de l'axe.

\begin{popfigure}
\begin{center}
\begin{tikzpicture}
\begin{axis}[
    width=0.9\linewidth, height=7cm,
    axis lines=middle,
    xlabel={$x$}, ylabel={$y$},
    xmin=-0.5, xmax=4.5, ymin=-2.5, ymax=2.5,
    xtick={1,2,3}, ytick=\empty,
    xlabel style={anchor=west}, ylabel style={anchor=south},
    clip=false,
]
% Zones positives (vert)
\fill[popgreen!30] plot[domain=1:2, samples=100, variable=\x] (\x, {\x^3-6*\x^2+11*\x-6}) -- (axis cs:2,0) -- (axis cs:1,0) -- cycle;
\fill[popgreen!30] plot[domain=3:3.75, samples=100, variable=\x] (\x, {\x^3-6*\x^2+11*\x-6}) -- (axis cs:3.75,0) -- (axis cs:3,0) -- cycle;
% Zones négatives (rose)
\fill[poppink!30] plot[domain=0.35:1, samples=100, variable=\x] (\x, {\x^3-6*\x^2+11*\x-6}) -- (axis cs:1,0) -- (axis cs:0.35,0) -- cycle;
\fill[poppink!30] plot[domain=2:3, samples=100, variable=\x] (\x, {\x^3-6*\x^2+11*\x-6}) -- (axis cs:3,0) -- (axis cs:2,0) -- cycle;
% Courbe
\addplot[popcurve, domain=0.35:3.75, samples=300, thick] {x^3-6*x^2+11*x-6};
% Racines
\addplot[only marks, mark=*, mark size=2.5pt, popdark] coordinates {(1,0) (2,0) (3,0)};
% Annotations
\node[popdark, anchor=south west, font=\small] at (axis cs:1.4, 0.6) {$P(x)\geq 0$};
\node[popdark, anchor=north, font=\small] at (axis cs:2.5, -0.8) {$P(x)\leq 0$};
\end{axis}
\end{tikzpicture}
\end{center}
\legende{Courbe de $P(x)=x^3-6x^2+11x-6$. En vert, les zones où la courbe est au-dessus de l'axe ($P(x)\geq 0$, solution $[1;2]\cup[3;+\infty[$) ; en rose, les zones en dessous.}
\end{popfigure}

On observe concrètement la propriété admise : la courbe \emph{traverse} l'axe en chacune des trois racines simples $1$, $2$, $3$, ce qui traduit le changement de signe. Si la racine était double, la courbe se contenterait de \emph{toucher} l'axe et de repartir du même côté, comme le fait une parabole en son sommet.

\begin{exemplebox}[Vérification graphique d'un tableau de signes]
Pour l'inéquation $P(x) < 0$ avec le même polynôme, on lit sur le graphique les zones roses : la courbe est strictement sous l'axe pour $x \in \;]-\infty\,;\,1[ \;\cup\; ]2\,;\,3[$. C'est exactement ce que donne la ligne « $-$ » du tableau de signes, avec crochets ouverts car l'inégalité est stricte. Prends l'habitude de ce contrôle croisé tableau/graphique : il détecte immédiatement une erreur de signe.
\end{exemplebox}

\courssub{Un second exemple avec racine double}

\begin{exoresolu}[Cas d'une racine double]
Résoudre dans $\mathbb{R}$ l'inéquation $-x^3 + 4x^2 - 4x > 0$.
\tcblower
\textbf{Factorisation.} On met d'abord $-x$ en facteur : $-x^3+4x^2-4x = -x(x^2-4x+4) = -x(x-2)^2$.

\textbf{Racines.} $0$ (racine simple) et $2$ (racine \emph{double}).

\textbf{Tableau de signes.} Le facteur $(x-2)^2$ est un carré : toujours positif, nul seulement en $2$.
\begin{center}
\begin{tabular}{|c|ccccccc|}
\hline
\rowcolor{popblueL}
$x$ & $-\infty$ & & $0$ & & $2$ & & $+\infty$ \\
\hline
$-x$ & & $+$ & $0$ & $-$ & $\vert$ & $-$ & \\
\hline
$(x-2)^2$ & & $+$ & $\vert$ & $+$ & $0$ & $+$ & \\
\hline
\rowcolor{popgreenL}
$-x(x-2)^2$ & & $+$ & $0$ & $-$ & $0$ & $-$ & \\
\hline
\end{tabular}
\end{center}
Remarque bien : en $x=2$, le polynôme s'annule mais \textbf{ne change pas de signe} — la racine double ne fait pas traverser l'axe à la courbe.

\textbf{Conclusion.} On veut $P(x) > 0$ (strict) : $\mathcal{S} = \;]-\infty\,;\,0[$.
\end{exoresolu}

\begin{savaistu}[Les inéquations de degré 3 en économie]
Les inéquations de degré 3 ne sont pas qu'un exercice scolaire : elles interviennent naturellement en \textbf{économie} pour déterminer les \cle{seuils de rentabilité}. Quand une entreprise — par exemple une unité de transformation de manioc à Bafoussam — produit $x$ tonnes, son bénéfice est souvent modélisé par une fonction $B(x) = R(x) - C(x)$ où le coût total $C$ comporte des termes cubiques (les coûts explosent quand on sature les machines). Résoudre $B(x) \geq 0$, c'est déterminer la \emph{plage de production rentable} : produire trop peu ne couvre pas les frais fixes, produire trop fait grimper les coûts plus vite que les recettes. Les bornes de l'intervalle solution sont précisément les seuils que tout gestionnaire doit connaître avant d'investir.
\end{savaistu}

\begin{exercice}[À toi de jouer]
\begin{enumerate}
\item Résoudre dans $\mathbb{R}$ : $x^3 - 2x^2 - 5x + 6 \geq 0$ (on vérifiera que $1$ est racine).
\item Résoudre dans $\mathbb{R}$ : $x^3 - 3x - 2 < 0$ (on vérifiera que $-1$ est racine ; attention à la racine double).
\end{enumerate}
\end{exercice}

\begin{corrige}[Exercice]
\textbf{1.} $P(1) = 1 - 2 - 5 + 6 = 0$, donc $P(x) = (x-1)(x^2 - x - 6)$. Le trinôme a pour discriminant $\Delta = 1 + 24 = 25$, racines $-2$ et $3$. Donc $P(x) = (x-1)(x+2)(x-3)$, racines rangées : $-2$, $1$, $3$. Coefficient dominant positif, trois racines simples : les signes alternent $-,+,-,+$. D'où $\mathcal{S} = [-2\,;\,1] \cup [3\,;\,+\infty[$.

\textbf{2.} $P(-1) = -1 + 3 - 2 = 0$, donc $P(x) = (x+1)(x^2 - x - 2) = (x+1)(x+1)(x-2) = (x+1)^2(x-2)$. Le facteur $(x+1)^2$ est toujours positif (nul en $-1$) : le signe de $P$ est celui de $(x-2)$, avec annulation en $-1$ et $2$. On veut $P(x) < 0$ strictement : $\mathcal{S} = \;]-\infty\,;\,2[$ \textbf{privé de} $-1$, c'est-à-dire $\mathcal{S} = ]-\infty\,;\,-1[ \;\cup\; ]-1\,;\,2[$. Voilà le piège de la racine double en action : $-1$ annule $P$ et doit être exclu, même s'il ne change pas le signe.
\end{corrige}

Tu maîtrises maintenant la résolution des inéquations de degré 3 par factorisation et tableau de signes. Dans la section suivante, nous verrons comment une \emph{inconnue auxiliaire} permet de ramener des équations et inéquations plus complexes (bicarrées, avec radicaux ou valeurs absolues) à ces situations que tu sais désormais traiter parfaitement.

\coursec{Résolution par inconnue auxiliaire}

Dans la section précédente, nous avons appris à résoudre des inéquations de degré $3$ grâce à la factorisation et au tableau de signes. Mais que faire face à une équation comme $x^4 - 5x^2 + 4 = 0$, de degré $4$, ou comme $x - 3\sqrt{x} + 2 = 0$, qui mélange $x$ et sa racine carrée ? Aucune des méthodes vues jusqu'ici ne s'applique directement. La bonne nouvelle, c'est qu'une seule et même idée va nous permettre de ramener toutes ces situations à des équations du second degré, que nous maîtrisons parfaitement : le \cle{changement d'inconnue}, aussi appelé \cle{inconnue auxiliaire}.

\courssub{Le principe général : repérer, poser, résoudre, revenir, vérifier}

\textbf{L'intuition.}\quad Imagine un commerçant de Douala qui note ses recettes dans un cahier en écrivant toujours la même longue formule. Pour aller plus vite, il décide de l'abréger par une seule lettre, fait tous ses calculs avec cette lettre, puis, à la fin, « traduit » son résultat dans la formule d'origine. C'est exactement ce que nous allons faire : quand une même expression en $x$ apparaît \emph{plusieurs fois} dans une équation, on la remplace provisoirement par une lettre $X$, on résout l'équation (devenue plus simple) en $X$, puis on revient à $x$.

\begin{definition}[Inconnue auxiliaire]
Soit une équation (ou inéquation) d'inconnue $x$ dans laquelle une même expression $\varphi(x)$ apparaît plusieurs fois. Poser l'\cle{inconnue auxiliaire} $X = \varphi(x)$ consiste à :
\begin{itemize}
\item réécrire l'équation en fonction de $X$ uniquement ;
\item résoudre cette nouvelle équation, dite \emph{équation auxiliaire} ;
\item résoudre ensuite, pour chaque valeur $X_0$ trouvée, l'équation $\varphi(x) = X_0$.
\end{itemize}
L'ensemble des solutions de l'équation initiale est la réunion des solutions de toutes les équations $\varphi(x) = X_0$ \emph{qui admettent des solutions réelles}.
\end{definition}

\begin{attention}[La condition d'existence, le piège numéro 1]
Chaque changement d'inconnue impose une \cle{condition} sur $X$, héritée de la nature de $\varphi(x)$ :
\begin{itemize}
\item si $X = x^2$, alors $X \geq 0$ ;
\item si $X = \sqrt{x}$, alors $x \geq 0$ \textbf{et} $X \geq 0$ ;
\item si $X = \dfrac{1}{x}$, alors $x \neq 0$ et $X \neq 0$.
\end{itemize}
Une valeur de $X$ qui viole sa condition (par exemple $X = -3$ alors que $X = x^2$) ne donne \textbf{aucune} solution en $x$ : l'oublier fait accepter des \emph{solutions fantômes}, erreur très sanctionnée au Baccalauréat. De même, il faut toujours vérifier que les valeurs finales de $x$ respectent les \cle{valeurs interdites} de l'énoncé (dénominateurs nuls, radicandes négatifs).
\end{attention}

\begin{methode}[Les 5 étapes du changement d'inconnue]
\begin{enumerate}
\item \textbf{Repérer} l'expression $\varphi(x)$ qui se répète, et écrire les \textbf{conditions d'existence} de l'équation (valeurs interdites).
\item \textbf{Poser} $X = \varphi(x)$ et préciser immédiatement la condition portant sur $X$.
\item \textbf{Résoudre} l'équation auxiliaire en $X$ (souvent du second degré : discriminant $\Delta = b^2 - 4ac$).
\item \textbf{Revenir} à $x$ : pour chaque solution $X_0$ \emph{acceptable}, résoudre $\varphi(x) = X_0$.
\item \textbf{Vérifier} que chaque $x$ obtenu respecte les conditions d'existence, puis conclure par l'ensemble $\mathcal{S}$.
\end{enumerate}
\end{methode}

\begin{popfigure}
\begin{tikzpicture}[
  node distance=0.55cm,
  etape/.style={rectangle, rounded corners=3pt, draw=popblue, fill=popblueL, text width=4.6cm, align=center, font=\small, inner sep=5pt},
  cond/.style={rectangle, rounded corners=3pt, draw=poporange, fill=poporangeL, text width=4.2cm, align=center, font=\footnotesize, inner sep=4pt},
  fleche/.style={-{Stealth[length=2.5mm]}, thick, popdark}
]
\node[etape] (e1) {\textbf{1. Repérer} l'expression répétée $\varphi(x)$};
\node[etape, below=of e1] (e2) {\textbf{2. Poser} $X = \varphi(x)$};
\node[etape, below=of e2] (e3) {\textbf{3. Résoudre} l'équation en $X$};
\node[etape, below=of e3] (e4) {\textbf{4. Revenir} à $x$ : résoudre $\varphi(x) = X_0$};
\node[etape, below=of e4] (e5) {\textbf{5. Vérifier} et conclure : $\mathcal{S} = \{\dots\}$};
\node[cond, right=1.1cm of e2, align=center] (c2) {condition sur $X$\\ ($X \geq 0$, $X \neq 0$\dots)};
\node[cond, right=1.1cm of e4, align=center] (c4) {rejeter les $X_0$ interdits\\ (solutions fantômes)};
\node[cond, right=1.1cm of e5, align=center] (c5) {valeurs interdites\\ de l'énoncé};
\draw[fleche] (e1) -- (e2);
\draw[fleche] (e2) -- (e3);
\draw[fleche] (e3) -- (e4);
\draw[fleche] (e4) -- (e5);
\draw[fleche, poporange, dashed] (c2.west) -- (e2.east);
\draw[fleche, poporange, dashed] (c4.west) -- (e4.east);
\draw[fleche, poporange, dashed] (c5.west) -- (e5.east);
\end{tikzpicture}
\legende{Organigramme de la démarche : repérer $\to$ poser $\to$ résoudre $\to$ revenir $\to$ vérifier. À chaque étape clé correspond un contrôle des conditions.}
\end{popfigure}

\courssub{Les équations bicarrées : $X = x^2$}

\begin{definition}[Équation bicarrée]
On appelle \cle{équation bicarrée} toute équation de la forme
\[ ax^4 + bx^2 + c = 0 \qquad (a \neq 0). \]
Elle ne contient que des puissances \emph{paires} de $x$ : en posant $X = x^2$ (avec $X \geq 0$), elle se ramène à l'équation du second degré $aX^2 + bX + c = 0$.
\end{definition}

\begin{propriete}[Nombre de solutions d'une bicarrée]
Après résolution de $aX^2 + bX + c = 0$ :
\begin{itemize}
\item toute racine $X_0 > 0$ donne \textbf{deux} solutions : $x = \sqrt{X_0}$ et $x = -\sqrt{X_0}$ ;
\item la racine $X_0 = 0$ donne \textbf{une} solution : $x = 0$ ;
\item toute racine $X_0 < 0$ ne donne \textbf{aucune} solution (un carré n'est jamais négatif dans $\mathbb{R}$).
\end{itemize}
Une équation bicarrée possède donc entre $0$ et $4$ solutions réelles.
\end{propriete}

\begin{exemplebox}[Résolution de $x^4 - 5x^2 + 4 = 0$]
L'expression $x^2$ apparaît deux fois : on pose $X = x^2$, avec la condition $X \geq 0$. L'équation devient :
\[ X^2 - 5X + 4 = 0. \]
Discriminant : $\Delta = (-5)^2 - 4 \times 1 \times 4 = 25 - 16 = 9 > 0$. D'où :
\[ X_1 = \frac{5 - 3}{2} = 1 \qquad \text{et} \qquad X_2 = \frac{5 + 3}{2} = 4. \]
Les deux racines sont positives : toutes deux sont acceptables. Retour à $x$ :
\begin{align*}
x^2 = 1 &\iff x = -1 \ \text{ou}\ x = 1 ;\\
x^2 = 4 &\iff x = -2 \ \text{ou}\ x = 2.
\end{align*}
Conclusion : $\mathcal{S} = \{-2\,;\,-1\,;\,1\,;\,2\}$.
\end{exemplebox}

\begin{attention}[Solutions fantômes]
Si l'on avait trouvé, par exemple, $X_1 = -1$ et $X_2 = 4$, il aurait fallu \emph{rejeter} $X_1$ : l'équation $x^2 = -1$ n'a pas de solution réelle. Écrire « $x = \pm\sqrt{-1}$ » ou « $x = \pm 1$ » serait une double faute. La condition $X \geq 0$ se vérifie \textbf{avant} de revenir à $x$.
\end{attention}

\courssub{Équations avec radicaux : $X = \sqrt{x}$ ou $X = \sqrt[3]{x}$}

Quand une équation mélange $x$ et $\sqrt{x}$, on remarque que $x = \left(\sqrt{x}\right)^2$ : tout s'exprime alors en fonction de $X = \sqrt{x}$.

\begin{methode}[Équations du type $ax + b\sqrt{x} + c = 0$]
\begin{enumerate}
\item Condition d'existence : $x \geq 0$ (présence de $\sqrt{x}$).
\item Poser $X = \sqrt{x}$, donc $X \geq 0$ et $x = X^2$.
\item Résoudre $aX^2 + bX + c = 0$.
\item Ne retenir que les racines $X_0 \geq 0$, puis calculer $x = X_0^2$.
\item Conclure.
\end{enumerate}
Pour une équation faisant intervenir $\sqrt[3]{x}$, on pose $X = \sqrt[3]{x}$ \emph{sans condition de signe} : la racine cubique existe pour tout réel, et $x = X^3$.
\end{methode}

\begin{exemplebox}[Résolution de $x - 3\sqrt{x} + 2 = 0$]
\textbf{Condition d'existence :} $x \geq 0$. Posons $X = \sqrt{x}$, avec $X \geq 0$ ; alors $x = X^2$ et l'équation devient :
\[ X^2 - 3X + 2 = 0. \]
Discriminant : $\Delta = 9 - 8 = 1 > 0$, d'où $X_1 = \dfrac{3-1}{2} = 1$ et $X_2 = \dfrac{3+1}{2} = 2$. Les deux sont positives : acceptables.
\begin{align*}
\sqrt{x} = 1 &\iff x = 1^2 = 1 ;\\
\sqrt{x} = 2 &\iff x = 2^2 = 4.
\end{align*}
Les deux valeurs vérifient $x \geq 0$. Conclusion : $\mathcal{S} = \{1\,;\,4\}$. Vérification rapide : $1 - 3 + 2 = 0$ et $4 - 6 + 2 = 0$.
\end{exemplebox}

\courssub{Équations avec dénominateurs : $X = \dfrac{1}{x}$}

Face à une équation contenant $\dfrac{1}{x}$ et $\dfrac{1}{x^2}$, on pose $X = \dfrac{1}{x}$. Le point capital, ici, est la chasse aux \cle{valeurs interdites} \emph{avant} tout calcul.

\begin{exoresolu}[Une équation fractionnaire]
Résoudre dans $\mathbb{R}$ l'équation $\dfrac{1}{x^2} - \dfrac{3}{x} + 2 = 0$.
\tcblower
\textbf{Étape 1 — Valeurs interdites.} Le dénominateur $x$ (et $x^2$) impose $x \neq 0$.

\textbf{Étape 2 — Changement d'inconnue.} Posons $X = \dfrac{1}{x}$ (donc $X \neq 0$). Alors $\dfrac{1}{x^2} = X^2$ et l'équation s'écrit :
\[ X^2 - 3X + 2 = 0. \]

\textbf{Étape 3 — Résolution.} $\Delta = 9 - 8 = 1$, d'où $X_1 = 1$ et $X_2 = 2$. Aucune n'est nulle : acceptables.

\textbf{Étape 4 — Retour à $x$.}
\[ \frac{1}{x} = 1 \iff x = 1 \qquad ; \qquad \frac{1}{x} = 2 \iff x = \frac{1}{2}. \]

\textbf{Étape 5 — Vérification.} $1$ et $\dfrac{1}{2}$ sont tous deux non nuls : aucune valeur interdite.

\textbf{Conclusion :} $\mathcal{S} = \left\{\dfrac{1}{2}\,;\,1\right\}$.
\end{exoresolu}

\courssub{Équations symétriques : $X = x + \dfrac{1}{x}$ (complément utile au Bac)}

Certaines équations présentent une \emph{symétrie} remarquable : elles associent $x^2 + \dfrac{1}{x^2}$ et $x + \dfrac{1}{x}$. L'identité suivante est la clé :
\[ \left(x + \frac{1}{x}\right)^2 = x^2 + 2 + \frac{1}{x^2} \qquad \text{donc} \qquad x^2 + \frac{1}{x^2} = X^2 - 2 \ \text{avec}\ X = x + \frac{1}{x}. \]

\begin{propriete}[Changement $X = x + \dfrac{1}{x}$]
Pour $x \neq 0$, en posant $X = x + \dfrac{1}{x}$, on a $x^2 + \dfrac{1}{x^2} = X^2 - 2$. Toute équation de la forme
\[ a\left(x^2 + \frac{1}{x^2}\right) + b\left(x + \frac{1}{x}\right) + c = 0 \]
se ramène alors au second degré : $a(X^2 - 2) + bX + c = 0$, soit $aX^2 + bX + (c - 2a) = 0$.
\end{propriete}

\begin{exemplebox}[Résolution de $2\left(x^2 + \dfrac{1}{x^2}\right) - 5\left(x + \dfrac{1}{x}\right) + 2 = 0$]
Valeur interdite : $x = 0$. Posons $X = x + \dfrac{1}{x}$ ; alors $x^2 + \dfrac{1}{x^2} = X^2 - 2$ et l'équation devient :
\[ 2(X^2 - 2) - 5X + 2 = 0 \iff 2X^2 - 5X - 2 = 0. \]
Discriminant : $\Delta = 25 + 16 = 41$, d'où $X_1 = \dfrac{5 - \sqrt{41}}{4}$ et $X_2 = \dfrac{5 + \sqrt{41}}{4}$.
Pour chaque $X_i$, on résout $x + \dfrac{1}{x} = X_i$, c'est-à-dire $x^2 - X_i\,x + 1 = 0$. Le discriminant de cette équation en $x$ est $\Delta_x = X_i^2 - 4$.
\begin{itemize}
\item Pour $X_1 = \frac{5 - \sqrt{41}}{4} \approx -0,35$, on a $\Delta_x < 0$ : pas de solution réelle.
\item Pour $X_2 = \frac{5 + \sqrt{41}}{4} \approx 2,85$, on a $\Delta_x > 0$ : deux solutions réelles $x = \frac{X_2 \pm \sqrt{X_2^2 - 4}}{2}$.
\end{itemize}
Ce type d'équation est un classique des sujets complémentaires : retiens surtout l'identité $x^2 + \dfrac{1}{x^2} = X^2 - 2$.
\end{exemplebox}

\courssub{Inéquations par changement d'inconnue}

La démarche est identique, avec une vigilance supplémentaire : on résout l'inéquation auxiliaire (tableau de signes du trinôme), on obtient un \emph{intervalle} pour $X$, puis on « traduit » cet intervalle en intervalles pour $x$.

\begin{exoresolu}[Une inéquation bicarrée]
Résoudre dans $\mathbb{R}$ l'inéquation $x^4 - 5x^2 + 4 \leq 0$.
\tcblower
Posons $X = x^2$ ($X \geq 0$) : l'inéquation devient $X^2 - 5X + 4 \leq 0$.

Les racines du trinôme sont $X_1 = 1$ et $X_2 = 4$ (calculées plus haut). Un trinôme est du signe de $a$ (ici $a = 1 > 0$) \emph{à l'extérieur} des racines, et du signe contraire \emph{entre} les racines. Donc :
\[ X^2 - 5X + 4 \leq 0 \iff 1 \leq X \leq 4. \]
Cet intervalle est bien inclus dans $[0\,;\,+\infty[$ : la condition $X \geq 0$ est respectée. Retour à $x$ :
\[ 1 \leq x^2 \leq 4. \]
Or $x^2 \geq 1 \iff x \leq -1$ ou $x \geq 1$, et $x^2 \leq 4 \iff -2 \leq x \leq 2$. L'intersection de ces deux conditions donne :
\[ \mathcal{S} = [-2\,;\,-1] \cup [1\,;\,2]. \]
\end{exoresolu}

\begin{attention}[Inégalités et retour à $x$]
Deux pièges classiques au retour à l'inconnue initiale :
\begin{itemize}
\item Traduire $1 \leq x^2 \leq 4$ par « $1 \leq x \leq 2$ » en oubliant la branche négative : c'est faux, car $x^2$ n'est pas monotone sur $\mathbb{R}$. Il faut toujours raisonner avec $|x|$ ou distinguer les deux branches.
\item Oublier d'intersecter avec la condition sur $X$ : si l'on avait trouvé $-1 \leq X \leq 4$, il aurait fallu retenir $0 \leq X \leq 4$ avant de revenir à $x$.
\end{itemize}
\end{attention}

\courssub{Synthèse : le retour à l'inconnue initiale}

Le tableau suivant récapitule les changements d'inconnue au programme et leurs conditions.

\begin{center}
\begin{tabularx}{\linewidth}{|l|l|X|}
\hline
\rowcolor{popblueL} \textbf{Situation} & \textbf{Changement} & \textbf{Conditions et retour à $x$} \\
\hline
$ax^4 + bx^2 + c = 0$ & $X = x^2$ & $X \geq 0$ ; si $X_0 > 0$ : $x = \pm\sqrt{X_0}$ \\
\hline
$ax + b\sqrt{x} + c = 0$ & $X = \sqrt{x}$ & $x \geq 0$, $X \geq 0$ ; $x = X_0^2$ \\
\hline
Équation en $\sqrt[3]{x}$ & $X = \sqrt[3]{x}$ & aucune condition de signe ; $x = X_0^3$ \\
\hline
Équation en $\dfrac{1}{x}$, $\dfrac{1}{x^2}$ & $X = \dfrac{1}{x}$ & $x \neq 0$, $X \neq 0$ ; $x = \dfrac{1}{X_0}$ \\
\hline
$a\left(x^2 + \frac{1}{x^2}\right) + b\left(x + \frac{1}{x}\right) + c = 0$ & $X = x + \frac{1}{x}$ & $x \neq 0$ ; $x^2 + \frac{1}{x^2} = X^2 - 2$ \\
\hline
\end{tabularx}
\end{center}

\begin{aretenir}[L'essentiel]
\begin{itemize}
\item Le changement d'inconnue transforme une équation compliquée en équation du second degré, à condition qu'\textbf{une même expression se répète}.
\item On écrit \textbf{toujours} la condition portant sur $X$ au moment où on le pose.
\item On ne revient à $x$ que pour les valeurs de $X$ \textbf{acceptables}, puis on vérifie les valeurs interdites de l'énoncé.
\item Pour une inéquation, on traduit l'\textbf{intervalle} solution en $X$ en intervalles en $x$, sans oublier les branches négatives.
\end{itemize}
\end{aretenir}

\begin{savaistu}[Une idée qui va très loin]
Le changement de variable est l'une des idées les plus puissantes de toutes les mathématiques. C'est lui qui permet, en classe de Terminale C puis à l'université, de calculer des \emph{intégrales} compliquées (on parle de « changement de variable dans une intégrale »), de résoudre des \emph{équations différentielles} qui décrivent la propagation d'une épidémie ou la charge d'un condensateur, et même de compresser les données transmises par les réseaux MTN ou Orange. Les ingénieurs de CAMTEL qui modélisent le trafic internet posent, sans le dire, des « $X = \varphi(x)$ » à longueur de journée. En maîtrisant cette technique dès maintenant, tu t'approprie un outil que les mathématiciens utilisent depuis plus de trois siècles — c'est déjà grâce à un changement d'inconnue que Cardan et Ferrari ont résolu, au XVI\textsuperscript{e} siècle, les équations de degré $3$ et $4$ !
\end{savaistu}

\begin{exercice}[Pour t'entraîner]
\begin{enumerate}
\item Résoudre dans $\mathbb{R}$ : $x^4 - 13x^2 + 36 = 0$.
\item Résoudre dans $\mathbb{R}$ : $2x + \sqrt{x} - 1 = 0$.
\item Résoudre dans $\mathbb{R}$ : $\dfrac{2}{x^2} + \dfrac{1}{x} - 1 = 0$.
\item Résoudre dans $\mathbb{R}$ : $x^4 - 3x^2 - 4 \geq 0$.
\end{enumerate}
\end{exercice}

\begin{corrige}[Exercices d'entraînement]
\textbf{1.} Posons $X = x^2$ ($X \geq 0$) : $X^2 - 13X + 36 = 0$, $\Delta = 169 - 144 = 25$, d'où $X_1 = 4$ et $X_2 = 9$, toutes deux positives. Alors $x = \pm 2$ ou $x = \pm 3$ : $\mathcal{S} = \{-3\,;\,-2\,;\,2\,;\,3\}$.

\textbf{2.} Condition : $x \geq 0$. Posons $X = \sqrt{x}$ ($X \geq 0$) : $2X^2 + X - 1 = 0$, $\Delta = 1 + 8 = 9$, d'où $X_1 = \dfrac{-1-3}{4} = -1$ (rejetée car $X \geq 0$) et $X_2 = \dfrac{1}{2}$. Alors $\sqrt{x} = \dfrac{1}{2}$, soit $x = \dfrac{1}{4}$ : $\mathcal{S} = \left\{\dfrac{1}{4}\right\}$.

\textbf{3.} Valeur interdite : $x = 0$. Posons $X = \dfrac{1}{x}$ : $2X^2 + X - 1 = 0$, d'où $X_1 = -1$ et $X_2 = \dfrac{1}{2}$ (aucune nulle). Alors $x = -1$ ou $x = 2$ : $\mathcal{S} = \{-1\,;\,2\}$.

\textbf{4.} Posons $X = x^2$ : $X^2 - 3X - 4 \geq 0$. Racines : $X_1 = -1$, $X_2 = 4$. Le trinôme est positif à l'extérieur des racines : $X \leq -1$ ou $X \geq 4$. Comme $X = x^2 \geq 0$, seul $X \geq 4$ convient, soit $x^2 \geq 4 \iff x \leq -2$ ou $x \geq 2$ : $\mathcal{S} = \left]-\infty\,;\,-2\right] \cup \left[2\,;\,+\infty\right[$.
\end{corrige}

\coursec{Systèmes linéaires $2\times 2$}

Dans la section précédente, tu as appris à résoudre des équations et des inéquations en introduisant une \cle{inconnue auxiliaire} : une seule inconnue, mais sous une forme parfois déguisée. Nous franchissons maintenant une étape décisive : que faire lorsqu'un problème fait intervenir \emph{deux} quantités inconnues liées par \emph{deux} conditions ?

Imagine un vendeur du marché Mokolo à Yaoundé : il a vendu au total 12 articles, des cahiers et des stylos, pour une recette de 4 000 FCFA, sachant qu'un cahier coûte 500 FCFA et un stylo 250 FCFA. Deux inconnues (le nombre de cahiers, le nombre de stylos), deux informations, deux équations. C'est exactement l'objet de cette section : les \cle{systèmes linéaires} à deux équations et deux inconnues, outil fondamental du programme de Terminale A et classique incontournable du Baccalauréat.

\courssub{Définitions et vocabulaire}

\begin{definition}[Système linéaire $2\times 2$]
On appelle \cle{système de deux équations linéaires à deux inconnues} tout système de la forme
\[
\begin{cases} ax + by = e \\ cx + dy = f \end{cases}
\]
où $a$, $b$, $c$, $d$, $e$, $f$ sont des nombres réels donnés (les \cle{coefficients}), avec $(a,b) \neq (0,0)$ et $(c,d) \neq (0,0)$, et où $x$ et $y$ sont les \cle{inconnues}.
\end{definition}

\begin{definition}[Solution d'un système]
On appelle \cle{solution} du système tout couple $(x\,;y)$ de réels qui vérifie \emph{simultanément} les deux équations. Résoudre le système, c'est déterminer l'ensemble $\mathcal{S}$ de tous les couples solutions.
\end{definition}

\begin{attention}[Les deux équations à la fois]
Un couple qui vérifie une seule des deux équations n'est \textbf{pas} une solution. Par exemple, le couple $(6\,;0)$ vérifie $2x + 3y = 12$ mais pas $x - y = 1$ : ce n'est pas une solution du système $\begin{cases} 2x+3y=12 \\ x-y=1 \end{cases}$. La vérification doit porter sur les \emph{deux} équations.
\end{attention}

\begin{definition}[Systèmes équivalents]
Deux systèmes sont dits \cle{équivalents} lorsqu'ils admettent exactement le même ensemble de solutions. On obtient un système équivalent en :
\begin{itemize}
\item échangeant deux équations ;
\item multipliant une équation par un réel non nul ;
\item ajoutant à une équation un multiple d'une autre équation ;
\item remplaçant une inconnue par son expression tirée d'une autre équation (substitution).
\end{itemize}
Toutes les méthodes de résolution reposent sur ces transformations : on transforme le système en systèmes équivalents de plus en plus simples, jusqu'à lire directement la solution.
\end{definition}

\courssub{Interprétation géométrique : intersection de deux droites}

Voici l'intuition qui éclaire toute la théorie. L'équation $ax + by = e$ (avec $(a,b)\neq(0,0)$) est l'équation d'une \cle{droite} du plan muni d'un repère. Un couple $(x\,;y)$ est solution du système si et seulement si le point $M(x\,;y)$ appartient \emph{aux deux droites} à la fois : les solutions du système sont donc les \textbf{points d'intersection} des deux droites.

Or, dans le plan, deux droites sont dans l'une des trois situations suivantes :

\begin{itemize}
\item \textbf{Droites sécantes} : elles se coupent en un unique point $\Rightarrow$ le système admet une \textbf{solution unique} ;
\item \textbf{Droites strictement parallèles} : elles ne se rencontrent jamais $\Rightarrow$ le système n'admet \textbf{aucune solution} ($\mathcal{S} = \emptyset$) ;
\item \textbf{Droites confondues} : tous leurs points sont communs $\Rightarrow$ le système admet une \textbf{infinité de solutions}.
\end{itemize}

\begin{popfigure}
\begin{center}
\begin{tikzpicture}[scale=0.85, >=Stealth]
% --- Cas 1 : sécantes ---
\begin{scope}[shift={(0,0)}]
\draw[->, popmuted] (-0.3,0) -- (3.4,0) node[below left] {\scriptsize $x$};
\draw[->, popmuted] (0,-0.4) -- (0,3.2) node[below left] {\scriptsize $y$};
\draw[popblue, thick, domain=-0.2:3.2] plot (\x, {0.9*\x+0.2}) node[right] {\scriptsize $d_1$};
\draw[poporange, thick, domain=-0.2:3.2] plot (\x, {-0.7*\x+3.4}) node[right] {\scriptsize $d_2$};
\fill[popink] (2,2) circle (2pt) node[above right] {\scriptsize $(x_0\,;y_0)$};
\node[popdark] at (1.6,-0.9) {\scriptsize \textbf{Sécantes : solution unique}};
\end{scope}
% --- Cas 2 : parallèles ---
\begin{scope}[shift={(5.2,0)}]
\draw[->, popmuted] (-0.3,0) -- (3.4,0) node[below left] {\scriptsize $x$};
\draw[->, popmuted] (0,-0.4) -- (0,3.2) node[below left] {\scriptsize $y$};
\draw[popblue, thick, domain=-0.2:3.2] plot (\x, {0.7*\x+0.4}) node[right] {\scriptsize $d_1$};
\draw[poporange, thick, domain=-0.2:3.2] plot (\x, {0.7*\x+1.7}) node[right] {\scriptsize $d_2$};
\node[popdark] at (1.6,-0.9) {\scriptsize \textbf{Parallèles : aucune solution}};
\end{scope}
% --- Cas 3 : confondues ---
\begin{scope}[shift={(10.4,0)}]
\draw[->, popmuted] (-0.3,0) -- (3.4,0) node[below left] {\scriptsize $x$};
\draw[->, popmuted] (0,-0.4) -- (0,3.2) node[below left] {\scriptsize $y$};
\draw[popblue, very thick, domain=-0.2:3.2] plot (\x, {0.6*\x+0.7}) node[right] {\scriptsize $d_1 = d_2$};
\node[popdark] at (1.6,-0.9) {\scriptsize \textbf{Confondues : infinité de solutions}};
\end{scope}
\end{tikzpicture}
\end{center}
\legende{Les trois positions relatives de deux droites du plan et le nombre de solutions du système associé.}
\end{popfigure}

\begin{aretenir}[À retenir]
Un système linéaire $2\times 2$ possède \emph{toujours} l'une de ces trois issues : une solution unique, aucune solution, ou une infinité de solutions. Il ne peut jamais avoir exactement deux ou trois solutions.
\end{aretenir}

\courssub{Méthode de substitution}

L'idée est enfantine : si l'une des équations permet d'exprimer facilement une inconnue en fonction de l'autre, on \emph{substitue} cette expression dans la seconde équation, qui ne contient alors plus qu'une seule inconnue.

\begin{methode}[Résolution par substitution]
\begin{enumerate}
\item Choisir l'équation et l'inconnue les plus simples (idéalement, une inconnue de coefficient $1$ ou $-1$).
\item Exprimer cette inconnue en fonction de l'autre (par exemple $x = \ldots$).
\item Remplacer (substituer) cette expression dans l'autre équation : on obtient une équation du premier degré à une inconnue.
\item Résoudre cette équation, puis revenir à l'expression de l'étape 2 pour obtenir la seconde inconnue.
\item Vérifier le couple obtenu dans les deux équations du système initial.
\end{enumerate}
\end{methode}

\courssub{Méthode de combinaison (ou d'addition)}

L'idée : multiplier les équations par des nombres bien choisis de sorte qu'en \emph{ajoutant} les deux équations membre à membre, l'une des inconnues disparaisse. C'est la méthode favorite des correcteurs du Bac car elle est rapide et systématique.

\begin{methode}[Résolution par combinaison]
\begin{enumerate}
\item Pour éliminer $x$ : multiplier la première équation par $c$ (coefficient de $x$ dans la seconde) et la seconde par $-a$ (ou l'inverse), de façon à obtenir des coefficients opposés.
\item Additionner membre à membre : il reste une équation en $y$ seul. Résoudre.
\item Recommencer en éliminant $y$ pour obtenir $x$ directement (ou substituer la valeur de $y$ trouvée).
\item Vérifier le couple obtenu dans les deux équations initiales.
\end{enumerate}
\end{methode}

\courssub{Méthode du déterminant et formules de Cramer}

\begin{definition}[Déterminant du système]
Le \cle{déterminant} du système $\begin{cases} ax + by = e \\ cx + dy = f \end{cases}$ est le nombre réel
\[
D = \begin{vmatrix} a & b \\ c & d \end{vmatrix} = ad - bc.
\]
\end{definition}

\begin{propriete}[Théorème — formules de Cramer]
Si $D = ad - bc \neq 0$, alors le système $\begin{cases} ax + by = e \\ cx + dy = f \end{cases}$ admet une \textbf{solution unique}, donnée par les \cle{formules de Cramer} :
\[
x = \frac{D_x}{D} = \frac{ed - bf}{ad - bc}, \qquad y = \frac{D_y}{D} = \frac{af - ec}{ad - bc},
\]
où $D_x = \begin{vmatrix} e & b \\ f & d \end{vmatrix}$ et $D_y = \begin{vmatrix} a & e \\ c & f \end{vmatrix}$ s'obtiennent en remplaçant dans $D$ la colonne de l'inconnue cherchée par la colonne des seconds membres. (Démonstration exigible : elle se fait par combinaison.)
\end{propriete}

\begin{experience}[Démonstration guidée des formules de Cramer]
On suppose $D = ad - bc \neq 0$.
\begin{enumerate}
\item \textbf{Éliminer $y$.} Multiplions la première équation par $d$ et la seconde par $-b$, puis additionnons :
\[ (ad - bc)\,x = ed - bf, \quad \text{c'est-à-dire} \quad D\,x = D_x. \]
\item \textbf{Éliminer $x$.} Multiplions la première équation par $-c$ et la seconde par $a$, puis additionnons :
\[ (ad - bc)\,y = af - ec, \quad \text{c'est-à-dire} \quad D\,y = D_y. \]
\item \textbf{Conclure.} Comme $D \neq 0$, on peut diviser : $x = \dfrac{D_x}{D}$ et $y = \dfrac{D_y}{D}$. Le système admet donc \emph{au plus} une solution. En reportant ce couple dans les deux équations (calcul direct), on vérifie qu'il convient : la solution existe et est unique. $\blacksquare$
\end{enumerate}
\end{experience}

\begin{propriete}[Discussion selon le signe de $D$]
\begin{itemize}
\item Si $D \neq 0$ : \textbf{solution unique} (droites sécantes), donnée par Cramer.
\item Si $D = 0$ : les droites sont parallèles (coefficients de $x$ et $y$ proportionnels). Deux sous-cas :
\begin{itemize}
\item les équations sont \emph{entièrement} proportionnelles (seconds membres compris) : droites \textbf{confondues}, \textbf{infinité de solutions} ;
\item sinon : droites \textbf{strictement parallèles}, \textbf{aucune solution}.
\end{itemize}
\end{itemize}
\end{propriete}

\begin{attention}[Piège classique : $D = 0$ ne signifie pas toujours « aucune solution »]
Erreur fréquente au Bac : calculer $D$, trouver $0$, et conclure immédiatement $\mathcal{S} = \emptyset$. C'est faux ! Compare :
\[
\text{(S}_1\text{)} \begin{cases} x + y = 3 \\ 2x + 2y = 6 \end{cases} \qquad \text{(S}_2\text{)} \begin{cases} x + y = 3 \\ 2x + 2y = 7 \end{cases}
\]
Dans les deux cas $D = 1\times 2 - 1\times 2 = 0$. Mais dans (S$_1$), la seconde équation est le double de la première : les droites sont confondues et $\mathcal{S} = \{(x\,;y) \in \mathbb{R}^2 \mid x + y = 3\}$ (infinité de solutions). Dans (S$_2$), les équations sont contradictoires : $\mathcal{S} = \emptyset$. Quand $D = 0$, il faut \textbf{toujours} examiner si les équations sont proportionnelles, seconds membres compris.
\end{attention}

\courssub{Exemple chiffré : un même système résolu par les trois méthodes}

\begin{exoresolu}[Résolution du système $\begin{cases} 2x + 3y = 12 \\ x - y = 1 \end{cases}$ par les trois méthodes]
Résoudre dans $\mathbb{R}^2$ le système $\begin{cases} 2x + 3y = 12 \quad (E_1)\\ x - y = 1 \quad (E_2) \end{cases}$ par substitution, par combinaison, puis par le déterminant.
\tcblower
\textbf{1. Par substitution.} Dans $(E_2)$, le coefficient de $x$ vaut $1$ : on isole $x$ :
\[ x = 1 + y. \]
On substitue dans $(E_1)$ : $2(1 + y) + 3y = 12$, soit $2 + 5y = 12$, d'où $5y = 10$ et $y = 2$. Alors $x = 1 + 2 = 3$.

\textbf{2. Par combinaison.} Pour éliminer $x$, on calcule $(E_1) - 2\times(E_2)$ :
\[ (2x + 3y) - 2(x - y) = 12 - 2 \times 1 \;\Longrightarrow\; 5y = 10 \;\Longrightarrow\; y = 2. \]
Pour éliminer $y$, on calcule $(E_1) + 3\times(E_2)$ :
\[ (2x + 3y) + 3(x - y) = 12 + 3 \;\Longrightarrow\; 5x = 15 \;\Longrightarrow\; x = 3. \]

\textbf{3. Par le déterminant (Cramer).}
\[ D = \begin{vmatrix} 2 & 3 \\ 1 & -1 \end{vmatrix} = 2\times(-1) - 3\times 1 = -5 \neq 0 \quad \Rightarrow \quad \text{solution unique.} \]
\[ D_x = \begin{vmatrix} 12 & 3 \\ 1 & -1 \end{vmatrix} = 12\times(-1) - 3\times 1 = -15, \qquad D_y = \begin{vmatrix} 2 & 12 \\ 1 & 1 \end{vmatrix} = 2 - 12 = -10. \]
\[ x = \frac{D_x}{D} = \frac{-15}{-5} = 3, \qquad y = \frac{D_y}{D} = \frac{-10}{-5} = 2. \]

\textbf{Vérification.} Pour $(3\,;2)$ : $2\times 3 + 3\times 2 = 6 + 6 = 12$ ✓ et $3 - 2 = 1$ ✓.
\[ \boxed{\mathcal{S} = \{(3\,;2)\}} \]
Les trois méthodes donnent le même résultat : rassurant, et excellent moyen de contrôle !
\end{exoresolu}

\begin{methode}[Choisir la méthode la plus rapide]
\begin{enumerate}
\item Une inconnue a un coefficient $1$ ou $-1$ (comme $x - y = 1$) ? \textbf{Substitution} : c'est le plus rapide.
\item Les coefficients d'une inconnue sont multiples ou opposés (ex. $2x$ et $-4x$) ? \textbf{Combinaison} : une seule addition suffit.
\item Système « quelconque » sans coefficient commode, ou question demandant de \emph{discuter} ? \textbf{Déterminant} : le calcul de $D$ tranche immédiatement sur l'unicité.
\item Dans tous les cas, termine par la \textbf{vérification} du couple dans les deux équations initiales : c'est gratuit et cela élimine les erreurs de calcul.
\end{enumerate}
\end{methode}

\begin{exoresolu}[Un système sans solution et un système à infinité de solutions]
Résoudre : (a) $\begin{cases} 2x - 4y = 5 \\ -x + 2y = 1 \end{cases}$ \qquad (b) $\begin{cases} 3x - 6y = 9 \\ x - 2y = 3 \end{cases}$
\tcblower
\textbf{(a)} $D = 2\times 2 - (-4)\times(-1) = 4 - 4 = 0$. Les coefficients des inconnues sont proportionnels ($-x + 2y = -\tfrac{1}{2}(2x - 4y)$), mais $-\tfrac{1}{2}\times 5 = -2{,}5 \neq 1$ : les seconds membres ne suivent pas la proportionnalité. Par combinaison, $(E_1) + 2(E_2)$ donne $0 = 7$, impossible. Les droites sont strictement parallèles : $\mathcal{S} = \emptyset$.

\textbf{(b)} $D = 3\times(-2) - (-6)\times 1 = -6 + 6 = 0$. Ici la première équation vaut exactement $3$ fois la seconde ($3x - 6y = 9 \Leftrightarrow x - 2y = 3$) : les droites sont confondues. Tout couple vérifiant $x = 3 + 2y$ convient :
\[ \mathcal{S} = \{(3 + 2y\,;\,y),\ y \in \mathbb{R}\}, \quad \text{une infinité de solutions.} \]
\end{exoresolu}

\begin{exercice}[Pour t'entraîner]
\begin{enumerate}
\item Résoudre par la méthode de ton choix : $\begin{cases} 3x + 2y = 7 \\ 5x - y = 16 \end{cases}$ puis $\begin{cases} 4x - 2y = 6 \\ -6x + 3y = -9 \end{cases}$.
\item Chez un opérateur téléphonique à Douala, un forfait A coûte 500 FCFA de fixe plus 50 FCFA par minute, un forfait B coûte 1 100 FCFA de fixe plus 30 FCFA par minute. Pour quelle durée de communication les deux forfaits coûtent-ils le même prix ? (Poser un système.)
\end{enumerate}
\end{exercice}

\begin{corrige}[Exercice 1]
\textbf{1.} Premier système : par combinaison, $(E_1) + 2(E_2)$ donne $13x = 39$, donc $x = 3$ ; alors $y = 5x - 16 = -1$. Vérification : $9 - 2 = 7$ ✓, $15 + 1 = 16$ ✓. $\mathcal{S} = \{(3\,;-1)\}$.

Second système : $D = 4\times 3 - (-2)\times(-6) = 12 - 12 = 0$. Or $(E_1) \times (-\tfrac{3}{2})$ redonne exactement $(E_2)$ : équations proportionnelles, droites confondues. $\mathcal{S} = \{(x\,;y) \mid y = 2x - 3\}$, infinité de solutions.

\textbf{2.} Soit $t$ la durée en minutes et $C$ le coût : $\begin{cases} C = 500 + 50t \\ C = 1100 + 30t \end{cases}$. Par substitution : $500 + 50t = 1100 + 30t$, soit $20t = 600$, $t = 30$ min, et $C = 2\,000$ FCFA. Les deux forfaits coûtent le même prix pour 30 minutes de communication.
\end{corrige}

\begin{savaistu}[Gabriel Cramer et l'histoire des systèmes]
Les \cle{formules de Cramer} portent le nom du mathématicien suisse Gabriel Cramer, qui les publia en 1750 dans son \emph{Introduction à l'analyse des lignes courbes algébriques}. Mais l'idée de résoudre des systèmes par élimination est bien plus ancienne : le traité chinois \emph{Les Neuf Chapitres sur l'art du calcul} (vers 200 av. J.-C.) décrivait déjà une méthode équivalente au pivot de Gauss, que tu découvriras dans la prochaine section pour les systèmes $3\times 3$. Aujourd'hui, ces méthodes tournent dans les ordinateurs qui gèrent les réseaux MTN et Orange, les prévisions météo ou les bilans comptables des entreprises : ce sont des systèmes à des millions d'inconnues que l'on résout par les mêmes principes d'élimination !
\end{savaistu}

Tu maîtrises désormais les trois méthodes de résolution des systèmes $2\times 2$ et la discussion complète selon le déterminant. La section suivante généralise ces idées aux systèmes de trois équations à trois inconnues, où la méthode du \cle{pivot de Gauss} deviendra ton outil le plus puissant.

\coursec{Systèmes linéaires 3×3 et pivot de Gauss}

Dans la section précédente, tu as appris à résoudre des systèmes de deux équations à deux inconnues, par substitution, par combinaison ou à l'aide du déterminant. Mais la vie réelle — et le Baccalauréat ! — propose souvent des problèmes à \emph{trois} inconnues : trois prix à déterminer au marché de Mokolo, trois débits dans un réseau d'eau de la CAMWATER, trois investissements d'un commerçant. Il faut alors résoudre un système de trois équations à trois inconnues. Bonne nouvelle : les idées que tu maîtrises déjà (substitution, combinaison) fonctionnent encore, et nous allons y ajouter une méthode systématique et infaillible, le \cle{pivot de Gauss}, qui est l'outil utilisé par tous les ordinateurs du monde.

\courssub{Systèmes de trois équations linéaires à trois inconnues}

\begin{definition}[Système linéaire 3×3]
On appelle \cle{système linéaire de trois équations à trois inconnues} tout système de la forme
\[
\begin{cases}
a_1 x + b_1 y + c_1 z = d_1 \\
a_2 x + b_2 y + c_2 z = d_2 \\
a_3 x + b_3 y + c_3 z = d_3
\end{cases}
\]
où les $a_i$, $b_i$, $c_i$, $d_i$ sont des réels donnés (avec, pour chaque ligne, au moins un des coefficients $a_i, b_i, c_i$ non nul) et où $x$, $y$, $z$ sont les \cle{inconnues}. \emph{Résoudre} le système, c'est trouver tous les triplets $(x\,;y\,;z)$ de $\mathbb{R}^3$ qui vérifient \emph{simultanément} les trois équations.
\end{definition}

\textbf{Interprétation géométrique (admise).} Dans l'espace muni d'un repère $(O\,;\vec{i},\vec{j},\vec{k})$, une équation du type $ax + by + cz = d$ (avec $(a,b,c) \neq (0,0,0)$) représente un \cle{plan}. Résoudre un système 3×3 revient donc à chercher l'intersection de \emph{trois plans}. Trois situations sont possibles :
\begin{itemize}
\item les trois plans se coupent en \textbf{un point unique} : le système a une solution unique ;
\item les trois plans n'ont \textbf{aucun point commun} (par exemple deux plans strictement parallèles) : le système est \emph{incompatible}, il n'a aucune solution ;
\item les trois plans ont \textbf{une droite ou un plan entier en commun} : le système est \emph{indéterminé}, il a une infinité de solutions.
\end{itemize}

\begin{popfigure}
\begin{tikzpicture}[scale=1.05]
% Plan 1 (bleu)
\fill[popblueL,opacity=0.85] (-2.6,-0.4) -- (0.4,-1.6) -- (2.8,-0.2) -- (-0.2,1.0) -- cycle;
\draw[popblue] (-2.6,-0.4) -- (0.4,-1.6) -- (2.8,-0.2) -- (-0.2,1.0) -- cycle;
\node[popblue] at (-1.9,0.6) {$P_1$};
% Plan 2 (orange)
\fill[poporangeL,opacity=0.8] (-1.8,-1.4) -- (1.4,-1.4) -- (2.6,1.2) -- (-0.6,1.2) -- cycle;
\draw[poporange] (-1.8,-1.4) -- (1.4,-1.4) -- (2.6,1.2) -- (-0.6,1.2) -- cycle;
\node[poporange] at (2.1,0.9) {$P_2$};
% Plan 3 (vert)
\fill[popgreenL,opacity=0.8] (-1.2,-1.8) -- (0.6,1.8) -- (2.0,1.0) -- (0.2,-2.6) -- cycle;
\draw[popgreen] (-1.2,-1.8) -- (0.6,1.8) -- (2.0,1.0) -- (0.2,-2.6) -- cycle;
\node[popgreen] at (0.9,1.7) {$P_3$};
% Point d'intersection
\fill[popink] (0.1,-0.35) circle (2.2pt);
\node[popink,below right=1pt] at (0.1,-0.35) {$(x\,;y\,;z)$};
\end{tikzpicture}
\legende{Trois plans de l'espace se coupant en un point unique : le système 3×3 admet une solution unique.}
\end{popfigure}

\begin{attention}[Trois équations, toujours trois équations]
Contrairement au cas 2×2, les calculs sont plus longs et le risque d'erreur augmente. Deux pièges classiques : \textbf{combiner une ligne avec elle-même} (par exemple $L_2 \leftarrow L_2 - L_2$, qui donne $0 = 0$ et détruit l'information) et \textbf{perdre une équation en route}. À chaque étape, ton système doit toujours comporter exactement \emph{trois} équations. À la fin, \textbf{vérifie toujours} ta solution en la reportant dans les trois équations initiales.
\end{attention}

\courssub{Méthode de substitution}

L'idée est identique à celle du cas 2×2 : on \emph{exprime} une inconnue en fonction des deux autres grâce à une équation, puis on \emph{reporte} cette expression dans les deux autres équations. On obtient ainsi un système 2×2, que l'on sait résoudre.

\begin{exemplebox}[Substitution]
Résolvons $\begin{cases} x + y + z = 6 \\ 2x - y + z = 3 \\ x + 2y - z = 2 \end{cases}$.

La première équation donne $x = 6 - y - z$. On reporte dans les deux autres :
\[
\begin{cases}
2(6 - y - z) - y + z = 3 \\
(6 - y - z) + 2y - z = 2
\end{cases}
\iff
\begin{cases}
-3y - z = -9 \\
y - 2z = -4
\end{cases}
\]
De la deuxième : $y = 2z - 4$, reporté dans la première : $-3(2z-4) - z = -9$, soit $-7z = -21$, donc $z = 3$. Alors $y = 2\times 3 - 4 = 2$ et $x = 6 - 2 - 3 = 1$. Solution : $(1\,;2\,;3)$.
\end{exemplebox}

La substitution est efficace lorsqu'une inconnue a pour coefficient $1$ ou $-1$ (sinon, les fractions alourdissent les calculs).

\courssub{Méthode de combinaisons}

On élimine une même inconnue entre deux \emph{paires} d'équations, ce qui ramène le problème à un système 2×2.

\begin{exemplebox}[Combinaisons]
Reprenons le même système. Éliminons $z$ :
\begin{itemize}
\item entre $L_1$ et $L_3$ : $L_1 + L_3$ donne $2x + 3y = 8$ ;
\item entre $L_2$ et $L_3$ : $L_2 + L_3$ donne $3x + y = 5$.
\end{itemize}
On résout $\begin{cases} 2x + 3y = 8 \\ 3x + y = 5 \end{cases}$ : de la seconde, $y = 5 - 3x$, reporté : $2x + 15 - 9x = 8$, soit $x = 1$, puis $y = 2$, et enfin $z = 6 - 1 - 2 = 3$.
\end{exemplebox}

\courssub{Le pivot de Gauss : une méthode systématique}

Substitution et combinaisons demandent un peu d'astuce. Le \cle{pivot de Gauss} est une méthode \emph{algorithmique} : elle fonctionne toujours de la même façon, sans réfléchir au cas par cas. C'est exactement pour cela qu'elle est programmable sur machine.

\begin{propriete}[Opérations élémentaires — admise]
Les trois opérations suivantes transforment un système en un système \cle{équivalent} (c'est-à-dire ayant exactement les mêmes solutions) :
\begin{enumerate}
\item \textbf{échanger} deux lignes : $L_i \leftrightarrow L_j$ ;
\item \textbf{multiplier} une ligne par un réel non nul : $L_i \leftarrow \lambda L_i$ avec $\lambda \neq 0$ ;
\item \textbf{ajouter} à une ligne un multiple d'une autre ligne : $L_i \leftarrow L_i + \lambda L_j$ avec $j \neq i$.
\end{enumerate}
\end{propriete}

\textbf{Idée de la méthode.} À l'aide de ces opérations, on élimine $x$ des lignes 2 et 3 (grâce au coefficient de $x$ dans la ligne 1, appelé \cle{pivot}), puis on élimine $y$ de la ligne 3. Le système devient \emph{triangulaire} (ou échelonné) : la dernière équation ne contient que $z$, la deuxième contient $y$ et $z$, la première contient tout. On résout alors « en remontant » : $z$ d'abord, puis $y$, puis $x$.

\begin{methode}[Algorithme du pivot de Gauss]
\begin{enumerate}
\item \textbf{Choisir le pivot} de la colonne $x$ : un coefficient non nul (de préférence $1$ ou $-1$). Si besoin, échanger deux lignes pour le placer en $L_1$.
\item \textbf{Éliminer $x$} dans $L_2$ et $L_3$ par des opérations du type $L_2 \leftarrow L_2 - a_2 L_1$ et $L_3 \leftarrow L_3 - a_3 L_1$ (si le pivot vaut 1).
\item \textbf{Éliminer $y$} dans la nouvelle ligne $L_3$ en utilisant le coefficient de $y$ dans $L_2$ comme nouveau pivot.
\item \textbf{Remontée} : résoudre la dernière équation (en $z$ seul), reporter dans la deuxième pour trouver $y$, puis dans la première pour trouver $x$.
\item \textbf{Vérifier} le triplet obtenu dans les trois équations initiales.
\end{enumerate}
\end{methode}

\begin{exoresolu}[Pivot de Gauss, pas à pas]
Résoudre dans $\mathbb{R}^3$ : $\begin{cases} x + y + z = 6 \\ 2x - y + z = 3 \\ x + 2y - z = 2 \end{cases}$
\tcblower
Le coefficient de $x$ dans $L_1$ vaut $1$ : c'est un pivot idéal.

\textbf{Étape 1 : éliminer $x$ dans $L_2$ et $L_3$.}
$L_2 \leftarrow L_2 - 2L_1$ : $(2x - y + z) - 2(x + y + z) = 3 - 12$, soit $-3y - z = -9$.
$L_3 \leftarrow L_3 - L_1$ : $(x + 2y - z) - (x + y + z) = 2 - 6$, soit $y - 2z = -4$.
\[
\begin{cases}
x + y + z = 6 \\
-3y - z = -9 \\
y - 2z = -4
\end{cases}
\]

\textbf{Étape 2 : éliminer $y$ dans $L_3$.}
$L_3 \leftarrow 3L_3 + L_2$ : $3(y - 2z) + (-3y - z) = -12 - 9$, soit $-7z = -21$.
\[
\begin{cases}
x + y + z = 6 \\
-3y - z = -9 \\
-7z = -21
\end{cases}
\]

\textbf{Étape 3 : remontée.} La dernière ligne donne $z = 3$. Reporté dans la deuxième : $-3y - 3 = -9$, donc $y = 2$. Reporté dans la première : $x + 2 + 3 = 6$, donc $x = 1$.

\textbf{Vérification.} $1 + 2 + 3 = 6$ ; $2 - 2 + 3 = 3$ ; $1 + 4 - 3 = 2$. Les trois équations sont satisfaites.

\textbf{Conclusion :} $\mathcal{S} = \{(1\,;2\,;3)\}$.
\end{exoresolu}

\begin{popfigure}
\begin{tikzpicture}[scale=0.92, every node/.style={font=\small}]
\node at (0,2.6) {\textbf{Système initial}};
\node[align=center] at (0,0) {$\begin{pmatrix} 1 & 1 & 1 \\ 2 & -1 & 1 \\ 1 & 2 & -1 \end{pmatrix} \begin{pmatrix} x \\ y \\ z \end{pmatrix} = \begin{pmatrix} 6 \\ 3 \\ 2 \end{pmatrix}$};
\draw[-{Stealth},very thick,poporange] (2.6,0) -- (4.0,0) node[midway,above,popdark]{$L_2 \leftarrow L_2 - 2L_1$} node[midway,below,popdark]{$L_3 \leftarrow L_3 - L_1$};
\node at (6.6,2.6) {\textbf{Après élimination de } $x$};
\node[align=center] at (6.6,0) {$\begin{pmatrix} \mathbf{1} & 1 & 1 \\ 0 & -3 & -1 \\ 0 & 1 & -2 \end{pmatrix} \begin{pmatrix} x \\ y \\ z \end{pmatrix} = \begin{pmatrix} 6 \\ -9 \\ -4 \end{pmatrix}$};
\draw[-{Stealth},very thick,poporange] (6.6,-1.5) -- (6.6,-2.8) node[midway,right,popdark]{$L_3 \leftarrow 3L_3 + L_2$};
\node[align=center] at (6.6,-4.6) {$\begin{pmatrix} \mathbf{1} & 1 & 1 \\ 0 & \mathbf{-3} & -1 \\ 0 & 0 & \mathbf{-7} \end{pmatrix} \begin{pmatrix} x \\ y \\ z \end{pmatrix} = \begin{pmatrix} 6 \\ -9 \\ -21 \end{pmatrix}$};
\node[popgreen] at (6.6,-6.2) {\textbf{Forme triangulaire : remontée } $z = 3,\ y = 2,\ x = 1$};
\end{tikzpicture}
\legende{Réduction triangulaire du système par pivot de Gauss : sous la diagonale (en gras), on ne laisse que des zéros.}
\end{popfigure}

\courssub{Cas particuliers : systèmes incompatible et indéterminé}

Le pivot de Gauss détecte automatiquement les cas sans solution ou avec une infinité de solutions.

\begin{aretenir}[Lecture des cas particuliers]
Après réduction triangulaire :
\begin{itemize}
\item si une ligne s'écrit $0x + 0y + 0z = k$ avec $k \neq 0$ (par exemple $0 = 5$), c'est \textbf{impossible} : le système est \cle{incompatible}, $\mathcal{S} = \emptyset$ ;
\item si une ligne s'écrit $0x + 0y + 0z = 0$ (elle disparaît), le système se ramène à moins d'équations que d'inconnues : il est \cle{indéterminé} et admet une \textbf{infinité de solutions}, que l'on exprime à l'aide d'un paramètre (par exemple $z = t$, $t \in \mathbb{R}$).
\end{itemize}
\end{aretenir}

\begin{exemplebox}[Un système incompatible]
$\begin{cases} x + y + z = 3 \\ 2x + 2y + 2z = 5 \\ x - y + z = 1 \end{cases}$

$L_2 \leftarrow L_2 - 2L_1$ donne $0x + 0y + 0z = -1$, c'est-à-dire $0 = -1$ : absurde. Le système n'a \textbf{aucune solution} (géométriquement, les deux premiers plans sont strictement parallèles).
\end{exemplebox}

\begin{exemplebox}[Un système indéterminé]
$\begin{cases} x + y + z = 3 \\ 2x + 2y + 2z = 6 \\ x - y + z = 1 \end{cases}$

$L_2 \leftarrow L_2 - 2L_1$ donne $0 = 0$ : la ligne disparaît. Il reste $\begin{cases} x + y + z = 3 \\ x - y + z = 1 \end{cases}$. En soustrayant : $2y = 2$, donc $y = 1$, puis $x + z = 2$. En posant $z = t$ ($t \in \mathbb{R}$), on obtient $\mathcal{S} = \{(2 - t\,;\,1\,;\,t),\ t \in \mathbb{R}\}$ : une infinité de solutions (les trois plans se coupent le long d'une droite).
\end{exemplebox}

\begin{attention}[Rédaction des opérations sur les lignes]
Annonce \emph{toujours} l'opération effectuée ($L_2 \leftarrow L_2 - 2L_1$) avant d'écrire le nouveau système : le correcteur du Bac doit pouvoir suivre ton raisonnement, et toi-même tu pourras vérifier tes calculs. Autre piège : dans $L_3 \leftarrow 3L_3 + L_2$, n'oublie pas de multiplier \emph{tous} les termes de $L_3$ par 3, y compris le second membre. Signalons enfin que les \textbf{systèmes à paramètres} (dont les coefficients dépendent d'une lettre $m$ à discuter) sont \emph{hors programme} en Terminale A : tu n'auras à résoudre que des systèmes à coefficients numériques.
\end{attention}

\begin{exercice}[À toi de jouer]
Résoudre par le pivot de Gauss : $\begin{cases} x + 2y + z = 8 \\ 2x + y - z = 1 \\ x + y + 2z = 9 \end{cases}$. Vérifier la solution.
\end{exercice}

\begin{corrige}[Exercice]
$L_2 \leftarrow L_2 - 2L_1$ : $-3y - 3z = -15$, soit $y + z = 5$. $L_3 \leftarrow L_3 - L_1$ : $-y + z = 1$. En additionnant ces deux lignes : $2z = 6$, donc $z = 3$, puis $y = 2$, et $x = 8 - 4 - 3 = 1$. Vérification : $1 + 4 + 3 = 8$ ; $2 + 2 - 3 = 1$ ; $1 + 2 + 6 = 9$. $\mathcal{S} = \{(1\,;2\,;3)\}$.
\end{corrige}

\begin{savaistu}[Gauss, la météo et des millions d'équations]
Carl Friedrich Gauss (1777--1855), surnommé le « prince des mathématiciens », a systématisé cette méthode pour ses calculs d'astronomie. Aujourd'hui, le pivot de Gauss est au cœur de tous les logiciels de calcul scientifique : les prévisions météorologiques, la conception des avions ou la simulation des réseaux électriques (comme celui géré par ENEO) se ramènent à la résolution de systèmes comportant des \emph{millions} d'équations à des millions d'inconnues — résolus par des versions ultra-rapides de la méthode que tu viens d'apprendre !
\end{savaistu}

\coursec{Problèmes concrets et modélisation}

Dans les sections précédentes, tu as appris à résoudre des systèmes linéaires à trois inconnues par substitution, combinaisons ou pivot de Gauss, ainsi que des équations et inéquations de degré 3. Ces techniques forment une véritable boîte à outils. Mais au Baccalauréat, l'énoncé ne te dira jamais « résous le système suivant » sans raison : il te racontera une histoire — un marchand de Mokolo, une cantine scolaire, une cuve d'eau à Ebolowa — et ce sera à \emph{toi} de traduire cette histoire en mathématiques. Cette traduction s'appelle la \cle{modélisation}. C'est la compétence reine de cette leçon, et c'est elle qui fait la différence entre une copie moyenne et une excellente copie.

\courssub{La démarche de modélisation en quatre étapes}

\textbf{Intuition.} Un problème concret est comme un texte écrit dans une langue étrangère : la langue de la vie quotidienne. Les mathématiques sont une autre langue, plus précise. Modéliser, c'est \emph{traduire} de la première langue vers la seconde, résoudre dans la langue mathématique (où l'on est très à l'aise), puis \emph{retraduire} la réponse dans la langue du problème. Le schéma suivant résume ce va-et-vient.

\begin{popfigure}
\begin{tikzpicture}[
  box/.style={draw=popblue, thick, rounded corners=6pt, fill=popblueL, text width=3.1cm, minimum height=1.5cm, align=center, font=\small},
  arr/.style={-{Stealth[length=3mm]}, very thick, poporange},
  lbl/.style={font=\footnotesize\itshape, text=popdark, midway, above}
]
\node[box, align=center] (a) at (0,0) {\textbf{Situation réelle}\\ problème concret};
\node[box, align=center] (b) at (5.2,0) {\textbf{Modèle mathématique}\\ équation, inéquation, système};
\node[box, align=center] (c) at (10.4,0) {\textbf{Solution mathématique}\\ valeurs de $x, y, z$};
\node[box, fill=popgreenL, draw=popgreen, align=center] (d) at (5.2,-3.2) {\textbf{Solution concrète}\\ interprétée et validée};
\draw[arr] (a) -- (b) node[lbl]{traduction};
\draw[arr] (b) -- (c) node[lbl]{résolution};
\draw[arr] (c.south) |- (d.east) node[lbl, above, xshift=-1.2cm, yshift=0.35cm]{interprétation};
\draw[arr, dashed, poppurple] (d.west) -| (a.south) node[lbl, above, xshift=1.6cm, yshift=-0.9cm]{validation};
\end{tikzpicture}
\legende{Le cycle de la modélisation : on traduit, on résout, on interprète, on valide.}
\end{popfigure}

\begin{methode}[Fiche de modélisation : les 4 étapes à rédiger dans une copie]
\begin{enumerate}
\item \textbf{Choix des inconnues.} Écris une phrase complète : « Soit $x$ le prix d'un cahier (en FCFA), $y$ le prix d'un stylo (en FCFA) ». Précise toujours l'unité et, si besoin, la nature de l'inconnue (entier, réel positif).
\item \textbf{Mise en équation(s).} Traduis chaque information de l'énoncé en une égalité ou une inégalité. Une information = une équation. Il faut autant d'équations indépendantes que d'inconnues.
\item \textbf{Résolution.} Résous le modèle avec les techniques du chapitre (factorisation, pivot de Gauss, tableau de signes...). Cette étape est purement mathématique.
\item \textbf{Interprétation et validation.} Vérifie que les solutions sont \emph{compatibles avec le contexte} (positives, entières, bornées), vérifie-les dans l'énoncé, puis conclus par une phrase en langage courant, sans symbole mathématique.
\end{enumerate}
\end{methode}

\begin{attention}[Solution exacte mais impossible dans le contexte]
Une solution mathématiquement correcte peut être \emph{concrètement absurde}. Exemples typiques :
\begin{itemize}
\item un prix de $-250$ FCFA : à rejeter, un prix est positif ;
\item $3{,}7$ élèves dans un bus : à rejeter, on compte des personnes entières ;
\item un terrain de longueur $x = -2$ m : à rejeter, une longueur est positive ;
\item une racine $x = 12$ alors que l'énoncé impose $x \le 10$ : à rejeter, hors contrainte.
\end{itemize}
Le correcteur attend que tu \textbf{écrives explicitement le rejet} : « $x = -2$ est impossible car une longueur est positive ». Oublier cette phrase, c'est perdre des points même avec un calcul juste.
\end{attention}

\courssub{Problèmes de prix et de ventes : vers un système linéaire}

\textbf{Intuition.} Quand un commerçant vend des lots mélangés (par exemple « 2 pains et 3 beignets pour 500 FCFA »), chaque achat fournit une équation. Deux achats différents donnent un système $2\times 2$ ; trois achats donnent un système $3\times 3$. Les prix unitaires sont les inconnues.

\begin{exemplebox}[Problème de la cantine scolaire (système $3\times 3$)]
La cantine d'un lycée de Bafoussam propose trois menus : menu A (riz + sauce), menu B (ndolé + miondo), menu C (poulet + frites). Le comptable a noté les recettes de trois jours :
\begin{itemize}
\item lundi : 40 menus A, 30 menus B, 20 menus C pour \num{115000} FCFA ;
\item mardi : 30 menus A, 50 menus B, 10 menus C pour \num{120000} FCFA ;
\item mercredi : 20 menus A, 20 menus B, 40 menus C pour \num{110000} FCFA.
\end{itemize}
\emph{Étape 1 — Inconnues.} Soit $x$, $y$, $z$ les prix unitaires (en milliers de FCFA) des menus A, B, C.

\emph{Étape 2 — Mise en équations.} Les recettes en milliers de FCFA donnent :
\[
\begin{cases}
40x + 30y + 20z = 115 \\
30x + 50y + 10z = 120 \\
20x + 20y + 40z = 110
\end{cases}
\]
On simplifie en divisant chaque équation par 10 :
\[
\begin{cases}
4x + 3y + 2z = 11{,}5 \quad (L_1)\\
3x + 5y + z = 12 \quad (L_2)\\
2x + 2y + 4z = 11 \quad (L_3)
\end{cases}
\]

\emph{Étape 3 — Résolution par pivot de Gauss.} On élimine $x$ :
$L_2 \leftarrow 4L_2 - 3L_1$ : $12x + 20y + 4z - (12x + 9y + 6z) = 48 - 34{,}5 \;\Rightarrow\; 11y - 2z = 13{,}5 \quad (L_2')$
$L_3 \leftarrow 2L_3 - L_1$ : $4x + 4y + 8z - (4x + 3y + 2z) = 22 - 11{,}5 \;\Rightarrow\; y + 6z = 10{,}5 \quad (L_3')$

On a maintenant un système $2\times 2$ en $y, z$ :
\[
\begin{cases}
11y - 2z = 13{,}5 \\
y + 6z = 10{,}5
\end{cases}
\]
De $L_3'$ : $y = 10{,}5 - 6z$. En injectant dans $L_2'$ :
$11(10{,}5 - 6z) - 2z = 13{,}5 \;\Rightarrow\; 115{,}5 - 68z = 13{,}5 \;\Rightarrow\; 68z = 102 \;\Rightarrow\; z = 1{,}5$.
Alors $y = 10{,}5 - 6(1{,}5) = 10{,}5 - 9 = 1{,}5$.
En remontant à $L_1$ : $4x = 11{,}5 - 3(1{,}5) - 2(1{,}5) = 11{,}5 - 4{,}5 - 3 = 4 \;\Rightarrow\; x = 1$.

\emph{Étape 4 — Validation et conclusion.} Les trois prix sont positifs : cohérent. Vérification lundi : $40\times1000 + 30\times1500 + 20\times1500 = \num{115000}$ FCFA. \textbf{Le menu A coûte \num{1000} FCFA, les menus B et C coûtent \num{1500} FCFA chacun.}
\end{exemplebox}

\begin{attention}[Bien choisir l'unité avant de poser le système]
Dans l'exemple précédent, travailler en milliers de FCFA a simplifié tous les coefficients. Mais alors la réponse $x = 1$ signifie \num{1000} FCFA, pas 1 FCFA ! Une erreur d'unité à l'étape 4 ruine une résolution parfaite. Annonce ton unité à l'étape 1 et respecte-la jusqu'au bout.
\end{attention}

\courssub{Problèmes de géométrie : vers une équation polynomiale (degré 2 ou 3)}

\textbf{Intuition.} Dès qu'un volume intervient, on multiplie trois longueurs : le volume devient naturellement un polynôme de degré 3 en la dimension inconnue. Les problèmes d'aire conduisent, eux, à des équations de degré 2. Résoudre le problème revient alors à chercher une racine évidente et à factoriser (degré 3) ou à utiliser le trinôme (degré 2).

\begin{exemplebox}[La cuve cubique (équation de degré 3)]
Une entreprise de Bafang construit une cuve cubique de côté $x$ (en mètres). Le cahier des charges impose : le volume brut $x^3$, diminué de $3x^2$ (épaisseur des parois), diminué de $4x$ (canalisation), augmenté de $12$ (fond renforcé), doit être exactement nul au point d'équilibre de la conception. On cherche donc $x > 0$ tel que
\[ x^3 - 3x^2 - 4x + 12 = 0. \]
\emph{Racine évidente.} On teste les diviseurs de $12$ : pour $x = 2$ : $8 - 12 - 8 + 12 = 0$. Donc $2$ est racine.

\emph{Factorisation} par identification : $x^3 - 3x^2 - 4x + 12 = (x - 2)(x^2 + bx - 6)$. En développant le terme en $x^2$ : $b - 2 = -3$, donc $b = -1$. Ainsi
\[ x^3 - 3x^2 - 4x + 12 = (x - 2)(x^2 - x - 6) = (x - 2)(x - 3)(x + 2). \]
\emph{Solutions mathématiques :} $x \in \{-2 ; 2 ; 3\}$.

\emph{Validation.} $x = -2$ est impossible (longueur négative). Il reste deux conceptions possibles : une cuve de côté $2$ m ou de côté $3$ m. Le choix final dépendra de critères extérieurs (coût, encombrement).
\end{exemplebox}

\begin{exoresolu}[Dimensions d'un terrain rectangulaire (équation de degré 2)]
Un terrain rectangulaire a une longueur qui dépasse sa largeur de $5$ m. Son aire est de $300$ m$^2$. Déterminer ses dimensions.
\tcblower
\emph{Étape 1.} Soit $x$ la largeur en mètres ($x > 0$) ; la longueur est $x + 5$.

\emph{Étape 2.} Aire : $x(x + 5) = 300$, soit $x^2 + 5x - 300 = 0$.

\emph{Étape 3.} $\Delta = 25 + 1200 = 1225 = 35^2$. D'où $x = \dfrac{-5 + 35}{2} = 15$ ou $x = \dfrac{-5 - 35}{2} = -20$.

\emph{Étape 4.} $x = -20$ est rejeté (longueur négative). Vérification : $15 \times 20 = 300$. \textbf{Le terrain mesure 15 m de large et 20 m de long.}
\end{exoresolu}

\courssub{Problèmes de seuil et comparaison d'offres : vers une inéquation}

\textbf{Intuition.} « À partir de combien de minutes l'offre B devient-elle plus avantageuse ? » : le mot-clé \emph{à partir de} signale une inéquation. On exprime le coût de chaque offre en fonction de la même variable, puis on compare.

\begin{exemplebox}[Comparaison de forfaits téléphoniques]
Un opérateur propose : offre A, abonnement de \num{2000} FCFA puis $50$ FCFA par minute ; offre B, sans abonnement, $100$ FCFA par minute. Pour $x$ minutes, les coûts sont $C_A(x) = 2000 + 50x$ et $C_B(x) = 100x$. L'offre A est plus avantageuse lorsque
\[ 2000 + 50x < 100x \iff 2000 < 50x \iff x > 40. \]
\emph{Validation.} $x$ est un temps, donc positif : la condition $x > 40$ est compatible. \textbf{Conclusion : l'offre A devient rentable dès que la consommation dépasse 40 minutes par mois ; en dessous, l'offre B est préférable.} À exactement 40 minutes, les deux offres coûtent le même prix (\num{4000} FCFA).
\end{exemplebox}

\begin{exoresolu}[Seuil de rentabilité d'un vendeur d'arachides]
Une vendeuse du marché Central de Yaoundé achète ses arachides à $300$ FCFA le sachet et paie \num{1500} FCFA de transport par jour. Elle revend chaque sachet $500$ FCFA. Combien de sachets doit-elle vendre par jour pour réaliser un bénéfice strictement positif ?
\tcblower
\emph{Étape 1.} Soit $n$ le nombre de sachets vendus ($n \in \mathbb{N}$).

\emph{Étape 2.} Bénéfice : $B(n) = 500n - 300n - 1500 = 200n - 1500$. On veut $B(n) > 0$.

\emph{Étape 3.} $200n - 1500 > 0 \iff n > 7{,}5$.

\emph{Étape 4.} Comme $n$ est un \emph{entier}, la plus petite solution est $n = 8$. Vérification : $B(8) = 1600 - 1500 = 100 > 0$, mais $B(7) = -100 < 0$. \textbf{La vendeuse doit vendre au moins 8 sachets par jour pour être bénéficiaire.}
\end{exoresolu}

\begin{attention}[Solutions entières : ne pas confondre $n > 7{,}5$ et $n \ge 7$]
Quand l'inconnue est un nombre d'objets, la solution d'une inéquation doit être arrondie \emph{dans le bon sens} : pour $n > 7{,}5$, la réponse est $n \ge 8$ (et non $n \ge 7$). Pour $n \le 7{,}5$, la réponse est $n \le 7$. Prends toujours le temps de tester la valeur frontière dans l'énoncé.
\end{attention}

\courssub{Contraintes du contexte et rédaction de la conclusion}

La dernière étape distingue les bonnes copies. Trois types de contraintes reviennent constamment :
\begin{itemize}
\item \textbf{positivité} : prix, longueurs, volumes, durées sont positifs ou nuls ;
\item \textbf{intégralité} : nombres de personnes, d'objets, de sachets sont des entiers naturels ;
\item \textbf{bornes} : capacité d'un bus, budget maximal, dimensions imposées.
\end{itemize}
La conclusion finale doit être une \emph{phrase complète en français}, reprenant les unités et la situation, sans la moindre lettre $x$ : non pas « $x = 15$ », mais « le terrain mesure 15 mètres de large ».

\begin{savaistu}[Les points faciles du Baccalauréat A]
Au Baccalauréat A, les problèmes de synthèse rapportent souvent des points faciles si la démarche de modélisation est bien rédigée : les barèmes attribuent fréquemment un point au choix des inconnues, un point à la mise en équations, et un point à la conclusion en langage courant — indépendamment de la résolution elle-même. Même si tu bloques sur les calculs, rédige proprement les étapes 1, 2 et 4 : tu peux récolter la moitié des points du problème sans avoir terminé la résolution !
\end{savaistu}

\begin{exercice}[Pour t'entraîner]
Un transporteur de la ligne Douala--Bafoussam propose deux tarifs : tarif 1, \num{3000} FCFA par voyage ; tarif 2, carte mensuelle à \num{10000} FCFA puis \num{2000} FCFA par voyage. À partir de combien de voyages par mois le tarif 2 devient-il avantageux ? Rédige les quatre étapes complètes de la modélisation.
\end{exercice}

\newpage

\coursec{Activité d'intégration}

\coursec{Activité d'intégration : Le défi du marché Mokolo}

\courssub{Objectifs de l'activité}

À l'issue de cette activité d'intégration, tu seras capable de :

\begin{itemize}
\item \cle{modéliser} une situation commerciale réelle par un système linéaire de trois équations à trois inconnues ;
\item \cle{résoudre} ce système par la méthode du \cle{pivot de Gauss}, en rédigeant proprement chaque étape ;
\item \cle{interpréter} une fonction polynôme de degré 3 : trouver une \cle{racine évidente}, \cle{factoriser} par identification des coefficients, puis dresser un \cle{tableau de signes} ;
\item déterminer l'ensemble des solutions d'une \cle{inéquation de degré 3} et en donner une \cle{interprétation concrète} ;
\item rédiger un mini-rapport argumenté et le présenter oralement, en justifiant chaque choix mathématique.
\end{itemize}

Cette activité mobilise donc l'ensemble des savoir-faire de la leçon dans une \emph{situation authentique} de la vie économique camerounaise.

\courssub{Situation}

\textbf{Le défi du marché Mokolo.}

Le marché Mokolo, au cœur de Yaoundé, est l'un des plus grands marchés du Cameroun. Ta classe a été contactée par l'association des commerçants du marché : trois jeunes gestionnaires d'un étal de produits vivriers viennent de reprendre l'affaire familiale et ont \emph{égaré le cahier des prix unitaires}. Ils vendent trois produits :

\begin{itemize}
\item le \textbf{manioc} (bâtons), à $x$ FCFA le kilogramme ;
\item les \textbf{arachides} (grillées), à $y$ FCFA le kilogramme ;
\item l'\textbf{huile de palme}, à $z$ FCFA le litre.
\end{itemize}

Heureusement, ils ont conservé trois \cle{tickets de vente} complets de la journée précédente, sur lesquels figurent les quantités vendues et les montants totaux encaissés :

\begin{center}
\begin{tabularx}{\linewidth}{|l|X|X|X|X|}
\hline
\rowcolor{popblueL} \textbf{Ticket} & \textbf{Manioc (kg)} & \textbf{Arachides (kg)} & \textbf{Huile (L)} & \textbf{Montant (FCFA)} \\
\hline
Ticket n\textsuperscript{o}1 & $2$ & $1$ & $1$ & $\num{2400}$ \\
\hline
Ticket n\textsuperscript{o}2 & $1$ & $3$ & $2$ & $\num{4250}$ \\
\hline
Ticket n\textsuperscript{o}3 & $4$ & $2$ & $3$ & $\num{6000}$ \\
\hline
\end{tabularx}
\end{center}

\textbf{Première mission :} retrouver les trois prix unitaires $x$, $y$ et $z$.

\textbf{Deuxième mission :} l'association a fait appel à un étudiant en économie de l'Université de Yaoundé I, qui a modélisé le bénéfice hebdomadaire de l'étal (en \emph{milliers de FCFA}) en fonction du prix de vente $x$ du kilogramme de manioc (en \emph{dizaines de FCFA}, avec $x \geq 0$) par :
\[ B(x) = -x^3 + 6x^2 - 9x + 4. \]
Les gestionnaires veulent savoir \textbf{pour quelles valeurs du prix $x$ le bénéfice hebdomadaire est strictement positif}, afin de fixer un prix qui garantisse la rentabilité de l'étal sans faire fuir la clientèle.

\courssub{Consignes}

Par groupes de 4, vous êtes les gestionnaires de l'étal. Vous rédigerez un \textbf{mini-rapport} (une page maximum) répondant aux questions suivantes, puis vous présenterez vos conclusions à la classe en 5 minutes.

\begin{enumerate}
\item \textbf{Modélisation.} Traduire chaque ticket de vente par une équation d'inconnues $x$, $y$, $z$. Écrire le système $(S)$ de trois équations à trois inconnues obtenu. Justifier pourquoi ce système est \emph{linéaire}.
\item \textbf{Résolution du système.} Résoudre $(S)$ par la méthode du \cle{pivot de Gauss} : éliminer $x$ dans les deux dernières équations, puis $y$ dans la dernière, et remonter le système triangulaire. Vérifier que la solution trouvée satisfait les trois tickets.
\item \textbf{Étude du bénéfice : racine évidente.} Vérifier que $x = 1$ est une racine du polynôme $B(x) = -x^3 + 6x^2 - 9x + 4$.
\item \textbf{Factorisation.} En déduire une factorisation de $B(x)$ sous la forme $B(x) = (x - 1)(ax^2 + bx + c)$ en déterminant $a$, $b$, $c$ par \cle{identification des coefficients}. Factoriser ensuite complètement le trinôme obtenu.
\item \textbf{Tableau de signes.} Dresser le tableau de signes de $B(x)$ sur $[0\,;\,+\infty[$.
\item \textbf{Conclusion commerciale.} Résoudre l'inéquation $B(x) > 0$ et interpréter le résultat : quelle fourchette de prix du manioc (en FCFA, sachant que $x$ est en dizaines de FCFA) conseillez-vous aux gestionnaires ? Que se passe-t-il aux valeurs $x = 1$ et $x = 4$ ?
\item \textbf{Esprit critique.} Le modèle prédit un bénéfice négatif pour un prix très élevé. Proposer une explication concrète de ce phénomène dans le contexte du marché Mokolo.
\end{enumerate}

\courssub{Ressources à mobiliser}

\begin{aretenir}[Boîte à outils de la leçon]
\begin{itemize}
\item \textbf{Système linéaire $2\times 2$ :} résolution par substitution, combinaison linéaire ou déterminant (règle de Cramer).
\item \textbf{Système linéaire $3\times 3$ — pivot de Gauss :} par des opérations élémentaires $L_i \leftarrow L_i + \lambda L_j$ (ou $L_i \leftarrow aL_i - bL_j$), on transforme le système en système \emph{triangulaire} équivalent, que l'on résout par substitutions successives en remontant.
\item \textbf{Équations/Inéquations par inconnue auxiliaire :} poser $X = \varphi(x)$ pour ramener l'équation à un polynôme de degré 2 ou 1 en $X$, résoudre en $X$, puis revenir à $x$ en vérifiant les conditions de définition.
\item \textbf{Racine évidente :} $\alpha$ est racine de $P$ si et seulement si $P(\alpha) = 0$, ce qui équivaut à l'existence d'un polynôme $Q$ tel que $P(x) = (x-\alpha)Q(x)$.
\item \textbf{Identification des coefficients :} deux polynômes sont égaux si et seulement si leurs coefficients de même degré sont égaux. On développe $(x-\alpha)(ax^2+bx+c)$ et on identifie.
\item \textbf{Tableau de signes :} un produit de facteurs est du signe du produit des signes ; un carré est toujours positif ou nul ; un binôme $x - \alpha$ s'annule en $\alpha$ et est du signe de $x$ après $\alpha$.
\item \textbf{Retour au problème :} toute solution mathématique doit être \emph{interprétée} dans le contexte et \emph{vérifiée} sur les données de départ.
\end{itemize}
\end{aretenir}

\courssub{Démarche et corrigé commenté}

\begin{exoresolu}[Corrigé commenté du défi du marché Mokolo]
\textbf{Énoncé rappelé.} Retrouver les prix unitaires $x$, $y$, $z$ à partir des trois tickets, puis déterminer les valeurs de $x$ pour lesquelles $B(x) = -x^3 + 6x^2 - 9x + 4 > 0$.

\tcblower

\textbf{Question 1 — Modélisation.}

Chaque ticket traduit : (quantité de manioc)$\times x$ + (quantité d'arachides)$\times y$ + (quantité d'huile)$\times z$ = montant. D'où le système :
\[ (S) : \begin{cases} 2x + y + z = \num{2400} & (L_1) \\ x + 3y + 2z = \num{4250} & (L_2) \\ 4x + 2y + 3z = \num{6000} & (L_3) \end{cases} \]
Chaque équation est de la forme $ax + by + cz = d$ : les inconnues apparaissent au degré 1, sans produit entre elles. Le système est donc bien \cle{linéaire}.

\textbf{Question 2 — Pivot de Gauss.}

\emph{Étape 1 : éliminer $x$ dans $(L_2)$ et $(L_3)$ en utilisant le pivot $2x$ de $(L_1)$.}
\begin{align*}
(L_2') &\leftarrow 2(L_2) - (L_1) : \quad 2x + 6y + 4z - (2x + y + z) = \num{8500} - \num{2400} \\
       &\quad\Longrightarrow\quad 5y + 3z = \num{6100} \\[4pt]
(L_3') &\leftarrow (L_3) - 2(L_1) : \quad 4x + 2y + 3z - (4x + 2y + 2z) = \num{6000} - \num{4800} \\
       &\quad\Longrightarrow\quad z = \num{1200}
\end{align*}
Le système est déjà triangulaire :
\[ \begin{cases} 2x + y + z = \num{2400} \\ 5y + 3z = \num{6100} \\ z = \num{1200} \end{cases} \]

\emph{Étape 2 : remontée.} De la dernière ligne : $z = \num{1200}$. On reporte dans la deuxième :
\[ 5y = \num{6100} - 3\times\num{1200} = \num{2500} \quad\Longrightarrow\quad y = 500. \]
Puis dans la première : $2x = \num{2400} - 500 - \num{1200} = 700$, donc $x = 350$.

\emph{Vérification (indispensable !) :}
\begin{itemize}
\item Ticket n\textsuperscript{o}2 : $350 + 3(500) + 2(\num{1200}) = 350 + \num{1500} + \num{2400} = \num{4250}$ \checkmark
\item Ticket n\textsuperscript{o}3 : $4(350) + 2(500) + 3(\num{1200}) = \num{1400} + \num{1000} + \num{3600} = \num{6000}$ \checkmark
\end{itemize}

\textbf{Conclusion :} le manioc coûte \textbf{350 FCFA/kg}, les arachides \textbf{500 FCFA/kg} et l'huile \textbf{\num{1200} FCFA/L}.

\textbf{Question 3 — Racine évidente.}

$B(1) = -1 + 6 - 9 + 4 = 0$. Donc $x = 1$ est racine de $B$, et $B(x)$ est factorisable par $(x - 1)$.

\textbf{Question 4 — Factorisation par identification.}

On cherche $B(x) = (x-1)(ax^2 + bx + c)$. En développant :
\[ (x-1)(ax^2+bx+c) = ax^3 + (b-a)x^2 + (c-b)x - c. \]
Par identification avec $-x^3 + 6x^2 - 9x + 4$ :
\[ a = -1,\quad b - a = 6 \Rightarrow b = 5,\quad c - b = -9 \Rightarrow c = -4,\quad -c = 4 \ \checkmark \]
Donc $B(x) = (x-1)(-x^2 + 5x - 4)$. Le trinôme $-x^2 + 5x - 4$ admet $x = 1$ comme racine évidente ($-1+5-4=0$) ; le produit des racines vaut $\dfrac{-4}{-1} = 4$, donc l'autre racine est $4$. Ainsi :
\[ B(x) = -(x-1)^2(x-4). \]

\textbf{Question 5 — Tableau de signes.}

$(x-1)^2 \geq 0$ pour tout $x$ (nul en $1$) ; $(x-4)$ est négatif avant $4$, positif après ; le signe moins devant inverse le tout.

\begin{center}
\begin{tabular}{|c|c|c|c|c|c|c|}
\hline
$x$ & $0$ & $]0,1[$ & $1$ & $]1,4[$ & $4$ & $]4,+\infty[$ \\
\hline
$-(x-1)^2$ & $-$ & $-$ & $0$ & $-$ & $-$ & $-$ \\
\hline
$x-4$ & $-$ & $-$ & $-$ & $-$ & $0$ & $+$ \\
\hline
$B(x)$ & $+$ & $+$ & $0$ & $+$ & $0$ & $-$ \\
\hline
\end{tabular}
\end{center}

\textbf{Question 6 — Conclusion commerciale.}

$B(x) > 0$ pour $x \in [0\,;\,1[ \cup ]1\,;\,4[$. 
Rappel : $x$ est exprimé en \emph{dizaines de FCFA}. Cela correspond à un prix du manioc strictement compris entre \textbf{0 et 40 FCFA/kg}, la valeur \textbf{10 FCFA/kg} (correspondant à $x=1$) étant exclue.
Le bénéfice est \emph{nul} pour $x = 1$ (soit \textbf{10 FCFA/kg}) et $x = 4$ (soit \textbf{40 FCFA/kg}) : ce sont les \emph{seuils de rentabilité}. Au-delà de 40 FCFA/kg, le modèle prédit des pertes.

\textbf{Question 7 — Esprit critique.}

Un prix trop élevé fait fuir la clientèle vers les étals concurrents du marché Mokolo : les quantités vendues s'effondrent, et le bénéfice devient négatif (coûts fixes non couverts). Le modèle traduit mathématiquement cette réalité commerciale : la fonction bénéfice, concave, atteint un maximum puis décroît, reflétant la loi de l'offre et de la demande sur un marché concurrentiel.
\end{exoresolu}

\begin{attention}[Pièges fréquents dans cette activité]
\begin{itemize}
\item Ne pas oublier de \textbf{vérifier} la solution du système dans les trois équations de départ : une erreur de calcul dans le pivot se propage silencieusement.
\item Dans l'opération $(L_2') \leftarrow 2(L_2) - (L_1)$, multiplier \emph{tous} les termes de $(L_2)$, y compris le second membre.
\item Attention au signe dans le tableau : $B(x) = -(x-1)^2(x-4)$ commence par un signe \emph{moins} ; beaucoup d'élèves l'oublient et obtiennent l'inéquation inverse.
\item Le bénéfice est \emph{strictement} positif : les valeurs $x = 1$ et $x = 4$ doivent être \textbf{exclues} de l'ensemble solution.
\item \textbf{Conversion d'unités :} $x$ est en dizaines de FCFA. $x=1 \Leftrightarrow 10$ FCFA, $x=4 \Leftrightarrow 40$ FCFA. Ne pas confondre la valeur de $x$ et le prix en FCFA.
\end{itemize}
\end{attention}

\courssub{Bilan de l'activité}

Cette activité a permis de mettre en évidence les points suivants :

\begin{itemize}
\item Une situation commerciale concrète (tickets de vente) se \cle{traduit} naturellement en un système linéaire $3\times 3$ : les mathématiques deviennent un outil d'\emph{investigation} pour retrouver une information perdue.
\item La méthode du \cle{pivot de Gauss} est un algorithme puissant et systématique : elle transforme un système « enchevêtré » en système triangulaire que l'on résout mécaniquement par remontée.
\item L'étude d'un polynôme de degré 3 suit toujours la même stratégie : \emph{racine évidente} $\rightarrow$ \emph{factorisation} $\rightarrow$ \emph{tableau de signes} $\rightarrow$ \emph{conclusion}. La factorisation $B(x) = -(x-1)^2(x-4)$ montre qu'une racine peut être \emph{double}, ce qui influence le tableau de signes (le signe ne change pas en $x = 1$).
\item Une solution mathématique n'a de valeur que si elle est \emph{interprétée} : les seuils $x = 1$ et $x = 4$ deviennent des seuils de rentabilité, et la fourchette de prix conseillée aux commerçants de Mokolo est une décision économique fondée sur un raisonnement rigoureux.
\item Enfin, la rédaction du mini-rapport et la présentation orale développent des compétences transversales essentielles : communiquer, argumenter, vérifier et travailler en équipe — des savoir-être attendus au Baccalauréat comme dans la vie professionnelle.
\end{itemize}

\begin{savaistu}[Le marché Mokolo et les mathématiques]
Les grands marchés comme Mokolo à Yaoundé ou Mboppi à Douala fonctionnent sur des équilibres de prix que les économistes modélisent précisément par des systèmes d'équations. La méthode du pivot que tu viens d'utiliser porte le nom du mathématicien allemand \emph{Carl Friedrich Gauss} (1777--1855) ; elle est aujourd'hui au cœur des logiciels qui gèrent les stocks des supermarchés, les réseaux de distribution d'ENEO ou la planification des flottes de Camrail. Les systèmes que tu résous à la main avec trois inconnues, les ordinateurs les résolvent avec des millions d'inconnues — mais c'est exactement le même algorithme !
\end{savaistu}

\newpage

\coursec{Méthodes et exercices résolus}

\courssub{Méthode 1 : Racine évidente et factorisation d'un polynôme de degré 3}

\begin{methode}[Factoriser un polynôme de degré 3 à l'aide d'une racine évidente]
Soit $P(x)=ax^3+bx^2+cx+d$ un polynôme de degré 3 ($a\neq 0$).
\begin{enumerate}
\item \textbf{Chercher une racine évidente} $\alpha$ parmi les valeurs simples : $0$, $1$, $-1$, $2$, $-2$, $3$, $-3$. On calcule $P(\alpha)$ jusqu'à trouver $P(\alpha)=0$, et on \textbf{écrit ce calcul} sur la copie.
\item \textbf{Poser la factorisation} : $P(x)=(x-\alpha)(a'x^2+b'x+c')$. On sait déjà que $a'=a$ (coefficient dominant) et que $-\alpha c'=d$ (termes constants).
\item \textbf{Déterminer $b'$ et $c'$ par identification} : développer $(x-\alpha)(ax^2+b'x+c')$, regrouper par puissances de $x$, puis identifier les coefficients avec ceux de $P(x)$. On obtient un petit système que l'on résout. (Alternative acceptée : la division euclidienne de $P(x)$ par $x-\alpha$.)
\item \textbf{Vérifier} en redéveloppant mentalement ou en testant une valeur (par exemple $x=0$ donne $P(0)=-\alpha c'$).
\item \textbf{Étudier le trinôme} $ax^2+b'x+c'$ : calculer son discriminant $\Delta$ pour obtenir, si possible, une factorisation complète de $P$.
\end{enumerate}
Réflexe : si le polynôme n'a pas de terme constant ($d=0$), alors $0$ est racine évidente et on factorise directement par $x$.
\end{methode}

\begin{exoresolu}[Factorisation d'un polynôme de degré 3]
Soit $P(x)=x^3-4x^2+x+6$.
\begin{enumerate}
\item Montrer que $-1$ est une racine de $P$.
\item Déterminer les réels $a$, $b$, $c$ tels que $P(x)=(x+1)(ax^2+bx+c)$.
\item En déduire la factorisation complète de $P(x)$ dans $\mathbb{R}$.
\end{enumerate}
\tcblower
\textbf{1.} Calculons $P(-1)$ :
\[ P(-1)=(-1)^3-4(-1)^2+(-1)+6=-1-4-1+6=0. \]
Donc $-1$ est bien une racine de $P$.

\textbf{2.} Puisque $P(-1)=0$, $P(x)$ se factorise par $(x-(-1))=(x+1)$. Développons :
\begin{align*}
(x+1)(ax^2+bx+c) &= ax^3+bx^2+cx+ax^2+bx+c\\
&= ax^3+(a+b)x^2+(b+c)x+c.
\end{align*}
Par identification avec $P(x)=x^3-4x^2+x+6$, on obtient le système :
\[
\begin{cases}
a=1\\
a+b=-4\\
b+c=1\\
c=6
\end{cases}
\quad\Longrightarrow\quad
a=1,\quad b=-4-a=-5,\quad c=6.
\]
Vérifions la troisième équation : $b+c=-5+6=1$ : c'est cohérent. Donc :
\[ P(x)=(x+1)(x^2-5x+6). \]

\textbf{3.} Factorisons le trinôme $x^2-5x+6$. Son discriminant est :
\[ \Delta=(-5)^2-4\times 1\times 6=25-24=1>0. \]
Il admet deux racines réelles :
\[ x_1=\frac{5-\sqrt{1}}{2}=2 \qquad\text{et}\qquad x_2=\frac{5+\sqrt{1}}{2}=3. \]
Donc $x^2-5x+6=(x-2)(x-3)$. Conclusion :
\[ \boxed{P(x)=(x+1)(x-2)(x-3)} \]
$P$ admet trois racines réelles : $-1$, $2$ et $3$.
\end{exoresolu}

\courssub{Méthode 2 : Résoudre une équation de degré 3}

\begin{methode}[Résoudre une équation de degré 3 dans $\mathbb{R}$]
Pour résoudre $P(x)=0$ où $P$ est de degré 3 :
\begin{enumerate}
\item \textbf{Factoriser} $P$ complètement (racine évidente puis trinôme, comme dans la méthode 1).
\item \textbf{Appliquer la règle du produit nul} : un produit de facteurs est nul si et seulement si l'un au moins de ses facteurs est nul.
\item \textbf{Résoudre} chaque équation du premier degré ou du second degré obtenue.
\item \textbf{Conclure} en donnant l'ensemble des solutions $S=\{\ldots\}$, en vérifiant qu'aucune solution n'est comptée deux fois.
\end{enumerate}
Cas particuliers à reconnaître : équation bicarrée déguisée, identité remarquable $x^3-\alpha^3=(x-\alpha)(x^2+\alpha x+\alpha^2)$.
\end{methode}

\begin{exoresolu}[Équation de degré 3 (type Bac)]
On considère l'équation $(E)$ : $2x^3-3x^2-3x+2=0$.
\begin{enumerate}
\item Vérifier que $2$ est solution de $(E)$.
\item Résoudre $(E)$ dans $\mathbb{R}$.
\end{enumerate}
\tcblower
\textbf{1.} Calculons :
\[ 2(2)^3-3(2)^2-3(2)+2=16-12-6+2=0. \]
Donc $2$ est bien solution de $(E)$.

\textbf{2.} On factorise par $(x-2)$ : il existe $a$, $b$, $c$ tels que
\[ 2x^3-3x^2-3x+2=(x-2)(ax^2+bx+c). \]
En développant le membre de droite : $ax^3+(b-2a)x^2+(c-2b)x-2c$. Par identification :
\[
\begin{cases}
a=2\\
b-2a=-3\\
c-2b=-3\\
-2c=2
\end{cases}
\quad\Longrightarrow\quad
a=2,\quad c=-1,\quad b=-3+2a=1.
\]
Vérification de la troisième ligne : $c-2b=-1-2=-3$ : correct. Donc :
\[ (E)\iff (x-2)(2x^2+x-1)=0. \]
Un produit de facteurs est nul si et seulement si l'un des facteurs est nul :
\begin{itemize}
\item $x-2=0\iff x=2$ ;
\item $2x^2+x-1=0$ : $\Delta=1^2-4\times 2\times(-1)=1+8=9>0$, d'où
\[ x=\frac{-1-3}{4}=-1 \qquad\text{ou}\qquad x=\frac{-1+3}{4}=\frac{1}{2}. \]
\end{itemize}
Conclusion : l'équation $(E)$ admet trois solutions réelles :
\[ S=\left\{-1\ ;\ \frac{1}{2}\ ;\ 2\right\}. \]
\end{exoresolu}

\courssub{Méthode 3 : Résoudre une inéquation de degré 3 (tableau de signes)}

\begin{methode}[Résoudre une inéquation polynomiale par tableau de signes]
Pour résoudre $P(x)\geq 0$ (ou $\leq 0$, $>0$, $<0$) où $P$ est de degré 3 :
\begin{enumerate}
\item \textbf{Se ramener à une comparaison à zéro} : transporter tous les termes dans un seul membre.
\item \textbf{Factoriser complètement} $P(x)$ en produit de facteurs du premier degré (et éventuellement d'un trinôme de signe constant).
\item \textbf{Dresser le tableau de signes} : une ligne par facteur, avec les racines rangées dans l'ordre croissant ; la dernière ligne donne le signe du produit par la règle des signes.
\item \textbf{Lire la solution} sur la dernière ligne. Attention aux bornes : crochet fermé si l'inégalité est large ($\geq$, $\leq$), crochet ouvert si elle est stricte ($>$, $<$).
\end{enumerate}
Rappel : un trinôme $ax^2+bx+c$ avec $\Delta<0$ est du signe de $a$ sur $\mathbb{R}$ tout entier.
\end{methode}

\begin{exoresolu}[Inéquation de degré 3]
Résoudre dans $\mathbb{R}$ l'inéquation $(I)$ : $x^3-4x^2+x+6\leq 0$.
\tcblower
On reconnaît le polynôme $P(x)=x^3-4x^2+x+6$ factorisé dans le premier exercice résolu :
\[ P(x)=(x+1)(x-2)(x-3). \]
Les trois facteurs s'annulent respectivement en $-1$, $2$ et $3$. Dressons le tableau de signes :
\begin{center}
\begin{tabular}{|c|ccccccc|}
\hline
$x$ & $-\infty$ & & $-1$ & & $2$ & & $3$ \\
\hline
$x+1$ & & $-$ & $0$ & $+$ & $+$ & $+$ & $+$ \\
\hline
$x-2$ & & $-$ & $-$ & $-$ & $0$ & $+$ & $+$ \\
\hline
$x-3$ & & $-$ & $-$ & $-$ & $-$ & $-$ & $0$ \\
\hline
$P(x)$ & & $-$ & $0$ & $+$ & $0$ & $-$ & $0$ \\
\hline
\end{tabular}
\end{center}
Précisons les intervalles : sur $]-\infty;-1[$ le produit est $(-)(-)(-)=-$ ; sur $]-1;2[$ : $(+)(-)(-)=+$ ; sur $]2;3[$ : $(+)(+)(-)=-$ ; sur $]3;+\infty[$ : $(+)(+)(+)=+$.

On cherche $P(x)\leq 0$, c'est-à-dire les intervalles où la dernière ligne porte le signe $-$, \textbf{bornes incluses} car l'inégalité est large et $P$ y est défini. Conclusion :
\[ S=\left]-\infty\,;\,-1\right]\cup\left[2\,;\,3\right]. \]
\end{exoresolu}

\courssub{Méthode 4 : Résolution par inconnue auxiliaire}

\begin{methode}[Poser une inconnue auxiliaire]
Certaines équations ou inéquations se ramènent à un degré inférieur par un changement d'inconnue.
\begin{enumerate}
\item \textbf{Repérer la forme} : l'équation contient une même « brique » répétée, par exemple $x^2$ et $x^4$ (équation bicarrée), ou $(x-1)^2$ et $(x-1)^4$, ou $\sqrt{x}$ et $x$.
\item \textbf{Poser} $X=$ la brique de base (par exemple $X=x^2$), \textbf{en précisant le domaine} de $X$ (si $X=x^2$ alors $X\geq 0$ ; si $X=\sqrt{x}$ alors $X\geq 0$ et $x\geq 0$).
\item \textbf{Résoudre} l'équation en $X$ (second degré le plus souvent).
\item \textbf{Écarter} les valeurs de $X$ hors domaine, puis \textbf{revenir à l'inconnue initiale} $x$ en résolvant $x^2=X$ (ou $\sqrt{x}=X$, etc.).
\item \textbf{Conclure} avec l'ensemble des solutions en $x$.
\end{enumerate}
\end{methode}

\begin{exoresolu}[Équation bicarrée]
Résoudre dans $\mathbb{R}$ l'équation $(E)$ : $x^4-5x^2+4=0$.
\tcblower
L'équation ne contient que des puissances paires de $x$ : c'est une équation \cle{bicarrée}.

\textbf{Étape 1 : changement d'inconnue.} Posons $X=x^2$, avec la condition $X\geq 0$. Alors $x^4=X^2$ et :
\[ (E)\iff X^2-5X+4=0. \]

\textbf{Étape 2 : résolution en $X$.} Discriminant : $\Delta=(-5)^2-4\times 1\times 4=25-16=9>0$. Deux racines :
\[ X_1=\frac{5-3}{2}=1 \qquad\text{et}\qquad X_2=\frac{5+3}{2}=4. \]

\textbf{Étape 3 : validité et retour à $x$.} Les deux valeurs $1$ et $4$ sont positives : toutes deux sont acceptables.
\begin{itemize}
\item $x^2=1\iff x=1$ ou $x=-1$ ;
\item $x^2=4\iff x=2$ ou $x=-2$.
\end{itemize}

\textbf{Conclusion.} L'équation admet quatre solutions réelles :
\[ S=\{-2\ ;\ -1\ ;\ 1\ ;\ 2\}. \]
Vérification rapide pour $x=2$ : $16-20+4=0$ : correct.
\end{exoresolu}

\courssub{Méthode 5 : Résoudre un système linéaire 2$\times$2}

\begin{methode}[Trois techniques pour un système 2$\times$2]
Soit le système $\begin{cases} ax+by=e\\ cx+dy=f \end{cases}$.
\begin{enumerate}
\item \textbf{Substitution} : exprimer une inconnue en fonction de l'autre dans l'équation la plus simple (coefficient $1$ ou $-1$ de préférence), puis reporter dans l'autre équation.
\item \textbf{Combinaison (addition)} : multiplier les équations par des coefficients bien choisis pour éliminer une inconnue par addition membre à membre, puis résoudre et remonter.
\item \textbf{Déterminant} : calculer $\Delta=\begin{vmatrix} a & b\\ c & d\end{vmatrix}=ad-bc$.
\begin{itemize}
\item Si $\Delta\neq 0$ : solution unique $x=\dfrac{ed-bf}{\Delta}$, $y=\dfrac{af-ec}{\Delta}$.
\item Si $\Delta=0$ : le système a soit aucune solution, soit une infinité (droites confondues) ; on tranche en comparant les équations.
\end{itemize}
\item \textbf{Toujours vérifier} le couple trouvé dans les deux équations de départ.
\end{enumerate}
\end{methode}

\begin{exoresolu}[Système 2$\times$2 en contexte]
À la gare routière de Bafoussam, une coopérative achète des sacs d'engrais et des sacs de semences. Un premier chargement de $3$ sacs d'engrais et $2$ sacs de semences coûte $\num{58000}$ FCFA ; un second chargement de $5$ sacs d'engrais et $4$ sacs de semences coûte $\num{108000}$ FCFA. Déterminer le prix d'un sac d'engrais et celui d'un sac de semences.
\tcblower
\textbf{Mise en équations.} Soit $x$ le prix (en milliers de FCFA) d'un sac d'engrais et $y$ celui d'un sac de semences. Les données se traduisent par :
\[
\begin{cases}
3x+2y=58 & (L_1)\\
5x+4y=108 & (L_2)
\end{cases}
\]

\textbf{Résolution par combinaison.} Éliminons $y$ : calculons $2\times(L_1)-(L_2)$ :
\[ 2(3x+2y)-(5x+4y)=2\times 58-108 \iff 6x+4y-5x-4y=116-108 \iff x=8. \]
Reportons $x=8$ dans $(L_1)$ :
\[ 3\times 8+2y=58 \iff 2y=58-24=34 \iff y=17. \]

\textbf{Vérification.} Dans $(L_2)$ : $5\times 8+4\times 17=40+68=108$ : correct.

\textbf{Contrôle par le déterminant.} $\Delta=3\times 4-2\times 5=12-10=2\neq 0$ : la solution est bien unique, et $x=\dfrac{58\times 4-2\times 108}{2}=\dfrac{232-216}{2}=8$, $y=\dfrac{3\times 108-58\times 5}{2}=\dfrac{324-290}{2}=17$.

\textbf{Conclusion.} Un sac d'engrais coûte $\num{8000}$ FCFA et un sac de semences coûte $\num{17000}$ FCFA.
\end{exoresolu}

\courssub{Méthode 6 : Résoudre un système linéaire 3$\times$3 par le pivot de Gauss}

\begin{methode}[Pivot de Gauss pour un système 3$\times$3]
\begin{enumerate}
\item \textbf{Choisir un pivot} : une équation où le coefficient de $x$ est simple (idéalement $1$), que l'on place en première ligne $(L_1)$.
\item \textbf{Éliminer $x$} des lignes $(L_2)$ et $(L_3)$ par des combinaisons du type $L_2\leftarrow L_2-aL_1$, en écrivant explicitement chaque opération effectuée.
\item \textbf{Éliminer $y$} de la troisième ligne à l'aide de la nouvelle deuxième ligne : on obtient un système \cle{triangulaire}.
\item \textbf{Remonter} : la dernière équation donne $z$, que l'on reporte dans la deuxième pour obtenir $y$, puis dans la première pour obtenir $x$.
\item \textbf{Vérifier} le triplet dans les trois équations initiales.
\end{enumerate}
Si à une étape une ligne devient $0=0$, le système a une infinité de solutions ; si elle devient $0=k$ avec $k\neq 0$, il n'a aucune solution.
\end{methode}

\begin{exoresolu}[Pivot de Gauss (type Bac)]
Résoudre dans $\mathbb{R}^3$ le système :
\[
(S)\quad
\begin{cases}
x+y+z=6 & (L_1)\\
2x-y+z=3 & (L_2)\\
x+2y-z=2 & (L_3)
\end{cases}
\]
\tcblower
\textbf{Étape 1 : élimination de $x$ dans $(L_2)$ et $(L_3)$.}

$L_2\leftarrow L_2-2L_1$ : $(2x-y+z)-(2x+2y+2z)=3-12$, soit $-3y-z=-9$, c'est-à-dire $3y+z=9$.

$L_3\leftarrow L_3-L_1$ : $(x+2y-z)-(x+y+z)=2-6$, soit $y-2z=-4$.

Le système équivaut à :
\[
\begin{cases}
x+y+z=6 & (L_1)\\
3y+z=9 & (L_2')\\
y-2z=-4 & (L_3')
\end{cases}
\]

\textbf{Étape 2 : élimination de $y$ dans $(L_3')$.}

$L_3'\leftarrow 3L_3'-L_2'$ : $3(y-2z)-(3y+z)=3\times(-4)-9$, soit $-7z=-21$, d'où $z=3$.

\textbf{Étape 3 : remontée.}

Dans $(L_2')$ : $3y+3=9\iff 3y=6\iff y=2$.

Dans $(L_1)$ : $x+2+3=6\iff x=1$.

\textbf{Vérification.} $(L_2)$ : $2\times 1-2+3=3$ : correct. $(L_3)$ : $1+2\times 2-3=2$ : correct.

\textbf{Conclusion.} Le système admet une unique solution :
\[ S=\{(1\,;\,2\,;\,3)\}. \]
\end{exoresolu}

\courssub{Méthode 7 : Modéliser un problème concret}

\begin{methode}[Mettre un problème en équation ou en système]
\begin{enumerate}
\item \textbf{Choisir les inconnues} : les nommer explicitement avec leur unité (« Soit $x$ le nombre de... »).
\item \textbf{Traduire chaque information} de l'énoncé en une égalité ou une inégalité (une information = une équation).
\item \textbf{Résoudre} avec la technique adaptée (factorisation, tableau de signes, pivot de Gauss...).
\item \textbf{Filtrer} les solutions : dans un problème concret, on ne garde souvent que les solutions entières, positives, ou dans un intervalle donné.
\item \textbf{Rédiger la réponse} en langage courant, avec les unités, et vérifier la cohérence avec l'énoncé.
\end{enumerate}
\end{methode}

\begin{exoresolu}[Problème de modélisation (type Bac)]
Un cybercafé de Yaoundé propose trois formules de connexion Internet : la formule A à $x$ FCFA l'heure, la formule B à $y$ FCFA l'heure et la formule C à $z$ FCFA l'heure. Trois clients paient respectivement :
\begin{itemize}
\item Client 1 : $1$ h en A, $1$ h en B et $1$ h en C pour $\num{1500}$ FCFA ;
\item Client 2 : $2$ h en A et $1$ h en B pour $\num{1400}$ FCFA ;
\item Client 3 : $1$ h en B et $2$ h en C pour $\num{1600}$ FCFA.
\end{itemize}
Déterminer le tarif horaire de chaque formule.
\tcblower
\textbf{Mise en équations.} Les trois informations se traduisent par le système :
\[
\begin{cases}
x+y+z=1500 & (L_1)\\
2x+y=1400 & (L_2)\\
y+2z=1600 & (L_3)
\end{cases}
\]

\textbf{Étape 1 : élimination de $x$.} $L_2\leftarrow L_2-2L_1$ :
\[ (2x+y)-(2x+2y+2z)=1400-3000 \iff -y-2z=-1600 \iff y+2z=1600 \quad (L_2'). \]

\textbf{Étape 2 : comparaison avec $(L_3)$.} $L_3-L_2'$ : $(y+2z)-(y+2z)=1600-1600$, soit $0=0$.

L'égalité $0=0$ est toujours vraie : les équations $(L_2')$ et $(L_3)$ sont \textbf{redondantes}. Le système se réduit à deux équations indépendantes pour trois inconnues : il admet une \textbf{infinité de solutions}, que l'on exprime à l'aide d'un paramètre. Posons $z=t$ ($t\in\mathbb{R}$) :
\[ y=1600-2t \quad\text{et}\quad x=1500-y-z=1500-(1600-2t)-t=t-100. \]
Donc $S=\{(t-100\,;\,1600-2t\,;\,t),\ t\in\mathbb{R}\}$.

\textbf{Étape 3 : exploitation concrète.} Les tarifs doivent être strictement positifs : $t-100>0$ et $1600-2t>0$, c'est-à-dire $100<t<800$. Le gérant précise alors que la formule C coûte $600$ FCFA l'heure : $t=600$, d'où $x=500$ et $y=400$.

\textbf{Vérification.} $(L_1)$ : $500+400+600=1500$ : correct. $(L_2)$ : $2\times 500+400=1400$ : correct. $(L_3)$ : $400+2\times 600=1600$ : correct.

\textbf{Conclusion.} La formule A coûte $500$ FCFA l'heure, la formule B $400$ FCFA l'heure et la formule C $600$ FCFA l'heure.

\textbf{Leçon à retenir.} Cet exercice illustre un point essentiel du programme : un système 3$\times$3 peut avoir une solution unique, \emph{aucune} solution (le pivot aboutit à $0=k$ avec $k\neq 0$ : équations incompatibles, $S=\varnothing$) ou une \emph{infinité} de solutions (le pivot aboutit à $0=0$ : équations redondantes, on paramètre). Ici, les deux derniers relevés de la cybercafé contenaient la même information : il a fallu une donnée supplémentaire pour déterminer les tarifs.
\end{exoresolu}

\begin{attention}[Les 5 erreurs qui coûtent des points au Bac]
\begin{enumerate}
\item \textbf{Oublier de vérifier la racine évidente avant de factoriser.} Affirmer « $P(x)=(x-2)(ax^2+bx+c)$ » sans avoir calculé et écrit $P(2)=0$ fait perdre le point de justification. Toujours montrer le calcul de $P(\alpha)$ explicitement.
\item \textbf{Se tromper de signe dans le facteur.} Si la racine est $\alpha=-1$, le facteur est $(x-(-1))=(x+1)$, et non $(x-1)$. C'est l'erreur de signe la plus fréquente dans les identifications de coefficients.
\item \textbf{Mal gérer les bornes dans une inéquation.} Avec $\leq$ ou $\geq$, les racines du polynôme \emph{sont} solutions (crochets fermés) ; avec $<$ ou $>$, elles ne le sont pas (crochets ouverts). De plus, on ne résout une inéquation par tableau de signes qu'après avoir tout ramené dans un membre : on ne compare jamais un produit à un nombre non nul facteur par facteur.
\item \textbf{Oublier le domaine de l'inconnue auxiliaire.} Si l'on pose $X=x^2$, toute valeur $X<0$ trouvée doit être rejetée ; et il faut revenir à $x$ jusqu'au bout : répondre « $X=4$ » au lieu de « $x=\pm 2$ » ne donne pas tous les points.
\item \textbf{Ne pas vérifier la solution d'un système et mal rédiger le pivot de Gauss.} Après substitution ou combinaison, il faut reporter le couple (ou le triplet) dans \emph{toutes} les équations initiales. Au pivot de Gauss, chaque opération doit être annoncée ($L_2\leftarrow L_2-2L_1$) : une ligne modifiée sans indication est sanctionnée, et une erreur de calcul non détectée fausse tout le système.
\end{enumerate}
\end{attention}

\newpage

\coursec{Exercices}

\coursec{Équations, inéquations et systèmes}

\courssub{Équations de degré 3 : racine évidente et factorisation}

\begin{definition}[Polynôme de degré 3]
Un polynôme de degré 3 est une fonction polynôme de la forme 
$P(x)=ax^3+bx^2+cx+d$ avec $a,b,c,d\in\mathbb{R}$ et $a\neq 0$.
\end{definition}

\begin{propriete}[Racine évidente et factorisation]
Soit $P(x)=ax^3+bx^2+cx+d$ un polynôme à coefficients entiers. 
Si $P$ admet une racine rationnelle $\frac{p}{q}$ sous forme irréductible, alors $p$ divise $d$ et $q$ divise $a$. 
En particulier, si $a=1$ (polynôme unitaire), toute racine rationnelle est un \textbf{diviseur du terme constant $d$}.
Si $\alpha$ est une racine de $P$ ($P(\alpha)=0$), alors il existe un unique polynôme $Q$ de degré 2 tel que 
$P(x)=(x-\alpha)Q(x)$.
\end{propriete}

\begin{methode}[Résoudre une équation de degré 3 $P(x)=0$]
\begin{enumerate}
\item Chercher une \cle{racine évidente} $\alpha$ en testant les diviseurs du terme constant (généralement $\pm 1, \pm 2, \pm 3, \dots$).
\item Factoriser $P(x)$ par $(x-\alpha)$ :
\begin{itemize}
\item soit par \textbf{identification des coefficients} : écrire $P(x)=(x-\alpha)(ax^2+bx+c)$, développer et identifier ;
\item soit par \textbf{division euclidienne} de $P(x)$ par $(x-\alpha)$.
\end{itemize}
\item Résoudre l'équation du second degré $Q(x)=0$ (discriminant $\Delta$, racines $\frac{-b\pm\sqrt{\Delta}}{2a}$).
\item Conclure : l'ensemble des solutions est $\{\alpha\} \cup S_Q$ où $S_Q$ est l'ensemble des solutions de $Q(x)=0$.
\end{enumerate}
\end{methode}

\begin{exemplebox}[Résolution de $x^3-2x^2-5x+6=0$]
$P(x)=x^3-2x^2-5x+6$. Les diviseurs de $6$ sont $\pm 1, \pm 2, \pm 3, \pm 6$.
$P(1)=1-2-5+6=0$, donc $1$ est une racine.
Factorisation : $P(x)=(x-1)(x^2+bx+c)=x^3+(b-1)x^2+(c-b)x-c$.
Identification : 
$\begin{cases} b-1=-2 \Rightarrow b=-1 \\ c-b=-5 \Rightarrow c-(-1)=-5 \Rightarrow c=-6 \\ -c=6 \Rightarrow c=-6 \text{ (cohérent)} \end{cases}$
$P(x)=(x-1)(x^2-x-6)=(x-1)(x-3)(x+2)$.
Solutions : $S=\{-2; 1; 3\}$.
\end{exemplebox}

\courssub{Inéquations de degré 3 et tableaux de signes}

\begin{propriete}[Signe d'un produit de facteurs du premier degré]
Pour des réels $a_1 < a_2 < \dots < a_n$, le signe du produit $(x-a_1)(x-a_2)\dots(x-a_n)$ sur $\mathbb{R}$ s'étudie à l'aide d'un \textbf{tableau de signes} :
\begin{itemize}
\item Sur chaque intervalle $]-\infty, a_1[$, $]a_1, a_2[$, \dots, $]a_n, +\infty[$, chaque facteur $(x-a_i)$ a un signe constant.
\item Le signe du produit est le produit des signes (règle des signes).
\item Aux valeurs $a_i$, le produit est nul.
\end{itemize}
\end{propriete}

\begin{methode}[Résoudre une inéquation $P(x) > 0$ (ou $<, \ge, \le$)]
\begin{enumerate}
\item Factoriser complètement $P(x)$ en facteurs du premier degré (éventuellement un trinôme de signe constant si $\Delta < 0$).
\item Construire le tableau de signes : lignes pour chaque facteur, dernière ligne pour $P(x)$.
\item Lire la solution sur la dernière ligne en tenant compte du sens de l'inéquation (strict ou large) et des valeurs interdites s'il y a lieu (dénominateur, racine carrée\dots).
\end{enumerate}
\end{methode}

\begin{exemplebox}[Inéquation $x^3-x^2-4x+4 > 0$]
$P(x)=x^3-x^2-4x+4$. $P(1)=0 \Rightarrow P(x)=(x-1)(x^2-4)=(x-1)(x-2)(x+2)$.
Racines : $-2, 1, 2$.
\begin{center}
\begin{tabular}{|c|ccccccccc|}\hline
$x$ & $-\infty$ &  & $-2$ &  & $1$ &  & $2$ &  & $+\infty$ \\\hline
$x+2$ &  & $-$ & $0$ & $+$ &  & $+$ &  & $+$ &  \\
$x-1$ &  & $-$ &  & $-$ & $0$ & $+$ &  & $+$ &  \\
$x-2$ &  & $-$ &  & $-$ &  & $-$ & $0$ & $+$ &  \\\hline
$P(x)$ &  & $-$ & $0$ & $+$ & $0$ & $-$ & $0$ & $+$ &  \\\hline
\end{tabular}
\end{center}
$P(x)>0 \Leftrightarrow x \in ]-2, 1[ \cup ]2, +\infty[$.
\end{exemplebox}

\courssub{Résolution par inconnue auxiliaire}

\begin{methode}[Résolution par changement de variable (inconnue auxiliaire)]
\begin{enumerate}
\item Identifier une expression qui se répète ou une structure quadratique (ex: $x^4, x^2$ ; $x+\frac{1}{x}$ ; $\sqrt{x}$).
\item Poser $X = \text{cette expression}$ et exprimer l'équation en fonction de $X$.
\item Résoudre l'équation en $X$ (souvent du 1er ou 2nd degré).
\item Revenir à $x$ : résoudre $X = \text{expression en } x$ pour chaque solution $X$ trouvée, \textbf{en respectant les conditions de définition} ($X\ge 0$ si $X=x^2$ ou $X=\sqrt{x}$ ; $x\neq 0$ si $X=1/x$, etc.).
\item Vérifier les solutions dans l'équation initiale.
\end{enumerate}
\end{methode}

\begin{exemplebox}[Équation $x^4-13x^2+36=0$]
Poser $X=x^2 \ge 0$. L'équation devient $X^2-13X+36=0$.
$\Delta=169-144=25$. $X_1=\frac{13-5}{2}=4$, $X_2=\frac{13+5}{2}=9$.
$X=4 \Rightarrow x^2=4 \Rightarrow x=\pm 2$.
$X=9 \Rightarrow x^2=9 \Rightarrow x=\pm 3$.
$S=\{-3; -2; 2; 3\}$.
\end{exemplebox}

\begin{exemplebox}[Équation $x+\frac{1}{x}=\frac{5}{2}$]
Condition : $x\neq 0$.
Multiplier par $2x$ : $2x^2+2=5x \Leftrightarrow 2x^2-5x+2=0$.
$\Delta=25-16=9$. $x_1=\frac{5-3}{4}=\frac{1}{2}$, $x_2=\frac{5+3}{4}=2$.
Les deux sont $\neq 0$. $S=\{\frac{1}{2}; 2\}$.
\end{exemplebox}

\courssub{Systèmes linéaires 2$\times$2}

\begin{definition}[Système linéaire 2$\times$2]
Un système de deux équations à deux inconnues $x,y$ s'écrit 
$\begin{cases} ax+by=c \\ a'x+b'y=c' \end{cases}$ avec $(a,b)\neq(0,0)$ et $(a',b')\neq(0,0)$.
Son \cle{déterminant} est $\Delta = ab' - a'b$.
\end{definition}

\begin{propriete}[Nombre de solutions (Cramer)]
\begin{itemize}
\item Si $\Delta \neq 0$ : le système admet une \textbf{unique solution} 
$x=\frac{cb'-c'b}{\Delta}$, $y=\frac{ac'-a'c}{\Delta}$.
\item Si $\Delta = 0$ : les droites sont parallèles ou confondues.
\begin{itemize}
\item Si les équations sont proportionnelles (même rapport pour les coefficients et les seconds membres) : \textbf{infinité de solutions} (droites confondues).
\item Sinon : \textbf{aucune solution} (droites parallèles distinctes).
\end{itemize}
\end{itemize}
\end{propriete}

\begin{methode}[Trois méthodes de résolution]
\begin{enumerate}
\item \textbf{Substitution} : exprimer $x$ (ou $y$) dans une équation, substituer dans l'autre.
\item \textbf{Combinaison linéaire (élimination)} : multiplier les équations pour obtenir des coefficients opposés sur une inconnue, additionner pour l'éliminer.
\item \textbf{Déterminant (Cramer)} : calculer $\Delta, \Delta_x, \Delta_y$ et appliquer les formules si $\Delta\neq 0$.
\end{enumerate}
\end{methode}

\courssub{Systèmes linéaires 3$\times$3 et pivot de Gauss}

\begin{methode}[Méthode du pivot de Gauss (échelonnement)]
Pour un système 
$\begin{cases} a_1x+b_1y+c_1z=d_1 \\ a_2x+b_2y+c_2z=d_2 \\ a_3x+b_3y+c_3z=d_3 \end{cases}$
\begin{enumerate}
\item Choisir un \textbf{pivot} (coefficient non nul, idéalement $1$ ou $-1$) sur la première inconnue ($x$) dans la première équation. Échanger les lignes si nécessaire.
\item Éliminer $x$ des deux autres équations : remplacer $L_2$ par $L_2 - \frac{a_2}{a_1}L_1$ et $L_3$ par $L_3 - \frac{a_3}{a_1}L_1$.
\item Sur le système $2\times 2$ obtenu en $y,z$, choisir un pivot sur $y$ (deuxième équation) et éliminer $y$ de la troisième.
\item On obtient un système \textbf{triangulaire} (en échelon). Résoudre par \textbf{remontée} : trouver $z$, puis $y$, puis $x$.
\end{enumerate}
\end{methode}

\begin{attention}[Pièges du pivot de Gauss]
\begin{itemize}
\item Ne jamais diviser par $0$ : si le pivot est $0$, échanger avec une ligne du dessous qui a un coefficient non nul.
\item Garder les fractions exactes le plus longtemps possible (éviter les approximations décimales).
\item Vérifier la solution dans les \textbf{trois} équations initiales.
\end{itemize}
\end{attention}

\courssub{Problèmes concrets et modélisation}

\begin{methode}[Démarche de modélisation]
\begin{enumerate}
\item \textbf{Variables} : identifier les grandeurs inconnues et leur donner un nom (lettre), avec leur unité.
\item \textbf{Contraintes} : noter les conditions (positivité, entier, domaine de définition\dots).
\item \textbf{Traduction} : écrire les relations mathématiques (équations, inéquations, système).
\item \textbf{Résolution} : appliquer les méthodes algébriques adaptées.
\item \textbf{Interprétation} : redonner du sens aux nombres trouvés (arrondir si entier, rejeter les solutions non conformes aux contraintes).
\item \textbf{Conclusion} : répondre par une phrase complète à la question posée.
\end{enumerate}
\end{methode}

\begin{savaistu}[Modélisation au Cameroun]
Les systèmes linéaires modélisent de nombreuses situations réelles : 
\begin{itemize}
\item \textbf{Transport} : calcul des tarifs (bus, taxi, train) à partir de recettes journalières (ex: ligne Douala–Yaoundé).
\item \textbf{Agriculture} : optimisation des achats d'intrants (engrais, semences) sous contrainte budgétaire (coopératives de l'Ouest, du Nord).
\item \textbf{Hydraulique} : dimensionnement de cuves, tuyaux, réseaux d'adduction d'eau (CAMWATER, SNEC).
\item \textbf{Télécommunications} : forfaits MTN/Orange (voix, data, SMS) à partir de factures.
\end{itemize}
La rigueur algébrique garantit des décisions techniques et économiques fiables.
\end{savaistu}

\courssub{Je vérifie mes connaissances}

\begin{exercice}[QCM : une seule bonne réponse]
Pour chaque question, choisir la bonne réponse (aucune justification n'est demandée).
\begin{enumerate}
\item Le nombre $2$ est racine évidente du polynôme $P(x)=x^3-3x^2-4x+12$ car :
\begin{enumerate}
\item $P(2)=0$ \quad \item $P(0)=2$ \quad \item $P(-2)=0$ \quad \item $P(2)=12$
\end{enumerate}
\item Si $P(x)=x^3-3x^2-4x+12=(x-2)(x^2+bx+c)$, alors :
\begin{enumerate}
\item $b=-1$ et $c=-6$ \quad \item $b=1$ et $c=6$ \quad \item $b=-5$ et $c=-6$ \quad \item $b=-1$ et $c=6$
\end{enumerate}
\item Le déterminant du système $\begin{cases} 2x-5y=1\\ 3x+y=4 \end{cases}$ vaut :
\begin{enumerate}
\item $-17$ \quad \item $17$ \quad \item $-13$ \quad \item $13$
\end{enumerate}
\item Pour résoudre l'équation $x^4-5x^2+4=0$, on pose comme inconnue auxiliaire :
\begin{enumerate}
\item $X=x$ \quad \item $X=x^2$ \quad \item $X=x^4$ \quad \item $X=\sqrt{x}$
\end{enumerate}
\end{enumerate}
\end{exercice}

\begin{exercice}[Vrai ou faux, en justifiant]
Répondre par Vrai ou Faux en justifiant chaque réponse par un calcul ou un argument.
\begin{enumerate}
\item L'équation $x^3+x^2+x+1=0$ admet $1$ pour racine évidente.
\item Si $(E)$ : $ax+by=c$ et $(E')$ : $a'x+b'y=c'$ avec $ab'-a'b\neq 0$, alors le système formé par $(E)$ et $(E')$ admet un unique couple solution.
\item L'équation $\dfrac{1}{x}+x=2$ est définie pour tout réel $x$.
\item Un système linéaire de trois équations à trois inconnues admet toujours une unique solution.
\end{enumerate}
\end{exercice}

\begin{exercice}[Racines évidentes]
Pour chacun des polynômes suivants, déterminer une racine évidente parmi les nombres $-3$ ; $-2$ ; $-1$ ; $0$ ; $1$ ; $2$ ; $3$.
\begin{enumerate}
\item $P(x)=x^3-4x^2+x+6$ ;
\item $Q(x)=x^3+3x^2-4$ ;
\item $R(x)=2x^3-5x^2-4x+3$.
\end{enumerate}
\end{exercice}

\begin{exercice}[Calculs directs]
\begin{enumerate}
\item Calculer le déterminant de chacun des systèmes suivants et préciser s'il admet un unique couple solution :
\begin{enumerate}
\item $\begin{cases} 4x-3y=2\\ x+2y=5 \end{cases}$ \qquad \item $\begin{cases} 6x-9y=1\\ 2x-3y=4 \end{cases}$
\end{enumerate}
\item Factoriser $x^2-7x+12$, puis en déduire le signe de $x^2-7x+12$ suivant les valeurs de $x$.
\item Donner l'ensemble de définition dans $\mathbb{R}$ de l'équation $\dfrac{x^2-1}{x-3}=2x+1$.
\end{enumerate}
\end{exercice}

\courssub{Je m'entraîne}

\begin{exercice}[Équation et inéquation de degré 3]
On considère le polynôme $P(x)=x^3+2x^2-5x-6$.
\begin{enumerate}
\item Vérifier que $-1$ est une racine de $P$.
\item Déterminer les réels $a$, $b$ et $c$ tels que $P(x)=(x+1)(ax^2+bx+c)$.
\item Résoudre dans $\mathbb{R}$ l'équation $P(x)=0$.
\item À l'aide d'un tableau de signes, résoudre dans $\mathbb{R}$ l'inéquation $x^3+2x^2-5x-6<0$.
\end{enumerate}
\end{exercice}

\begin{exercice}[Inconnue auxiliaire]
\begin{enumerate}
\item En posant $X=x^2$, résoudre dans $\mathbb{R}$ l'équation $x^4-13x^2+36=0$.
\item On considère l'équation $(E)$ : $x+\dfrac{1}{x}=\dfrac{5}{2}$.
\begin{enumerate}
\item Préciser la valeur interdite.
\item Après réduction au même dénominateur, montrer que $(E)$ équivaut à $2x^2-5x+2=0$ avec $x\neq 0$.
\item Résoudre $(E)$ dans $\mathbb{R}$.
\end{enumerate}
\end{enumerate}
\end{exercice}

\begin{exercice}[Système 2$\times$2 : les trois méthodes]
On considère le système $(S)$ : $\begin{cases} 3x-2y=7\\ 5x+y=16 \end{cases}$.
\begin{enumerate}
\item Résoudre $(S)$ par substitution.
\item Résoudre $(S)$ par combinaison linéaire (élimination).
\item Résoudre $(S)$ par la méthode du déterminant (formules de Cramer).
\item Vérifier que le couple obtenu satisfait les deux équations.
\end{enumerate}
\end{exercice}

\begin{exercice}[Système 3$\times$3 : pivot de Gauss]
On considère le système $(S)$ : $\begin{cases} x-y+z=2\\ 2x+y-z=1\\ x+2y+z=8 \end{cases}$.
\begin{enumerate}
\item En utilisant la première équation comme pivot, éliminer $x$ dans les deux dernières équations.
\item Achever la résolution par remontée et donner le triplet solution.
\item Vérifier que le triplet obtenu satisfait les trois équations.
\end{enumerate}
\end{exercice}

\courssub{Je me prépare au Bac}

\begin{exercice}[La cuve de stockage d'eau — type Bac]
Une entreprise d'adduction d'eau de la région du Littoral fabrique des cuves dont le volume, en mètres cubes, est modélisé par
\[ V(x)=x^3-4x^2+x+6, \]
où $x$ (exprimé en mètres) est une dimension caractéristique de la cuve, avec $x>0$.
\begin{enumerate}
\item \begin{enumerate}
\item Vérifier que $2$ est une racine du polynôme $V$.
\item Déterminer les réels $a$, $b$ et $c$ tels que $V(x)=(x-2)(ax^2+bx+c)$.
\item Résoudre dans $\mathbb{R}$ l'équation $V(x)=0$.
\end{enumerate}
\item \begin{enumerate}
\item Dresser le tableau de signes de $V(x)$ sur $\mathbb{R}$.
\item En déduire l'ensemble des valeurs de $x$ pour lesquelles le volume de la cuve est strictement positif.
\item Sachant que la cuve doit avoir un volume strictement positif et que sa dimension $x$ ne peut pas dépasser \SI{2,5}{m} pour des raisons d'encombrement, déterminer l'intervalle des valeurs de $x$ retenues par l'entreprise.
\end{enumerate}
\item L'entreprise souhaite que le volume de la cuve soit supérieur ou égal à \SI{6}{m^3}. Montrer que cela revient à résoudre l'inéquation $x(x^2-4x+1)\geq 0$, puis déterminer les valeurs de $x$ convenables (on donnera les bornes sous forme exacte).
\end{enumerate}
\end{exercice}

\begin{exercice}[Le transporteur de la ligne Douala–Yaoundé — type Bac]
Une agence de transport interurbain vend des tickets de trois catégories : \emph{Classique} (prix $x$ FCFA), \emph{Confort} (prix $y$ FCFA) et \emph{VIP} (prix $z$ FCFA). Les recettes de trois journées sont consignées dans le tableau suivant :

\begin{center}
\begin{tabularx}{\linewidth}{|l|X|X|X|X|}
\hline
\rowcolor{popblueL} Journée & Classique & Confort & VIP & Recette (FCFA) \\
\hline
Lundi & 40 & 20 & 5 & \num{225000} \\
\hline
Mardi & 30 & 25 & 10 & \num{237500} \\
\hline
Mercredi & 50 & 15 & 8 & \num{239000} \\
\hline
\end{tabularx}
\end{center}

\begin{enumerate}
\item Montrer que les prix $x$, $y$ et $z$ vérifient le système
\[ (S) : \begin{cases} 8x+4y+z=\num{45000}\\ 6x+5y+2z=\num{47500}\\ 50x+15y+8z=\num{239000} \end{cases} \]
(on divisera la première équation par $5$ et la deuxième par $5$).
\item Résoudre le système $(S)$ par la méthode du pivot de Gauss.
\item En déduire le prix de vente d'un ticket de chaque catégorie (en FCFA, valeur exacte ou approchée au franc près).
\item Le jeudi, l'agence vend uniquement des tickets \emph{Classique} et \emph{Confort}. Elle vend $12$ tickets \emph{Confort}. Déterminer le nombre minimal de tickets \emph{Classique} à vendre pour que la recette du jeudi dépasse \num{100000} FCFA (on résoudra une inéquation du premier degré et l'on tiendra compte du fait que le nombre de tickets est un entier naturel).
\end{enumerate}
\end{exercice}

\begin{exercice}[Étude d'un système et modélisation — type Bac]
\textbf{Partie A.} On considère le système $(S_k)$ : $\begin{cases} 2x+3y=5\\ 4x+6y=k \end{cases}$ où $k$ est un nombre réel fixé.
\begin{enumerate}
\item Calculer le déterminant du système $(S_k)$. Que peut-on en conclure ?
\item On prend $k=10$.
\begin{enumerate}
\item Montrer que la deuxième équation équivaut à $2x+3y=5$.
\item En déduire l'ensemble des couples $(x\,;\,y)$ solutions de $(S_{10})$ (on pourra exprimer $y$ en fonction de $x$).
\end{enumerate}
\item On prend $k=12$. Montrer que le système $(S_{12})$ n'admet aucune solution.
\end{enumerate}
\textbf{Partie B.} Une coopérative agricole de l'Ouest achète deux types d'engrais. Le sac d'engrais A coûte $x$ milliers de FCFA et le sac d'engrais B coûte $y$ milliers de FCFA. Un premier achat de $2$ sacs de A et $3$ sacs de B coûte $95$ milliers de FCFA ; un second achat de $5$ sacs de A et $4$ sacs de B coûte $160$ milliers de FCFA.
\begin{enumerate}
\item Traduire ces données par un système de deux équations à deux inconnues.
\item Résoudre ce système par la méthode du déterminant et donner le prix de chaque sac d'engrais en FCFA.
\item La coopérative dispose d'un budget de \num{500000} FCFA. Elle souhaite acheter $10$ sacs de A et $n$ sacs de B. Déterminer le nombre maximal $n$ de sacs de B qu'elle peut acheter sans dépasser son budget.
\end{enumerate}
\end{exercice}

\courssub{Corrigés détaillés}



\newpage

\coursec{Corrigés des exercices}

\begin{corrige}[Exercice 1 — QCM]
\begin{enumerate}
\item \textbf{Réponse a.} $P(2)=2^3-3\times 2^2-4\times 2+12=8-12-8+12=0$. Dire que $2$ est racine de $P$ signifie exactement que $P(2)=0$.
\item \textbf{Réponse a.} On développe : $(x-2)(x^2+bx+c)=x^3+(b-2)x^2+(c-2b)x-2c$. Par identification avec $x^3-3x^2-4x+12$ : $b-2=-3\Rightarrow b=-1$ ; $-2c=12\Rightarrow c=-6$ ; vérification du coefficient de $x$ : $c-2b=-6+2=-4$ \checkmark. Donc $b=-1$ et $c=-6$.
\item \textbf{Réponse b.} $\Delta=2\times 1-3\times(-5)=2+15=17$.
\item \textbf{Réponse b.} L'équation s'écrit $(x^2)^2-5(x^2)+4=0$ : on pose $X=x^2$ (avec la condition $X\ge 0$), ce qui ramène à $X^2-5X+4=0$.
\end{enumerate}
\end{corrige}

\begin{corrige}[Exercice 2 — Vrai ou Faux]
\begin{enumerate}
\item \textbf{Faux.} $P(1)=1+1+1+1=4\neq 0$, donc $1$ n'est pas racine. En revanche $P(-1)=-1+1-1+1=0$ : c'est $-1$ qui est racine évidente (et l'on vérifie que $x^3+x^2+x+1=(x+1)(x^2+1)$).
\item \textbf{Vrai.} Le déterminant $\Delta=ab'-a'b$ est non nul : d'après les formules de Cramer, le système admet un unique couple solution $\left(\dfrac{cb'-c'b}{\Delta}\,;\,\dfrac{ac'-a'c}{\Delta}\right)$.
\item \textbf{Faux.} Le terme $\dfrac{1}{x}$ impose $x\neq 0$. L'ensemble de définition est $\mathbb{R}^*$ (et non $\mathbb{R}$ tout entier).
\item \textbf{Faux.} Si le déterminant du système est nul, le système peut n'avoir aucune solution (équations incompatibles) ou une infinité de solutions (équations proportionnelles). Contre-exemple : $\begin{cases} x+y+z=1\\ x+y+z=2\\ x+y+z=3 \end{cases}$ n'a aucune solution, car un même triplet $(x,y,z)$ ne peut pas donner à $x+y+z$ trois valeurs différentes.
\end{enumerate}
\end{corrige}

\begin{corrige}[Exercice 3 — Racines évidentes]
\begin{enumerate}
\item $P(-1)=(-1)^3-4(-1)^2-(-1)+6=-1-4+1+6=2\neq 0$ : reprenons avec les valeurs usuelles. $P(1)=1-4-1+6=2\neq 0$ ; $P(-2)=-8-16+2+6=-16\neq 0$ ; $P(2)=8-16+2+6=0$. \textbf{Racine :} $2$. (On vérifie aussi $P(3)=27-36+3+6=0$ : $3$ est également racine.)
\item $Q(1)=1+3-4=0$. \textbf{Racine :} $1$.
\item $R(-1)=2(-1)^3-5(-1)^2-4(-1)+3=-2-5+4+3=0$. \textbf{Racine :} $-1$. (Vérification : $R(3)=54-45-12+3=0$ convient aussi.)
\end{enumerate}
\textbf{Astuce :} pour un polynôme de degré $3$, la somme nulle des coefficients signale la racine $1$ ; pour $R$, la somme des coefficients des termes de degré pair ($-5+3=-2$) égale l'opposée de celle des termes de degré impair prise avec le signe alterné ($2-4=-2$), c'est-à-dire $a-b+c-d=0$ : cela signale la racine $-1$.
\end{corrige}

\begin{corrige}[Exercice 4 — Calculs directs]
\begin{enumerate}
\item \begin{enumerate}
\item $\Delta=4\times 2-1\times(-3)=8+3=11\neq 0$ : le système admet \textbf{un unique couple solution}.
\item $\Delta=6\times(-3)-2\times(-9)=-18+18=0$ : pas de solution unique. Ici $6x-9y=3(2x-3y)$ mais $1\neq 3\times 4$ : les équations sont incompatibles, le système n'a \textbf{aucune solution}.
\end{enumerate}
\item $\Delta=(-7)^2-4\times 1\times 12=49-48=1$, racines $x_1=\frac{7-1}{2}=3$ et $x_2=\frac{7+1}{2}=4$, donc $x^2-7x+12=(x-3)(x-4)$. Le trinôme est du signe de $a=1>0$ à l'extérieur des racines : positif ou nul sur $]-\infty,3]\cup[4,+\infty[$ ; strictement négatif sur $]3,4[$ ; nul en $3$ et en $4$.
\item Le dénominateur impose $x-3\neq 0$, c'est-à-dire $x\neq 3$. Ensemble de définition : $\mathbb{R}\setminus\{3\}$.
\end{enumerate}
\end{corrige}

\begin{corrige}[Exercice 5 — Équation et inéquation de degré 3]
\begin{enumerate}
\item $P(-1)=(-1)^3+2(-1)^2-5(-1)-6=-1+2+5-6=0$ : $-1$ est bien racine de $P$.
\item On cherche $P(x)=(x+1)(ax^2+bx+c)=ax^3+(a+b)x^2+(b+c)x+c$. Par identification avec $x^3+2x^2-5x-6$ :
$a=1$ ; $a+b=2\Rightarrow b=1$ ; $b+c=-5\Rightarrow c=-6$ ; vérification du terme constant : $c=-6$ \checkmark.
\[ P(x)=(x+1)(x^2+x-6). \]
\item $x^2+x-6=0$ : $\Delta=1+24=25$, $x=\frac{-1\pm 5}{2}$, soit $x=2$ ou $x=-3$. Donc $P(x)=(x+1)(x-2)(x+3)$ et
\[ S=\{-3\,;\,-1\,;\,2\}. \]
\item Tableau de signes de $P(x)=(x+3)(x+1)(x-2)$ :
\begin{center}
\begin{tabular}{|c|ccccccccc|}\hline
$x$ & $-\infty$ &  & $-3$ &  & $-1$ &  & $2$ &  & $+\infty$ \\\hline
$x+3$ &  & $-$ & $0$ & $+$ &  & $+$ &  & $+$ &  \\
$x+1$ &  & $-$ &  & $-$ & $0$ & $+$ &  & $+$ &  \\
$x-2$ &  & $-$ &  & $-$ &  & $-$ & $0$ & $+$ &  \\\hline
$P(x)$ &  & $-$ & $0$ & $+$ & $0$ & $-$ & $0$ & $+$ &  \\\hline
\end{tabular}
\end{center}
Inéquation stricte : $P(x)<0 \Leftrightarrow x\in ]-\infty,-3[\,\cup\,]-1,2[$.
\end{enumerate}
\end{corrige}

\begin{corrige}[Exercice 6 — Inconnue auxiliaire]
\begin{enumerate}
\item Posons $X=x^2$ (donc $X\ge 0$). L'équation devient $X^2-13X+36=0$ : $\Delta=169-144=25$, $X_1=\frac{13-5}{2}=4$, $X_2=\frac{13+5}{2}=9$ (toutes deux $\ge 0$, donc acceptables).
$X=4\Rightarrow x^2=4\Rightarrow x=\pm 2$ ; $X=9\Rightarrow x^2=9\Rightarrow x=\pm 3$.
\[ S=\{-3\,;\,-2\,;\,2\,;\,3\}. \]
\item \begin{enumerate}
\item Valeur interdite : $x=0$ (à cause de $\frac{1}{x}$).
\item Pour $x\neq 0$ : $x+\dfrac{1}{x}=\dfrac{5}{2}\Leftrightarrow \dfrac{x^2+1}{x}=\dfrac{5}{2}\Leftrightarrow 2(x^2+1)=5x \Leftrightarrow 2x^2-5x+2=0$.
\item $\Delta=25-16=9$ ; $x_1=\frac{5-3}{4}=\frac{1}{2}$, $x_2=\frac{5+3}{4}=2$. Aucune des deux n'est la valeur interdite.
\[ S=\left\{\tfrac{1}{2}\,;\,2\right\}. \]
\end{enumerate}
\end{enumerate}
\end{corrige}

\begin{corrige}[Exercice 7 — Système 2$\times$2 : les trois méthodes]
\begin{enumerate}
\item \textbf{Substitution.} La deuxième équation donne $y=16-5x$. On reporte dans la première :
$3x-2(16-5x)=7 \Leftrightarrow 3x-32+10x=7 \Leftrightarrow 13x=39 \Leftrightarrow x=3$. Alors $y=16-15=1$.
\item \textbf{Combinaison.} On élimine $y$ : $L_1+2L_2$ donne $(3x-2y)+2(5x+y)=7+32$, soit $13x=39$, d'où $x=3$. Puis $y=16-5\times 3=1$.
\item \textbf{Cramer.} $\Delta=3\times 1-5\times(-2)=3+10=13\neq 0$.
$\Delta_x=7\times 1-16\times(-2)=7+32=39 \Rightarrow x=\frac{39}{13}=3$.
$\Delta_y=3\times 16-5\times 7=48-35=13 \Rightarrow y=\frac{13}{13}=1$.
\item Vérification : $3\times 3-2\times 1=9-2=7$ \checkmark ; $5\times 3+1=16$ \checkmark.
\[ S=\{(3\,;\,1)\}. \]
\end{enumerate}
\end{corrige}

\begin{corrige}[Exercice 8 — Système 3$\times$3 : pivot de Gauss]
\begin{enumerate}
\item Pivot : le coefficient de $x$ dans $L_1$ vaut $1$ (idéal).
$L_2\leftarrow L_2-2L_1$ : $(2x-2x)+(y+2y)+(-z-2z)=1-4$, soit $3y-3z=-3$, c'est-à-dire $y-z=-1$.
$L_3\leftarrow L_3-L_1$ : $(x-x)+(2y+y)+(z-z)=8-2$, soit $3y=6$.
Système échelonné : $\begin{cases} x-y+z=2\\ 3y=6\\ y-z=-1 \end{cases}$
\item \textbf{Remontée.} $3y=6\Rightarrow y=2$. Puis $2-z=-1\Rightarrow z=3$. Enfin $x-2+3=2\Rightarrow x=1$.
\[ S=\{(1\,;\,2\,;\,3)\}. \]
\item Vérification : $1-2+3=2$ \checkmark ; $2\times 1+2-3=1$ \checkmark ; $1+2\times 2+3=8$ \checkmark.
\end{enumerate}
\end{corrige}

\begin{corrige}[Exercice 9 — La cuve de stockage d'eau (type Bac)]
\textbf{Barème indicatif :} 1.a) 0,5 pt ; 1.b) 1 pt ; 1.c) 1 pt ; 2.a) 1,5 pts ; 2.b) 0,5 pt ; 2.c) 1 pt ; 3) 2,5 pts.
\begin{enumerate}
\item \begin{enumerate}
\item $V(2)=8-16+2+6=0$ : $2$ est racine de $V$.
\item $V(x)=(x-2)(ax^2+bx+c)=ax^3+(b-2a)x^2+(c-2b)x-2c$. Identification avec $x^3-4x^2+x+6$ : $a=1$ ; $b-2=-4\Rightarrow b=-2$ ; $-2c=6\Rightarrow c=-3$ ; vérification du coefficient de $x$ : $c-2b=-3+4=1$ \checkmark.
\[ V(x)=(x-2)(x^2-2x-3). \]
\item $x^2-2x-3=0$ : $\Delta=4+12=16$, $x=\frac{2\pm 4}{2}$, soit $x=3$ ou $x=-1$. Donc $V(x)=(x-2)(x-3)(x+1)$ et $S=\{-1\,;\,2\,;\,3\}$.
\end{enumerate}
\item \begin{enumerate}
\item Tableau de signes de $V(x)=(x+1)(x-2)(x-3)$ :
\begin{center}
\begin{tabular}{|c|ccccccccc|}\hline
$x$ & $-\infty$ &  & $-1$ &  & $2$ &  & $3$ &  & $+\infty$ \\\hline
$x+1$ &  & $-$ & $0$ & $+$ &  & $+$ &  & $+$ &  \\
$x-2$ &  & $-$ &  & $-$ & $0$ & $+$ &  & $+$ &  \\
$x-3$ &  & $-$ &  & $-$ &  & $-$ & $0$ & $+$ &  \\\hline
$V(x)$ &  & $-$ & $0$ & $+$ & $0$ & $-$ & $0$ & $+$ &  \\\hline
\end{tabular}
\end{center}
\item $V(x)>0 \Leftrightarrow x\in ]-1,2[\,\cup\,]3,+\infty[$.
\item Contraintes : $x>0$, $x\le 2{,}5$ et $V(x)>0$. L'intersection de $]-1,2[\cup]3,+\infty[$ avec $]0\,;\,2{,}5]$ est $]0,2[$. L'entreprise retient $x\in ]0\,;\,2[$ (en mètres).
\end{enumerate}
\item $V(x)\ge 6 \Leftrightarrow x^3-4x^2+x+6\ge 6 \Leftrightarrow x^3-4x^2+x\ge 0 \Leftrightarrow x(x^2-4x+1)\ge 0$.
Racines de $x^2-4x+1$ : $\Delta=16-4=12$, $x=2-\sqrt{3}$ ou $x=2+\sqrt{3}$. On ordonne : $0<2-\sqrt{3}<2+\sqrt{3}$.
\begin{center}
\begin{tabular}{|c|ccccccccc|}\hline
$x$ & $-\infty$ &  & $0$ &  & $2-\sqrt{3}$ &  & $2+\sqrt{3}$ &  & $+\infty$ \\\hline
$x$ &  & $-$ & $0$ & $+$ &  & $+$ &  & $+$ &  \\
$x-(2-\sqrt{3})$ &  & $-$ &  & $-$ & $0$ & $+$ &  & $+$ &  \\
$x-(2+\sqrt{3})$ &  & $-$ &  & $-$ &  & $-$ & $0$ & $+$ &  \\\hline
$x(x^2-4x+1)$ &  & $-$ & $0$ & $+$ & $0$ & $-$ & $0$ & $+$ &  \\\hline
\end{tabular}
\end{center}
Solution dans $\mathbb{R}$ : $[0\,;\,2-\sqrt{3}]\cup[2+\sqrt{3},+\infty[$.
Avec $x>0$ et $x\le 2{,}5$ : comme $2+\sqrt{3}\approx 3{,}73>2{,}5$, seul le premier intervalle convient. Valeurs convenables : $x\in ]0\,;\,2-\sqrt{3}]$.
\end{enumerate}
\textbf{Points de vigilance :} citer la condition $x>0$ à chaque intersection ; inéquation large $\Rightarrow$ crochets fermés sur les racines ; ne pas oublier de rejeter $2+\sqrt{3}$ au vu de la contrainte d'encombrement.
\end{corrige}

\begin{corrige}[Exercice 10 — Le transporteur Douala–Yaoundé (type Bac)]
\textbf{Barème indicatif :} 1) 1 pt ; 2) 3 pts ; 3) 1 pt ; 4) 2 pts.
\begin{enumerate}
\item Lundi : $40x+20y+5z=\num{225000}$, et en divisant par $5$ : $8x+4y+z=\num{45000}$.
Mardi : $30x+25y+10z=\num{237500}$, et en divisant par $5$ : $6x+5y+2z=\num{47500}$.
Mercredi : $50x+15y+8z=\num{239000}$. D'où le système $(S)$.
\item De $(1)$ : $z=\num{45000}-8x-4y$. On reporte dans $(2)$ et $(3)$ :
$(2)$ : $6x+5y+2(\num{45000}-8x-4y)=\num{47500} \Leftrightarrow -10x-3y=-\num{42500}$, soit $10x+3y=\num{42500}$ $(2')$.
$(3)$ : $50x+15y+8(\num{45000}-8x-4y)=\num{239000} \Leftrightarrow -14x-17y=-\num{121000}$, soit $14x+17y=\num{121000}$ $(3')$.
Système $2\times 2$ : $\Delta=10\times 17-14\times 3=170-42=128\neq 0$ : solution unique.
$\Delta_x=\num{42500}\times 17-\num{121000}\times 3=\num{722500}-\num{363000}=\num{359500}$, donc $x=\dfrac{\num{359500}}{128}=\num{2808,59375}$.
$\Delta_y=10\times\num{121000}-14\times\num{42500}=\num{1210000}-\num{595000}=\num{615000}$, donc $y=\dfrac{\num{615000}}{128}=\num{4804,6875}$.
$z=\num{45000}-8x-4y=\num{45000}-\num{22468,75}-\num{19218,75}=\num{3312,5}$.
\item Prix des tickets : \emph{Classique} $\approx \num{2809}$ FCFA, \emph{Confort} $\approx \num{4805}$ FCFA, \emph{VIP} $=\num{3312,5}$ FCFA (soit $\num{3313}$ FCFA au franc près).
Vérification (lundi) : $40\times\num{2808,59375}+20\times\num{4804,6875}+5\times\num{3312,5}=\num{112343,75}+\num{96093,75}+\num{16562,5}=\num{225000}$ \checkmark.
\item Soit $n$ le nombre de tickets \emph{Classique} ($n\in\mathbb{N}$). Recette du jeudi :
$n\times\num{2808,59375}+12\times\num{4804,6875}>\num{100000}$
$\Leftrightarrow \num{2808,59375}\,n+\num{57656,25}>\num{100000}$
$\Leftrightarrow n>\dfrac{\num{42343,75}}{\num{2808,59375}}\approx 15{,}08$.
Comme $n$ est un entier naturel, le nombre minimal est $n=16$ tickets \emph{Classique}.
\end{enumerate}
\textbf{Points de vigilance :} garder les valeurs exactes (fractions ou décimaux exacts) jusqu'au bout ; l'inéquation est stricte, donc $n>15{,}08$ impose $n\ge 16$ (et non $15$) ; conclure par une phrase en contexte.
\end{corrige}

\begin{corrige}[Exercice 11 — Étude d'un système et modélisation (type Bac)]
\textbf{Barème indicatif :} Partie A : 1) 1 pt ; 2) 1,5 pts ; 3) 1 pt. Partie B : 1) 0,5 pt ; 2) 1,5 pts ; 3) 1,5 pts.

\textbf{Partie A.}
\begin{enumerate}
\item $\Delta=2\times 6-4\times 3=12-12=0$. Le système n'admet pas de solution unique : il a soit aucune solution, soit une infinité.
\item \begin{enumerate}
\item Pour $k=10$ : $4x+6y=10 \Leftrightarrow 2(2x+3y)=2\times 5 \Leftrightarrow 2x+3y=5$ : c'est exactement la première équation.
\item Les deux équations sont équivalentes : infinité de solutions. En exprimant $y$ en fonction de $x$ : $y=\dfrac{5-2x}{3}$.
\[ S=\left\{\left(x\,;\,\tfrac{5-2x}{3}\right),\ x\in\mathbb{R}\right\}. \]
\end{enumerate}
\item Pour $k=12$ : la deuxième équation équivaut à $2x+3y=6$, alors que la première impose $2x+3y=5$. Or $5\neq 6$ : contradiction. $(S_{12})$ n'admet \textbf{aucune solution}.
\end{enumerate}

\textbf{Partie B.}
\begin{enumerate}
\item $\begin{cases} 2x+3y=95\\ 5x+4y=160 \end{cases}$ (en milliers de FCFA).
\item $\Delta=2\times 4-5\times 3=8-15=-7\neq 0$ : solution unique.
$\Delta_x=95\times 4-160\times 3=380-480=-100 \Rightarrow x=\dfrac{-100}{-7}=\dfrac{100}{7}$.
$\Delta_y=2\times 160-5\times 95=320-475=-155 \Rightarrow y=\dfrac{-155}{-7}=\dfrac{155}{7}$.
Prix : sac A $=\dfrac{\num{100000}}{7}\approx \num{14286}$ FCFA ; sac B $=\dfrac{\num{155000}}{7}\approx \num{22143}$ FCFA.
\item Budget : $500$ milliers de FCFA. $10x+ny\le 500 \Leftrightarrow 10\times\dfrac{100}{7}+n\times\dfrac{155}{7}\le 500 \Leftrightarrow 1000+155n\le 3500 \Leftrightarrow n\le \dfrac{2500}{155}=\dfrac{500}{31}\approx 16{,}13$.
$n$ entier naturel maximal : $n=16$ sacs d'engrais B.
\end{enumerate}
\textbf{Points de vigilance :} en Partie A, bien distinguer les cas $\Delta=0$ « infinité » (équations proportionnelles, seconds membres compris) et « aucune solution » ; en Partie B, préciser l'unité (milliers de FCFA) et convertir le résultat final en FCFA ; pour la question 3, l'inéquation est large (budget « sans dépasser ») et $n$ doit être entier.
\end{corrige}

\newpage

\coursec{Fiche bilan}

\begin{popfigure}
\begin{tikzpicture}[mindmap, grow cyclic,
every node/.style={concept, font=\small},
root concept/.append style={concept color=popblue, font=\bfseries},
level 1/.append style={level distance=4.6cm, sibling angle=60},
level 2/.append style={level distance=2.6cm, sibling angle=40, font=\scriptsize}]
\node[root concept]{Équations, inéquations et systèmes}
  child[concept color=poporange]{ node {Équations de degré 3}
    child { node {Racine évidente $\alpha$} }
    child { node {$P(x)=(x-\alpha)(ax^2+bx+c)$} }
  }
  child[concept color=popgreen]{ node {Inéquations de degré 3}
    child { node {Tableau de signes} }
  }
  child[concept color=poppurple]{ node {Inconnue auxiliaire}
    child { node {Poser $X=\varphi(x)$} }
    child { node {Revenir à $x$} }
  }
  child[concept color=poppink]{ node {Systèmes $2\times 2$}
    child { node {Substitution} }
    child { node {Combinaison} }
    child { node {Déterminant} }
  }
  child[concept color=popgold]{ node {Systèmes $3\times 3$}
    child { node {Pivot de Gauss} }
    child { node {Remontée} }
  }
  child[concept color=popblue!60]{ node {Modélisation}
    child { node {Problèmes concrets} }
  };
\end{tikzpicture}
\legende{Carte mentale de la leçon : les six compétences clés autour des équations, inéquations et systèmes.}
\end{popfigure}

\begin{aretenir}[L'essentiel de la leçon]
\textbf{Définitions clés}
\begin{itemize}
  \item Une \cle{équation de degré 3} est une équation de la forme $ax^3+bx^2+cx+d=0$ avec $a\neq 0$.
  \item Un réel $\alpha$ est une \cle{racine} du polynôme $P$ si $P(\alpha)=0$.
  \item Un \cle{système linéaire} est un ensemble d'équations du premier degré à plusieurs inconnues, résolu simultanément.
\end{itemize}

\textbf{Formules et théorèmes à connaître par cœur}
\begin{center}
\begin{tabularx}{\linewidth}{|l|X|}
\hline
\rowcolor{popblueL} \textbf{Notion} & \textbf{Résultat} \\
\hline
Factorisation & Si $P(\alpha)=0$, alors $P(x)=(x-\alpha)(ax^2+bx+c)$, obtenu par \textbf{identification} ou \textbf{division euclidienne}. \\
\hline
Discriminant & $\Delta=b^2-4ac$ ; $\Delta>0$ : deux racines $x_{1}=\dfrac{-b-\sqrt{\Delta}}{2a}$ et $x_{2}=\dfrac{-b+\sqrt{\Delta}}{2a}$ ; $\Delta=0$ : racine double $x_0=-\dfrac{b}{2a}$ ; $\Delta<0$ : pas de racine réelle. \\
\hline
Signe de $ax^2+bx+c$ & Du signe de $a$ à l'extérieur des racines, du signe contraire entre les racines (si $\Delta>0$). \\
\hline
Déterminant $2\times2$ & $D=\begin{vmatrix} a & b \\ a' & b' \end{vmatrix}=ab'-a'b$. Si $D\neq0$ : solution unique $x=\dfrac{D_x}{D}$, $y=\dfrac{D_y}{D}$. \\
\hline
Pivot de Gauss & Éliminer $x$ dans $(E_2)$ et $(E_3)$, puis $y$ dans $(E_3)$, puis \textbf{remonter} : $z$, puis $y$, puis $x$. \\
\hline
\end{tabularx}
\end{center}

\textbf{Méthodes express}
\begin{itemize}
  \item \textbf{Racine évidente} : tester $\pm1$, $\pm2$, $\pm3$ et les diviseurs de $d$ dans $P(x)=0$.
  \item \textbf{Identification} : développer $(x-\alpha)(ax^2+bx+c)$ et égaler les coefficients.
  \item \textbf{Inéquation de degré 3} : factoriser complètement, puis tableau de signes (une ligne par facteur).
  \item \textbf{Inconnue auxiliaire} : poser $X=x^2$ (équations bicarrées), $X=\sqrt{x}$, $X=\dfrac{1}{x}$, résoudre en $X$, \emph{puis revenir à $x$ sans oublier les conditions d'existence}.
  \item \textbf{Système $2\times2$} : combinaison si les coefficients s'y prêtent, substitution si une inconnue a pour coefficient $1$.
  \item \textbf{Système $3\times3$} : toujours vérifier la solution trouvée dans les trois équations de départ.
\end{itemize}
\end{aretenir}

\begin{exemplebox}[Illustration : factorisation par identification]
Soit $P(x)=x^3-4x^2+x+6$. On teste $x=-1$ : $P(-1)=-1-4-1+6=0$, donc $-1$ est racine évidente. On écrit $P(x)=(x+1)(x^2+bx+c)$. En développant : $(x+1)(x^2+bx+c)=x^3+(b+1)x^2+(b+c)x+c$. Par identification : $b+1=-4$ donc $b=-5$ ; $c=6$ ; vérification : $b+c=-5+6=1$, conforme au coefficient de $x$. Ainsi $P(x)=(x+1)(x^2-5x+6)=(x+1)(x-2)(x-3)$.
\end{exemplebox}

\begin{exemplebox}[Illustration : pivot de Gauss]
Résolvons $\begin{cases} x+y+z=6 \\ 2x-y+z=3 \\ x+2y-z=2 \end{cases}$. On effectue $(E_2)\leftarrow(E_2)-2(E_1)$ : $-3y-z=-9$, et $(E_3)\leftarrow(E_3)-(E_1)$ : $y-2z=-4$. Puis $(E_3)\leftarrow 3(E_3)+(E_2)$ : $-7z=-21$, donc $z=3$. Remontée : $y=-4+2z=2$, puis $x=6-y-z=1$. Vérification : $1+2+3=6$, $2-2+3=3$, $1+4-3=2$. Solution : $S=\{(1\,;\,2\,;\,3)\}$.
\end{exemplebox}

\begin{aretenir}[Checklist avant le Bac]
Avant l'épreuve, vérifie que tu sais :
\begin{itemize}
  \item[$\square$] Tester rapidement les valeurs $\pm1$, $\pm2$, $\pm3$ pour trouver une racine évidente.
  \item[$\square$] Factoriser $P(x)$ en $(x-\alpha)(ax^2+bx+c)$ par identification des coefficients.
  \item[$\square$] Effectuer la division euclidienne d'un polynôme de degré 3 par $(x-\alpha)$.
  \item[$\square$] Calculer un discriminant et conclure sur le nombre de racines réelles.
  \item[$\square$] Dresser un tableau de signes complet pour résoudre une inéquation de degré 3.
  \item[$\square$] Poser correctement une inconnue auxiliaire et penser à revenir à l'inconnue initiale.
  \item[$\square$] Résoudre un système $2\times2$ par substitution et par combinaison.
  \item[$\square$] Calculer un déterminant $2\times2$ et discuter le cas $D=0$.
  \item[$\square$] Appliquer le pivot de Gauss à un système $3\times3$ et effectuer la remontée.
  \item[$\square$] Vérifier une solution en la substituant dans toutes les équations de départ.
  \item[$\square$] Traduire un problème concret (prix, âges, mélanges) en équation ou en système.
  \item[$\square$] Rédiger proprement : ensemble de définition, conclusion avec $S=\{\dots\}$ ou intervalle.
\end{itemize}
\end{aretenir}

\findecours

\end{document}
